\chapter{Realizaciones espectrales, escalares y nulas de la Ley General de Masas: familias materiales, transportes CKM–PMNS y sectores fotónico y gluónico}
\label{chap:realizaciones-familias-mezclas-masas}

\section{Codominios de las realizaciones materiales: operador espectral, escalar, nulidad de calibre, compuesto y polo resonante}
Una fibra puede publicar un operador espectral, un escalar después de una
compresión física, una nulidad exacta de calibre, un operador compuesto o un
polo resonante con anchura. Estas salidas poseen codominios distintos y no se
cuentan como clases del dominio \(81/56/13/324\).

\section{Extensiones nulas, propagación entre fronteras y clausura electrónica}

Los sectores fotónico y gluónico se realizan mediante nulidad exacta, mientras
la clausura electrónica ocupa el centro material. El desarrollo completo
distingue propagación nula, entrelazamiento ida--retorno, interfaz física y
criterios de cierre sin transformar una nulidad en ausencia de estructura.

\begingroup
\let\chapter\section
\let\section\subsection
\let\subsection\subsubsection
\let\subsubsection\paragraph
\let\clearpage\relax
\let\cleardoublepage\relax
\chapter{Extensiones nulas y clausura material}
\label{chap:realizaciones-nula-material}
\index{fotón!interfaz física}\index{electrón!clausura material}
\index{extensión nula}\index{cono causal partido}

\section{\(4\) realizaciones del estado material común: nula, masiva, compuesta y topológica}
\label{sec:cuatro-modos-realizacion}

Una cifra aislada determina una coordenada, pero no identifica la estructura
que la produce, la memoria que conserva ni la operación que la convierte en
magnitud. HMT se sitúa antes de esa proyección: su objeto primario es una
extensión compatible que conserva las leyes aritméticas, la orientación
TRIT, los cocientes, la fase y la memoria de frontera. Mecánica Dimensional
especifica después qué componente se realiza como magnitud, cuál permanece
como memoria y qué transformación produce el observable.

Los desarrollos de esta parte distinguen cuatro modos por su operación de
cierre:
\begin{enumerate}
 \item una extensión abierta y nula transporta una perturbación entre
 fronteras;
 \item una extensión cerrada y persistente adquiere norma interior y define
 el candidato material;
 \item una familia de extensiones ponderada por acción define un ensamble
 térmico;
 \item una conexión que compara secciones en puntos distintos proporciona el
 codominio de la realización gravitatoria.
\end{enumerate}
La fuente común no identifica los cuatro codominios. Los mapas físicos
correspondientes se mantienen discontinuos:
\[
 \mathfrak R_\gamma:
 \Omega_{\HMT}^{\rm surv}\dashrightarrow\mathcal H_{\rm rad},
 \qquad
 \mathfrak R_m:
 \Omega_{\HMT}^{\rm surv}\dashrightarrow\mathcal H_{\rm mat},
\]
\[
 \mathfrak R_\Theta:
 \mathcal P(\Omega_{\HMT}^{\rm surv})\dashrightarrow\mathcal S_{\rm th},
 \qquad
 \mathfrak R_{\rm grav}:
 \Omega_{\HMT}^{\rm surv}\dashrightarrow\mathsf{Cartan}_{1,3}.
\]
Su unidad reside en la genealogía; su diferencia, en la lectura y en las
estructuras adicionales exigidas por cada realización.

\section{Propagación nula entre fronteras}
\label{sec:propagacion-nula-fotonica}

Sea \(B\) una base de transporte. Una \emph{sección nula} de signo
\(\pm\) es una aplicación
\[
 s_\pm:B\longrightarrow\mathcal A_{\rm sp}P_\pm
\]
tal que \(s_\pm(x)\in\mathcal A_{\rm sp}P_\pm\) para todo \(x\in B\).
Por la \cref{prop:descomposicion-espectral-partida},
\[
 N_{\rm sp}(s_\pm(x))=0
 \qquad(x\in B).
\]
Si \(\gamma\) es una ruta orientada entre dos fronteras,
\[
 \partial\gamma=x_{\rm abs}-x_{\rm em},
\]
la interfaz fotónica propuesta asigna a \(\gamma\) una sección nula
extendida. El adjetivo \emph{extendida} indica que su identidad pertenece al
transporte completo y no a un punto aislado.

La ruta conjugada se obtiene por inversión,
\[
 \gamma^\vee=C_{\rm rev}\gamma,
 \qquad
 \partial\gamma^\vee=x_{\rm em}-x_{\rm abs}.
\]
Relativamente al lector de orientación declarado, el par
\((\gamma,\gamma^\vee)\) representa las lecturas avanzada y retardada. Las
dos piernas se ordenan secuencialmente o se alojan en fibras distintas. No se
suman automáticamente dentro de una misma fibra: si se activan
simultáneamente las dos componentes conjugadas, la identidad exacta
\[
 \boxed{N_{\rm sp}(uP_++vP_-)=uv}
\]
muestra que, para \(u,v>0\), la sección abandona la frontera nula y entra en
el interior positivo.

\subsection{Longitud, ida y retorno}

Fijada la carta dimensional de velocidad,
\(\widehat c=e^{-\mathcal E}\), y una base positiva \(u_L\) de la recta de
longitud, se define
\[
 L_{\rm int}(\gamma)=\widehat c\,|\gamma|\,u_L.
\]
La inversión conserva el número de pasos:
\[
 L_{\rm int}(\gamma^\vee)=L_{\rm int}(\gamma),
 \qquad
 L_{\leftrightarrow}=2L_{\rm int}(\gamma).
\]
Estas identidades son cinemáticas. Convertirlas en propagación
electromagnética exige que un mismo mapa reproduzca la ley de Gauss, la
dinámica de calibre, la dispersión nula y la reducción a dos grados
transversales.

\subsection{Entropía de entrelazamiento entre las \(2\) lecturas}
\label{sec:entrelazamiento-ida-retorno}

La inversión de una ruta no basta para afirmar entrelazamiento cuántico. Una
formulación mínima exige primero una completación de Hilbert. Sean
\(\mathcal H_{\rm adv}\) y \(\mathcal H_{\rm ret}\) dos espacios complejos
con familias ortonormales
\(\{|i\rangle_{\rm adv}\}_{i\in I}\) y
\(\{|i\rangle_{\rm ret}\}_{i\in I}\), donde \(I\) es finito. Para una
probabilidad \(p=(p_i)_{i\in I}\), definimos
\[
 |\Psi_p\rangle
 :=\sum_{i\in I}\sqrt{p_i}\,
 |i\rangle_{\rm adv}\otimes|i\rangle_{\rm ret}.
 \tag{E1}\label{eq:estado-entrelazado-ida-retorno}
\]

\begin{proposition}[Entropía bicapa de ida y retorno]
\label{prop:entropia-entrelazamiento-ida-retorno}
Las matrices reducidas del estado
\(\rho_p=|\Psi_p\rangle\langle\Psi_p|\) son
\[
 \rho_{\rm adv}
 =\sum_{i\in I}p_i|i\rangle_{\rm adv}\langle i|,
 \qquad
 \rho_{\rm ret}
 =\sum_{i\in I}p_i|i\rangle_{\rm ret}\langle i|,
 \tag{E2}\label{eq:reducidas-ida-retorno}
\]
y poseen la misma entropía de von Neumann:
\[
 \boxed{
 S_{\rm ent}(p)
 =-\operatorname{Tr}(\rho_{\rm adv}\log\rho_{\rm adv})
 =-\sum_{i\in I}p_i\log p_i.}
 \tag{E3}\label{eq:entropia-ida-retorno}
\]
Aquí se adopta la convención \(0\log0:=0\). El estado es separable
exactamente cuando una sola probabilidad vale uno.
\end{proposition}

\begin{proof}
Al efectuar la traza parcial, la ortogonalidad anula los términos con
\(i\ne j\). Los valores propios no nulos de cada matriz reducida son, por
tanto, los \(p_i\), lo que da \eqref{eq:entropia-ida-retorno}. El rango de
Schmidt es uno exactamente en el caso determinista.
\end{proof}

La proposición formaliza, una vez elegida la completación, una entropía entre
las lecturas avanzada y retardada. No deduce las probabilidades \(p_i\), no
identifica esas lecturas con los propagadores avanzado y retardado de una
teoría de campos y no convierte una correlación clásica de rutas en
entrelazamiento. La construcción de un isomorfismo que envíe la memoria HMT
a estados físicos bipartitos, preserve la evolución y fije \(p\) pertenece a
la \textsc{interfaz física}. La entropía nula de una distribución
determinista constituye el control negativo inmediato.

\begin{statusbox}
La proposición está
\textsc{demostrada con estructura de partida explícita}: presupone la
completación de Hilbert, las bases ortonormales y la distribución \(p\). La construcción
del entrelazador que realice físicamente ambas capas es una
\textsc{interfaz física}; no forma parte de la proposición.
\end{statusbox}

\begin{definition}[Interfaz fotónica de HMT--MD]
\label{def:interfaz-fotonica-hmt-md}
Una realización fotónica de HMT--MD es una sección abierta sostenida sobre
una dirección nula del álgebra partida, transportada entre fronteras
conjugadas y protegida por una simetría que conserva exactamente la norma
nula.
\end{definition}

Los idempotentes \(P_\pm\) no son, por sí solos, fotones físicos ni las dos
polarizaciones de Maxwell. La realización completa requiere un
entrelazador
\[
 \mathfrak R_\gamma:
 \partial_{\rm null}\mathcal C_{\rm sp}^{+}
 \dashrightarrow\mathcal H_{\rm Maxwell}^{\rm phys}
\]
que preserve el producto interno, la evolución y la acción de calibre, y que
seleccione exactamente dos polarizaciones transversales en \(3+1\)
dimensiones.

\begin{criterion}[Sector fotónico]
\label{crit:sector-fotonico}
La identificación física queda refutada si la evolución abandona
\(\partial_{\rm null}\mathcal C_{\rm sp}^{+}\), genera masa sin el mecanismo
de ruptura declarado, no reproduce la dispersión nula o no deja exactamente
dos polarizaciones físicas tras la reducción de calibre.
\end{criterion}

\section{Clausura material y centro electrónico}
\label{sec:clausura-material-electron}

El acoplamiento simultáneo de dos direcciones nulas conjugadas produce, para
\(u,v>0\),
\[
 N_{\rm sp}(uP_++vP_-)=uv>0.
\]
Ninguna componente aislada posee norma distinta de cero; la norma interior
surge de su acoplamiento. Esta afirmación algebraica no suma masas
preexistentes y no identifica todavía las direcciones nulas con fotones
\emph{on shell}.

   La clasificación estructural de HMT contiene un único punto fijo de la inversión celular,
\[
 (5,5).
\]
El origen relativo de acción del electrón es \(v_e=0\). En la realización
partida normalizada se fija
\[
 Z_e=P_++P_-=I,
 \qquad
 N_{\rm sp}(Z_e)=1.
\]
El punto \((5,5)\), el origen \(v_e=0\), la firma cruda
\((24,0,12,0,-12)\) y la palabra transversal \(555555\) son cuatro objetos
distintos.

\begin{definition}[Candidato de clausura electrónica]
\label{def:candidato-clausura-electronica}
La clausura resonante de modos nulos electromagnéticos conjugados se
formaliza como el programa que busca una sección material central normalizada
obtenida por el acoplamiento de dos modos nulos conjugados, estabilizada por
una clausura de ruta y por su holonomía.
\end{definition}

La existencia de la extensión, del retorno y del centro del atlas no prueba
la compactitud mínima ni la estabilidad del cierre. El teorema dinámico
pendiente debe construir
\[
 \gamma_e=\xi_1\circ\cdots\circ\xi_N,
 \qquad
 \partial\gamma_e=0,
 \qquad
 \varepsilon(\gamma_e)=-1,
\]
donde cada \(\xi_j\) sea localmente nulo, el borde total se anule y
\(\varepsilon\) sea la holonomía de signo. La masa se asociaría al cociclo
global de clausura, no a una suma de masas nulas locales. Este programa no
reabre el electrón central \((5,5)\), el atlas ni la cadena
extensión--ruta--registro--carácter ya construidos.

\begin{proposition}[Criterio relativista de interior causal]
\label{prop:cierre-nulos-interior-causal}
Sean \(p_1,\ldots,p_N\) momentos nulos futuros en un espacio de Minkowski con
firma temporal positiva, \(p_i^2=0\), y sea
\[
 P=\sum_{i=1}^{N}p_i.
\]
Si todos los \(p_i\) son colineales, entonces \(P^2=0\). Si al menos dos son
futuros y no colineales, entonces \(P^2>0\), y puede definirse
\[
 m^2c^2=P^2.
\]
\end{proposition}

\begin{proof}
Para vectores nulos futuros se tiene
\(p_i\cdot p_j\geq0\), con igualdad exactamente cuando son colineales.
Como
\[
 P^2=\sum_i p_i^2+2\sum_{i<j}p_i\cdot p_j,
\]
la primera suma es nula. Todos los productos cruzados se anulan en el caso
colineal; si existe un par no colineal, al menos uno es estrictamente
positivo.
\end{proof}

Por tanto, una ruta cerrada no es masiva por mera topología, ni una ruta
abierta es nula por el solo hecho de poseer borde. La realización exige que
el cierre geométrico y la suma de momentos entren conjuntamente en el
interior causal.

\subsection{Möbius, conjugación y ocupación posterior}

Bajo las hipótesis de levantamiento y holonomía del
\cref{thm:dos-cuatro-pi-estadistica}, una vuelta invierte la orientación y
dos la restauran. La banda de Möbius organiza partícula y antipartícula como
orientaciones conjugadas de un mismo tipo de cierre. Su identificación física
requiere una representación que invierta la carga, preserve la norma y
entrelace la evolución. La geometría de Möbius no implica por sí sola el
álgebra de ocupación fermiónica; sus relaciones operatorias se
introducen como hipótesis operatorias en el capítulo siguiente.

\begin{criterion}[Cierre electrónico]
\label{crit:cierre-electronico}
La realización queda refutada si el centro \((5,5)\) no se eleva a una
extensión cerrada, si la holonomía no selecciona el sector impar, si el
Hamiltoniano carece de una solución ligada estable y única salvo simetría, o
si carga, dispersión y factor de forma no concuerdan con la realización
electrónica.
\end{criterion}

\endgroup

\begingroup
\let\chapter\section
\let\section\subsection
\let\subsection\subsubsection
\let\subsubsection\paragraph
\let\clearpage\relax
\section{Realizaciones físicas por objeto, representación y codominio}
Cada realización declara su objeto primario, su representación o fibra y su salida nativa. Una coordenada escalar sólo se publica cuando el propio objeto y su morfismo de realización la fuerzan.


\endgroup
\begingroup
\let\chapter\section
\let\section\subsection
\let\subsection\subsubsection
\let\subsubsection\paragraph
\let\clearpage\relax

\clearpage
\section{Realizaciones elementales del Modelo Estándar}
\begingroup
\footnotesize
\setlength{\tabcolsep}{1.9pt}
\begin{longtable}{L{2.03cm}L{2.03cm}L{1.87cm}L{2.34cm}L{6.55cm}}
\toprule
\rowcolor{HMTNavy}\HMTtablehead{Partícula} & \HMTtablehead{Carga} & \HMTtablehead{Espín} & \HMTtablehead{Estadística} & \HMTtablehead{Realización HMT--MD} \\
\midrule
\endfirsthead
\multicolumn{5}{l}{\sffamily\scriptsize\color{HMTGray}Continuación de la tabla}\\[-1pt]
\toprule
\rowcolor{HMTNavy}\HMTtablehead{Partícula} & \HMTtablehead{Carga} & \HMTtablehead{Espín} & \HMTtablehead{Estadística} & \HMTtablehead{Realización HMT--MD} \\
\midrule
\endhead
\midrule
\endfoot
\bottomrule
\endlastfoot
\(e^-\) & \(-1\) & \(1/\allowbreak{}2\) & fermiónica & punto fijo \((5,5)\), vacancia orientada \\
\rowcolor{HMTLight}\(\mu^-\) & \(-1\) & \(1/\allowbreak{}2\) & fermiónica & bloque leptónico de rango 8 \\
\(\tau^-\) & \(-1\) & \(1/\allowbreak{}2\) & fermiónica & bloque leptónico de rango 16 \\
\rowcolor{HMTLight}\(\nu_e\) & \(0\) & \(1/\allowbreak{}2\) & fermiónica & vector de interacción en soporte neutro 20 \\
\(\nu_\mu\) & \(0\) & \(1/\allowbreak{}2\) & fermiónica & vector de interacción en soporte neutro 20 \\
\rowcolor{HMTLight}\(\nu_\tau\) & \(0\) & \(1/\allowbreak{}2\) & fermiónica & vector de interacción en soporte neutro 20 \\
\(u\) & \(+2/\allowbreak{}3\) & \(1/\allowbreak{}2\) & fermiónica & rama up, generación 1, triplete de color \\
\rowcolor{HMTLight}\(d\) & \(-1/\allowbreak{}3\) & \(1/\allowbreak{}2\) & fermiónica & rama down, generación 1, triplete de color \\
\(c\) & \(+2/\allowbreak{}3\) & \(1/\allowbreak{}2\) & fermiónica & rama up, generación 2, triplete de color \\
\rowcolor{HMTLight}\(s\) & \(-1/\allowbreak{}3\) & \(1/\allowbreak{}2\) & fermiónica & rama down, generación 2, triplete de color \\
\(t\) & \(+2/\allowbreak{}3\) & \(1/\allowbreak{}2\) & fermiónica & rama up, generación 3, triplete de color \\
\rowcolor{HMTLight}\(b\) & \(-1/\allowbreak{}3\) & \(1/\allowbreak{}2\) & fermiónica & rama down, generación 3, triplete de color \\
\(\gamma\) & \(0\) & \(1\) & bosónica & sección nula, dos polarizaciones \\
\rowcolor{HMTLight}\(g^a\) & \(0\) & \(1\) & bosónica & adjunta de color, ocho generadores \\
\(W^\pm\) & \(\pm1\) & \(1\) & bosónica & bloque vectorial cargado conjugado \\
\rowcolor{HMTLight}\(Z^0\) & \(0\) & \(1\) & bosónica & bloque vectorial neutral \\
\(H\) & \(0\) & \(0\) & bosónica & sección escalar persistente \\
\end{longtable}
\endgroup
Las antipartículas se obtienen por \(C_M\). Los antiquarks pertenecen al triplete conjugado; \(W^-\) es el conjugado de \(W^+\); fotón, \(Z\), Higgs y gluones poseen las conjugaciones propias de sus representaciones.

El soporte B1 exacto de las clases elementales tiene rangos

\[
 1,\quad8,\quad16,\quad
 3\times20,\quad
 3\times16,\quad
 3\times20.
\]
Por tanto, el número de incidencias clase--ruta es

\[
 1+8+16+60+48+60=193.
\]
Las incidencias no son 193 rutas distintas: los tres sabores neutrínicos comparten el soporte neutro de rango veinte; los sabores up reutilizan un soporte de rango dieciséis y los sabores down, uno de rango veinte. El proyector de sabor debe refinar estos soportes. En la unión quark de 36 celdas, el color queda exactamente

\[
 12_R+12_G+12_B.
\]
Esta igualdad certifica el soporte cromático global, pero no autoriza todavía a elegir una celda por sabor.

\Needspace{7\baselineskip}
\section{Descomposición del sector leptónico por generación, carga, espín y carácter de masa}
\Needspace{7\baselineskip}
\subsection{Electrón: centro, vacancias y orientación}
El electrón es excepcional porque su fibra contiene el único punto fijo \((5,5)\). Su reducción a masa escalar no requiere romper una órbita. La ruta activa es

\[
 e_{\rm act}=(2,2,5,-4,-3,-2)^2,
\]
mientras que la envolvente histórica

\[
 e_{\rm env}=(4,3,5,-4,-3,-5)^2
\]
conserva signo y máximo absoluto de las dos hojas. La envolvente retiene la palabra

\[
 \tau_e=+++---+++---,
\]
pero no sustituye la historia activa.

Cuatro datos que deben mantenerse distintos son:

\[
 \text{firma cruda}=(24,0,12,0,-12),
\]
\[
 \text{firma par CT--108}=(-12,-4,12,-6,12),
\]
\[
 \text{vector de acontecimientos}=e_{\rm act},
\qquad
 \text{origen afín}=v_e=0.
\]
El selector basal \((-22,+15)\) posee dos lecturas internas complementarias. La descomposición regional da

\[
 -22=-(27-5),\qquad +15=5\cdot3=\binom62,
\]
y el marco de vacancias

\[
 V_e=
 \begin{pmatrix}
 21&22\\
 6&7
 \end{pmatrix}
\]
satisface

\[
 \det V_e=21\cdot7-22\cdot6=15.
\]
El \(22\) mide el déficit orientado de la extensión; el \(15\), la retención de área que sobrevive al cierre. Así, la partícula electrónica aparece como persistencia central de una vacancia orientada, no como una masa insertada en la celda central.

La ecuación completa es

\[
 \mathfrak e_e^\beta=
 \frac{\sqrt3}{4}
 \left[
 \exp\!\left(\frac{\varphi}{\pi^2}\right)
 -22\alpha_{\rm HMT}^{3}
 \right]
 (1+15\Delta_4)
 R_{\rm act}^{80-54\alpha_{\rm HMT}+6\Delta_4}.
\]
En la carta electrónica,

\[
 \mathfrak e_e^\beta
 =0.5109989506900008086082571648\ldots\ {\rm MeV}.
\]
\Needspace{7\baselineskip}
\subsection{Realización del muón como sección leptónica de 2.ª generación: fibra, firma y carácter}
El muón reside en una multisección leptónica de rango ocho. Su huella estructural recuperada comprende las celdas

\[
 \{21,23,25,39,43,57,59,61\},
\]
los conteos de giro/cambio \(22/25\), el eje norte--sur,

\[
 (N,E,S,O)=(62,25,7,14),
\qquad
 (h_{369},r_5)=(57,9),
\]
y la coordenada central

\[
 (n_A,s_C)=(42,-3).
\]
El objeto propio es el operador comprimido

\[
 \widehat M_\mu=P_\mu\widehat M P_\mu.
\]
Su reducción a un número exige probar que el proyector de sabor y generación conmuta con el operador, o que el estado preparado tiene varianza nula. Dada la firma enriquecida

\[
 (42,-3,8,0,2),
\]
la evaluación del carácter es

\[
 105.658360294435829\ldots\ {\rm MeV}.
\]
Esta cifra es un control exacto de la evaluación de esa firma; la derivación física completa debe reproducir la firma desde el registro genealógico seleccionado por el acoplamiento. La anchura muónica se obtiene del operador temporal de decaimiento y no de \(\widehat M_\mu\).

\Needspace{7\baselineskip}
\subsection{Realización del tau como sección leptónica de 3.ª generación: fibra, firma y carácter}
El tau ocupa una multisección leptónica de rango dieciséis. Su huella comprende

\[
 17/24,\qquad \text{eje este--oeste},\qquad
 (N,E,S,O)=(29,55,7,17),
\]
\[
 (h_{369},r_5)=(59,6),\qquad
 (n_A,s_C)=(64,0).
\]
La firma enriquecida

\[
 (64,0,25,-1,6)
\]
produce

\[
 1776.8609176314323\ldots\ {\rm MeV}.
\]
Como en el muón, la estructura firme es el bloque operatorial; la masa escalar se obtiene cuando el acoplamiento tau selecciona un bloque escalar o un polo aislado.

\Needspace{7\baselineskip}
\subsection{Fibra neutrínica común y separación de sabores mediante el lector PMNS}
Los tres sabores neutrínicos comparten el espacio

\[
 \mathcal H_\nu=
 \bigoplus_{j=1}^{4}\mathcal H_{\nu,G_j},
\qquad
 \dim\mathcal H_\nu=2+4+6+8=20.
\]
Los vectores

\[
 f_e,\quad f_\mu,\quad f_\tau
\]
constituyen una base de interacción; no son tres autovectores de masa. La profundidad \(G_j\) tampoco es un cuarto sabor estéril: una interpretación estéril requiere un proyector físico adicional.

El extractor neutro se define por

\[
 T_\nu(m)=\sum_{t\in E_m}
 \omega(t)\varepsilon(t)\theta(t),
\qquad
 e_m^\nu=T_\nu(m)-T_\nu(m+6),
\]
\[
 \tau_m^\nu=\operatorname{sgn}(e_m^\nu).
\]
De él se obtienen

\[
 b_{120}^\nu
 =\sum_{a=0}^{3}
 \operatorname{sgn}\!\left(\sum_{m\in I_a}e_m^\nu\right),
\]
\[
 b_\zeta^\nu=\sum_m\tau_m^\nu\tau_{m+3}^\nu.
\]
La firma neutrínica posee la forma

\[
 v_\nu^{(i)}
 =(0,0,k_{\Delta_4}^{(i)},b_{120}^{\nu_i},b_\zeta^{\nu_i}).
\]
La inversión bidireccional cambia \(b_{120}\) y conserva \(b_\zeta\); esta asimetría produce internamente las dos orientaciones de jerarquía.

\Needspace{7\baselineskip}
\subsection{Naturaleza fermiónica del neutrino}
La neutralidad no determina la estadística. El sector neutrínico vive en la completación exterior y porta la holonomía espinorial:

\[
 U_\nu(2\pi)=-I,\qquad U_\nu(4\pi)=I,
\]
\[
 \{a_{\nu,i},a_{\nu,j}^\dagger\}=\delta_{ij}.
\]
Por tanto, cada modo neutrínico físico es fermiónico. La pregunta Dirac/Majorana es posterior: decide si la conjugación intercambia dos secciones distintas o identifica una sección con su conjugada dentro de una representación real. Ninguna de las dos posibilidades convierte el neutrino en bosón.

\Needspace{7\baselineskip}
\subsection{Operador de masa y PMNS}
Sea \(P_\nu\) el proyector de rango tres sobre el subespacio de modos másicos. Primero se diagonaliza

\[
 P_\nu\widehat M_\nu P_\nu
 =\sum_{i=1}^{3}m_i\,|n_i\rangle\langle n_i|.
\]
Después se construye el transporte entre bases:

\[
 (f_e,f_\mu,f_\tau)
 =(n_1,n_2,n_3)U_{\rm PMNS}^\dagger.
\]
Un núcleo de rango completo que lee hoja, orientación, acarreo, vacancia y holonomía puede escribirse

\[
 G_{ab}=
 \frac1{\sqrt{|L_a||N_b|}}
 \sum_{c\in L_a}\sum_{d\in N_b}K_{\rm reg}(c,d),
\]
\[
 U_{\rm PMNS}=\operatorname{polar}(G),
\qquad \det G\ne0.
\]
La fase interna de rango uno produce los ángulos

\[
 (19.6807^\circ,56.4970^\circ,29.6224^\circ)
\]
y, por ello, no puede ser el núcleo PMNS físico completo. El rango completo es una exigencia matemática, no un retoque numérico.

Los observables

\[
 m_1,m_2,m_3,\qquad
 \Delta m_{ij}^2,\qquad
 m_\beta,\qquad
 m_{\beta\beta},\qquad
 \sum_i m_i
\]
son mapas diferentes del mismo espectro. No deben colapsarse en una sola “masa del neutrino”.

\Needspace{7\baselineskip}
\section{Quarks, color, generaciones y mezcla CKM}
\Needspace{7\baselineskip}
\subsection{Posición estructural de los \(6\) sabores de quark}
Cada quark se define por

\[
 \mathfrak q_X=
 (\text{rama de carga},\text{generación},
 \mathbb Z_3^{\rm color},\text{multisección},
 \text{registro genealógico},\text{hoja}).
\]
La posición estructural de los seis sabores queda resumida por sus huellas anteriores a toda carta de masa:

\Needspace{7\baselineskip}
\begingroup
\fontsize{7.2}{8.2}\selectfont
\setlength{\tabcolsep}{1.9pt}
\begin{longtable}{L{0.94cm}L{2.81cm}L{2.03cm}L{1.87cm}L{3.28cm}L{2.18cm}L{1.87cm}}
\toprule
\rowcolor{HMTNavy}\HMTtablehead{Sabor} & \HMTtablehead{Rama/\allowbreak{}generación} & \HMTtablehead{Giro/\allowbreak{}cambio} & \HMTtablehead{Eje} & \HMTtablehead{\((N,E,S,O)\)} & \HMTtablehead{\((h_{369},r_5)\)} & \HMTtablehead{\((n_A,s_C)\)} \\
\midrule
\endfirsthead
\multicolumn{7}{l}{\sffamily\scriptsize\color{HMTGray}Continuación de la tabla}\\[-1pt]
\toprule
\rowcolor{HMTNavy}\HMTtablehead{Sabor} & \HMTtablehead{Rama/\allowbreak{}generación} & \HMTtablehead{Giro/\allowbreak{}cambio} & \HMTtablehead{Eje} & \HMTtablehead{\((N,E,S,O)\)} & \HMTtablehead{\((h_{369},r_5)\)} & \HMTtablehead{\((n_A,s_C)\)} \\
\midrule
\endhead
\midrule
\endfoot
\bottomrule
\endlastfoot
\(u\) & positiva/\allowbreak{}1 & \(21/\allowbreak{}30\) & N--S & \((34,31,31,12)\) & \((51,13)\) & \((11,7)\) \\
\rowcolor{HMTLight}\(d\) & negativa/\allowbreak{}1 & \(20/\allowbreak{}32\) & E--O & \((27,35,18,28)\) & \((47,17)\) & \((17,8)\) \\
\(c\) & positiva/\allowbreak{}2 & \(21/\allowbreak{}29\) & E--O & \((20,38,24,26)\) & \((49,16)\) & \((61,8)\) \\
\rowcolor{HMTLight}\(s\) & negativa/\allowbreak{}2 & \(17/\allowbreak{}27\) & N--S & \((50,32,12,14)\) & \((57,11)\) & \((41,-3)\) \\
\(t\) & positiva/\allowbreak{}3 & \(21/\allowbreak{}26\) & N--S & \((46,27,16,19)\) & \((57,9)\) & \((100,-1)\) \\
\rowcolor{HMTLight}\(b\) & negativa/\allowbreak{}3 & \(21/\allowbreak{}30\) & E--O & \((12,63,9,24)\) & \((47,7)\) & \((71,-5)\) \\
\end{longtable}
\endgroup
Estas posiciones no son seis números elegidos: son seis clases de representación que actúan sobre tripletes de color. La ruta física puede ser una órbita o un bloque de rutas; no tiene que ser un punto.

\Needspace{7\baselineskip}
\subsection{Invarianza escalar del carácter de masa bajo la acción de color}
Sea \(P_q\) el proyector de sabor y generación y \(P_{\rm RGB}\) el proyector a la órbita de color. La masa quark es

\[
 \widehat M_q=
 P_qP_{\rm RGB}\widehat M P_{\rm RGB}P_q.
\]
Si

\[
 [\widehat M_q,U(g)]=0,\qquad g\in SU(3),
\]
entonces

\[
 \widehat M_q=m_qI_3
\]
sobre el triplete irreducible. Esta igualdad explica por qué el sabor porta una masa común aunque sus tres colores ocupen posiciones distintas en el atlas.

\Needspace{7\baselineskip}
\subsection{Carta de esquema y escala}
La coordenada HMT interna debe transportarse a una definición renormalizada:

\[
 m_q^{\mathcal S}(\mu)
 =Z_m^{\mathcal S}(\mu,\mu_0)\,
 \mathcal R_q(\widehat M_q),
\]
con ley de composición

\[
 Z_m(\mu_2,\mu_0)
 =Z_m(\mu_2,\mu_1)Z_m(\mu_1,\mu_0).
\]
Para \(u,d,s\) la comparación natural usa \(\overline{\rm MS}\) a \(2\,{\rm GeV}\); para \(c,b\), la escala de la propia masa renormalizada; para \(t\), debe declararse si la carta es de polo, MSR o cinemática. Sin este transporte, una diferencia numérica no es todavía un residual de la ley.

\Needspace{7\baselineskip}
\subsection{Evaluaciones de las firmas declaradas}
Las coordenadas centrales históricas producen:

\Needspace{7\baselineskip}
\begingroup
\footnotesize
\setlength{\tabcolsep}{2.1pt}
\begin{longtable}{L{4.68cm}L{10.61cm}}
\toprule
\rowcolor{HMTNavy}\HMTtablehead{Sabor} & \HMTtablehead{Evaluación} \\
\midrule
\endfirsthead
\multicolumn{2}{l}{\sffamily\scriptsize\color{HMTGray}Continuación de la tabla}\\[-1pt]
\toprule
\rowcolor{HMTNavy}\HMTtablehead{Sabor} & \HMTtablehead{Evaluación} \\
\midrule
\endhead
\midrule
\endfoot
\bottomrule
\endlastfoot
\(u\) & \(2.165924\allowbreak{}536173\allowbreak{}907\ldots\ {\rm MeV}\) \\
\rowcolor{HMTLight}\(d\) & \(4.679550\allowbreak{}162948\allowbreak{}874\ldots\ {\rm MeV}\) \\
\(c\) & \(1.270500\allowbreak{}838641\allowbreak{}911\ldots\ {\rm GeV}\) \\
\rowcolor{HMTLight}\(s\) & \(92.940093\allowbreak{}907845\allowbreak{}64\ldots\ {\rm MeV}\) \\
\(t\) & \(172.597479\allowbreak{}544019\allowbreak{}1\ldots\ {\rm GeV}\) \\
\rowcolor{HMTLight}\(b\) & \(4.190119\allowbreak{}763299\allowbreak{}790\ldots\ {\rm GeV}\) \\
\end{longtable}
\endgroup
Estas evaluaciones comprueban el carácter una vez dada la firma. La deducción completa exige obtener cada firma, su refinamiento de torre y su carta de esquema a partir del proyector de sabor y del registro genealógico correspondiente.

\Needspace{7\baselineskip}
\subsection{Determinación de los ángulos CKM desde la reflexión racional de las rutas quark}
Sean

\[
 A=1000\alpha_{\rm HMT},\qquad
 h=\frac{C^*}{6},\qquad
 \Delta=\frac{180}{\pi}\Delta_4.
\]
La gramática angular da

\[
 \begin{pmatrix}
 \theta_{12}\\
 \theta_{23}\\
 \theta_{13}
 \end{pmatrix}
 =
 \begin{pmatrix}
 2&-5&28\\
 0&7&-25/2\\
 1&-20&-2/3
 \end{pmatrix}
 \begin{pmatrix}
 A\\h\\\Delta
 \end{pmatrix},
\qquad
 \delta_{\rm CP}=9A.
\]
El falsador lineal exacto es

\[
 3857\,\delta_{\rm CP}
 =9\bigl(
 1528\theta_{12}
 +3380\theta_{23}
 +801\theta_{13}
 \bigr).
\]
A partir de los tres ángulos y de \(\delta_{\rm CP}\) se construye la matriz unitaria \(V_{\rm CKM}\). Su función en la ley de masas no es producir autovalores, sino transportar los operadores down a la base up:

\[
 B_d^{(u)}=V_{\rm CKM}B_dV_{\rm CKM}^*.
\]
\Needspace{7\baselineskip}
\subsection{Invariante de Jarlskog}
El invariante de violación CP es

\[
 J_{\rm CKM}
 =c_{12}c_{23}c_{13}^2
 s_{12}s_{23}s_{13}\sin\delta_{\rm CP},
\]
y la evaluación interna es

\[
 J_{\rm CKM}
 =3.1189723326632947\times10^{-5}.
\]
Equivalentemente, la compatibilidad entre espectros y mezcla puede leerse en el determinante del conmutador de los operadores hermíticos up/down. Un valor no nulo de \(J_{\rm CKM}\) registra orientación física no eliminable; no es un corrector de masa añadido a cada quark.

\Needspace{7\baselineskip}
\section{Sectores de calibre y escalar}
\Needspace{7\baselineskip}
\subsection{Carta del sector nulo: rutas de masa \(0\), representaciones de calibre y persistencia}
Sea

\[
 \mathcal A_{\rm sp}=\mathbb R[R]/(R^2-1),
 \qquad
 P_\pm=\frac{1\pm R}{2}.
\]
Para

\[
 x=r_+P_++r_-P_-,
\]
la norma partida es

\[
 N_{\rm sp}(x)=r_+r_-.
\]
La sección nula es una dirección de norma cero. Su masa es exactamente

\[
 m=0,
\]
sin pasar por un límite de caracteres positivos y sin asignar \(\kappa=0\). La nulidad pertenece al tipo de la sección.

\Needspace{7\baselineskip}
\subsection{Realización del fotón en el sector nulo electromagnético: helicidad, calibre y carácter}
El fotón es la realización vectorial transversal de una sección de calibre nula. El proyector físico debe dejar exactamente dos polarizaciones:

\[
 \operatorname{rg}P_\gamma^{\rm trans}=2,
\qquad
 P_\gamma^{\rm trans}k=0.
\]
La masa de reposo es

\[
 m_\gamma=0.
\]
La propagación bidireccional puede portar una anisotropía orientada

\[
 \epsilon=q_-^{90}-q_+^{90},
\qquad
 c_\pm=\frac{c}{1\pm\epsilon},
\]
mientras el promedio armónico conserva

\[
 \frac{2}{1/c_++1/c_-}=c.
\]
Así, la ida y la vuelta pueden diferir sin alterar el invariante two--way. Este sesgo no es una masa de fotón.

Una realización Proca/Stückelberg distinta aparece si un brecha espectral del registro genealógico satisface

\[
 \omega^2=k^2+m_T^2
\]
y existe un entrelazador dimensional

\[
 m_T\longmapsto m_\gamma.
\]
La demora inducida obedecería

\[
 \frac{\Delta T}{T}
 \simeq\frac12
 \left(\frac{m_\gamma c^2}{\hbar\omega}\right)^2.
\]
El sector Maxwell nulo y la realización Proca masiva son dos ramas distintas; una no invalida a la otra.

\Needspace{7\baselineskip}
\subsection{Realización cromodinámica del sector nulo: octete de gluones y representación adjunta de \(SU(3)\)}
El trit residual genera el espacio de color, cuya parte sin traza realiza

\[
 \mathfrak{sl}_3\longrightarrow\mathfrak{su}_3.
\]
El gluón ocupa la representación adjunta de dimensión ocho. Mientras la simetría de color permanezca exacta,

\[
 m_g=0.
\]
Confinamiento, apantallamiento y brecha espectral colectivo son propiedades del espectro interactivo; no constituyen una masa de polo elemental del gluón. Una masa efectiva debe indicar observable, calibre, régimen y operador que la define.

\Needspace{7\baselineskip}
\subsection{Realización de los bosones \(W^\pm\): carga, polarización y carácter electrodébil}
El sector \(W\) es un bloque vectorial cargado con conjugación

\[
 C_M:W^+\longleftrightarrow W^-.
\]
Su huella histórica es

\[
 15/37,\qquad \text{eje este--oeste},\qquad
 (n_A,s_C)=(94,-1),
\]
y la firma enriquecida declarada es

\[
 (94,-1,0,-2,4).
\]
La evaluación correspondiente es

\[
 80.37896490970455\ldots\ {\rm GeV}.
\]
La realización completa es el polo del resolvente del bloque \(P_W\widehat M P_W\); su parte imaginaria determina la anchura. La degeneración de \(W^+\) y \(W^-\) procede de la conjugación del operador.

\Needspace{7\baselineskip}
\subsection{Realización del bosón \(Z^0\): mezcla neutra, polarización y carácter electrodébil}
El \(Z^0\) es el bloque vectorial neutral del acoplamiento electrodébil. Su huella histórica es

\[
 20/23,\qquad \text{eje este--oeste},\qquad
 (n_A,s_C)=(95,-1),
\]
y la firma enriquecida

\[
 (95,-1,-11,0,-4).
\]
La evaluación del carácter da

\[
 91.18753344431341\ldots\ {\rm GeV}.
\]
El proyector neutral debe construirse antes de la carta de polo; masa y anchura son dos coordenadas del mismo polo, no una única cifra.

\Needspace{7\baselineskip}
\subsection{Realización del sector de Higgs: modo escalar, acoplamiento y carácter de masa}
El Higgs es una sección escalar persistente del bloque de calibre con paridad CP declarada. Su huella es

\[
 24/29,\qquad \text{eje norte--sur},\qquad
 (n_A,s_C)=(97,9).
\]
La evaluación condicionada del carácter produce

\[
 125.2924660425361\ldots\ {\rm GeV}.
\]
La definición HMT no identifica “escalar” con “masa”: el espín cero fija la representación rotacional; la masa procede del espectro del bloque escalar acoplado. La anchura y los canales de decaimiento requieren el operador temporal.

\Needspace{7\baselineskip}
\subsection{El momento magnético anómalo del muón}
El funcional interno

\[
 a_\mu=
 \frac{\alpha_{\rm HMT}}{2\pi}
 +\frac{\alpha_{\rm HMT}(\varphi-1)}{1000}
 +\frac{\alpha_{\rm HMT}^2}{1000(54+1)}
\]
da

\[
 a_\mu=0.001165920713008014\ldots,
\]
y

\[
 g_\mu=2(1+a_\mu)
 =2.002331841426016028\ldots.
\]
Los tres términos corresponden, respectivamente, al giro basal, al retorno orientado en la completación \(1000\) y a la corrección cuadrática mediada por el defecto \(54\). Este observable comprueba la estructura interna del sector muónico, pero no selecciona por sí solo su ruta de masa.

\Needspace{7\baselineskip}
\section{Construcción de partículas compuestas mediante composición no aditiva de rutas y enlaces}
\Needspace{7\baselineskip}
\subsection{Operador de composición}
La composición de rutas no es estrictamente aditiva porque frontera, orden temporal y acarreo interactúan. Para dos registros,

\[
 \mathcal L_1\star\mathcal L_2
 =\mathcal L_1+\mathcal L_2
 +\mathfrak b(\mathcal L_1,\mathcal L_2),
\]
donde \(\mathfrak b\) es el cociclo de ligadura. La asociatividad exige

\[
 \mathfrak b(L_1,L_2)
 +\mathfrak b(L_1\star L_2,L_3)
 =
 \mathfrak b(L_2,L_3)
 +\mathfrak b(L_1,L_2\star L_3).
\]
El lector de masa actúa después de esta composición:

\[
 (\Gamma_i,\mathcal T,\partial)
 \xrightarrow{\mathfrak C_{12}}
 \mathcal L_{\rm comp}
 \xrightarrow{L_{\rm tick}}
 \sigma_{\rm comp}
 \xrightarrow{\chi_{\rm full}}
 \widehat M_{\rm comp}.
\]
\Needspace{7\baselineskip}
\subsection{Realización mesónica mediante composición quark–antiquark y ley cocíclica del enlace}
El dominio mesónico contiene 432 expresiones orientadas

\[
 M(q,q')=q\otimes(q')^\vee
\]
con color compatible. El proyector físico impone singlete de color y fija

\[
 (J^{PC},I,I_3,S,C,B,\ldots),
\]
junto con el árbol temporal de enlace. La masa de un mesón no es \(m_q+m_{\bar q}\); es un autovalor del registro compuesto.

Las seis evaluaciones históricas de control incluyen

\[
 m_{\pi^0}=134.9767243278721\ldots\ {\rm MeV},
\]
\[
 m_{\pi^\pm}=139.5704851715722\ldots\ {\rm MeV},
\]
\[
 m_{K^\pm}=493.6770856495306\ldots\ {\rm MeV},
\]
\[
 m_{K^0}=497.6111864977505\ldots\ {\rm MeV}.
\]
Estas cifras son evaluaciones exactas de registros declarados. El cierre genealógico estado por estado requiere regenerar cada registro mediante \(\mathfrak C_{12}\), incluida su frontera de ligadura.

\Needspace{7\baselineskip}
\subsection{Realización bariónica mediante composición trilineal de rutas quark y neutralización de color}
El dominio bariónico contiene 1728 expresiones RGB. El proyector físico debe imponer simultáneamente:

\begin{enumerate}
\item singlete de color;
\item estadística fermiónica;
\item sabor e isospín;
\item espín y simetrización;
\item árbol triple con acarreo y memoria.
\end{enumerate}
Para protón y neutrón, las firmas históricas

\[
 \sigma_p=(59,0,9,1,0),\qquad
 \sigma_n=(59,1,-34,0,1)
\]
dan

\[
 m_p=938.2717515469849\ldots\ {\rm MeV},
\]
\[
 m_n=939.5658104725070\ldots\ {\rm MeV}.
\]
La precisión aparente no sustituye la derivación del registro genealógico de ligadura. El signo y la magnitud de

\[
 m_n-m_p
\]
deben salir de orientación, isospín, vacancias y frontera de enlace antes de abrir la comparación externa.

\Needspace{7\baselineskip}
\subsection{Estados excitados y resonancias hadrónicas}
Una identidad compuesta puede poseer varios polos. El estado físico queda individualizado por

\[
 (\text{constituyentes},\mathcal T,\partial,
 J^{PC},I,\text{radial},\text{orbital},\text{canal}).
\]
La misma composición de sabores no basta para distinguir excitaciones. Cada polo

\[
 z_k=M_k-\frac{i}{2}\Gamma_k
\]
debe proceder de un autovector del operador compuesto y de un canal de decaimiento. Esto extiende la ley desde unas pocas masas estables a todo el espectro de mesones y bariones sin convertir el catálogo externo en generador.

\Needspace{7\baselineskip}
\subsection{Inventario exhaustivo de anchuras y tasas de transición de las realizaciones del atlas}
Sea \(P\) el subespacio preparado y \(Q=I-P\). El Hamiltoniano efectivo es

\[
 H_{\rm eff}(z)
 =PHP+PHQ(z-QHQ)^{-1}QHP.
\]
Los polos de

\[
 \det(z-H_{\rm eff}(z))=0
\]
determinan masas y anchuras. El carácter de masa aporta la parte real basal del operador; el término de autoenergía de \(Q\) aporta desplazamiento y decaimiento. Por ello, las 384 anchuras del inventario requieren un lector de canales adicional al lector de las 471 masas.

\Needspace{7\baselineskip}
\section{Sectores topológicos y partículas hipotéticas}
\Needspace{7\baselineskip}
\subsection{Sector topológico de Fibonacci: anillo basado, fusión y representación de trenzas}
La incidencia areal--volumétrica

\[
 F_{\rm av}=
 \begin{pmatrix}
 0&1\\1&1
 \end{pmatrix}
\]
realiza el anillo

\[
 \tau^2=1+\tau,\qquad
 \operatorname{FPdim}\tau=\varphi.
\]
La salida primaria es una regla de fusión y una dimensión cuántica. Una partícula anyónica de Fibonacci exige además matrices \(F\) y \(R\), pentágono, hexágono, quiralidad, Hamiltoniano con brecha espectral y un mapa desde el estado TPK al espacio físico. Sólo entonces puede definirse

\[
 m_{\rm brecha espectral}=\frac{\Delta E}{c_{\rm eff}^2}.
\]
\Needspace{7\baselineskip}
\subsection{Sector de Ising–Majorana: objetos simples, reglas de fusión y representación de trenzas}
La categoría de Ising satisface

\[
 \psi^2=1,\qquad
 \psi\sigma=\sigma,\qquad
 \sigma^2=1+\psi,\qquad
 d_\sigma=\sqrt2.
\]
Deben distinguirse cuatro niveles:

\begin{enumerate}
\item la involución de rutas \(C_M=-\operatorname{rev}\);
\item la paridad de la completación CAR;
\item un espinor relativista autorconjugado;
\item un modo cero topológico de Ising.
\end{enumerate}
Un fermión Majorana relativista requiere una representación espinorial real, una dinámica y un término de masa. Un modo cero de Majorana se caracteriza por brecha espectral, splitting y protección; no posee una masa universal deducible sólo del anillo de fusión. La identidad de una ruta con su conjugada es condición necesaria, pero no suficiente, para la realización Majorana.

\Needspace{7\baselineskip}
\subsection{Representación de trenzas sobre \(\mathbb Z/6\mathbb Z\) y caracteres topológicos}
Sea

\[
 w_6:B_n\longrightarrow\mathbb Z/6,
\qquad
 \chi_k(b)=e^{2\pi ikw(b)/6}.
\]
Las representaciones

\[
 \rho_q(\sigma_i)=qP_i
\]
descienden al grupo simétrico únicamente cuando \(q^2=1\). Los otros cuatro valores son genuinamente trenzados. Su existencia matemática no implica cuatro partículas: una realización física requiere Hamiltoniano, brecha espectral, operadores de trenza, estabilidad y entrelazador con una multisección.

\Needspace{7\baselineskip}
\subsection{Imposibilidad de un taquión persistente}
La recta de acción es positiva,

\[
 \hbar_{\rm int}>0,
\]
el carácter satisface

\[
 \chi_{\rm full}>0,
\]
y el operador beta se obtiene por conjugación positiva:

\[
 \widehat M^\beta
 =R_{\rm act}^{\widehat K/2}
 \widehat M^{(0)}
 R_{\rm act}^{\widehat K/2}\ge0.
\]
Por tanto,

\[
 \operatorname{Spec}(\widehat M^\beta)\subset[0,\infty).
\]
Una sección persistente estable no puede generar

\[
 m^2<0.
\]
Si la linealización de un estado produce curvatura negativa o frecuencia imaginaria, ese estado es una inestabilidad que abandona la sección, no una partícula superlumínica. Si la prolongación termina, es una sección virtual con obstrucción terminal. Así, “taquión” se descompone en dos objetos HMT bien tipados ---inestabilidad o virtualidad--- y no constituye una especie persistente.

\Needspace{7\baselineskip}
\subsection{Gravedad sin partícula elemental obligatoria de espín \(2\)}
La realización gravitatoria sigue la cadena

\[
 \text{holonomía de ruta}
 \longrightarrow(e,\omega)
 \longrightarrow(T,R)
 \longrightarrow
 S_{\rm PCH},
\]
donde \(e\) es la coestructura, \(\omega\) la conexión,

\[
 T=de+\omega\wedge e,\qquad
 R=d\omega+\omega\wedge\omega,
\]
y \(S_{\rm PCH}\) es la acción de Palatini--Cartan--Holst. La gravedad aparece como respuesta global de torsión y curvatura de la red persistente. La vibración colectiva del universo modelo puede poseer una componente tensorial de espín dos, pero HMT no necesita postular una partícula elemental adicional para transportar la gravedad.

La constante gravitatoria interna se tipa mediante

\[
 G_{\rm HMT}
 =\frac{c_{\rm int}^3}{\hbar_{\rm int}}
 \bigl(\delta_\theta\Lambda_G\bigr)^2,
\]
donde \(\delta_\theta\) es la torsión angular elemental y \(\Lambda_G\) la longitud espectral seleccionada por la holonomía gravitatoria. Esta ecuación convierte torsión y escala de ruta en acoplamiento gravitatorio sin usar una masa como normalización. La vibración total se representa por el Hessiano

\[
 \mathcal H_{\rm univ}
 =\operatorname{Hess}_{(e,\omega)}S_{\rm PCH}
\]
sobre el estado global. Sus modos son excitaciones colectivas del universo modelo; una componente tensorial no es automáticamente una partícula separada.

Una excitación que se denomine “gravitón” debe demostrar tres cosas:

\begin{enumerate}
\item un subespacio tensorial persistente;
\item un polo aislado del operador linealizado;
\item un acoplamiento físico que no duplique la torsión ya presente.
\end{enumerate}
Sin esas tres condiciones, el objeto correcto es el modo colectivo de la geometría, no una nueva fila de masa.

\Needspace{7\baselineskip}
\subsection{Barbero–Immirzi y el enlace hexada–octada}
Con

\[
 A=1000\alpha_{\rm HMT},
\qquad
 q_\pm=\exp\!\left[-\frac{\pi}{180}(A\pm C^*)\right],
\]
el funcional hexada--octada usa dos incidencias distintas de las torres temporales. La combinación

\[
 \varepsilon_{\rm BI}
 =12\Delta S_{90}-\Delta S_{120}
\]
y el término angular producen

\[
 \gamma_{\rm BI}
 =\frac{\pi}{180}AC^*
 +\varepsilon_{\rm BI}
 =0.2740076726944593\ldots.
\]
El peso \(12:-1\) procede de la lectura de los doce sectores dodecafásicos
y de la única costura orientada de retorno. La ecuación
\(4a-3b=45\) verifica posteriormente que el par \((12,1)\) es la solución
entera positiva mínima. Este objeto no debe confundirse con el extractor dodecafásico de masa ni con el residuo angular \(90\)--\(120\); comparte el operador armónico, pero posee otro dominio y otra función.

\Needspace{7\baselineskip}
\subsection{Extensiones del atlas hacia sectores oscuros: dominios, lectores y estatuto físico}
Un sector oscuro HMT es una multisección persistente con

\[
 P_{\rm vis}P_{\rm dark}=0
\]
para los acoplamientos electromagnético y cromático, pero con posible incidencia gravitatoria o de frontera:

\[
 P_{\rm grav}P_{\rm dark}\ne0.
\]
Su “oscuridad” es una propiedad del núcleo de acoplamiento, no una etiqueta de masa. La ley predice primero un operador y su espectro; sólo un canal físico determina si el modo es estable, metaestable o resonante. Esta definición permite buscar multisecciones invisibles sin inventar una masa mediante el valor observado de una densidad cosmológica.

\Needspace{7\baselineskip}
\subsection{Sección neutrínica estéril: compatibilidad de fibra, mezcla y ausencia de carga de calibre}
Un neutrino estéril no es la profundidad \(G_4\). Es un autovector \(n_s\in\mathcal H_\nu\) que satisface

\[
 P_{\rm weak}n_s=0,\qquad
 P_{\rm mass}n_s=n_s.
\]
Puede mezclarse con los tres modos activos mediante un núcleo neutral ampliado, pero su existencia requiere

\[
 \dim\ker(P_{\rm weak}|_{\mathcal H_\nu})>0
\]
y un autovalor de masa en ese núcleo. Esta definición produce un criterio de búsqueda interno: diagonalizar primero el bloque neutral y comprobar después su acoplamiento débil.

\Needspace{7\baselineskip}
\subsection{Secciones axiónica y pseudoescalares: carácter topológico y acoplamiento de borde}
Un candidato axiónico es una sección pseudoescalar cuya orientación cambia bajo CP,

\[
 CP\,a\,CP^{-1}=-a,
\]
y cuyo desplazamiento compensa una fase topológica del sector de color. En HMT debe corresponder a un modo de torsión/holonomía impar y a un proyector que entrelace ese modo con la densidad topológica de calibre. Su masa se obtiene del segundo derivado del potencial efectivo,

\[
 m_a^2=
 \left.\frac{\partial^2V_{\rm eff}}{\partial a^2}\right|_{a=a_0},
\]
una vez que \(V_{\rm eff}\) haya sido construido desde el registro genealógico de vacancias y frontera. La mera existencia de un canal \(C^*\mapsto-C^*\) no basta para afirmar una partícula axiónica.

\Needspace{7\baselineskip}
\subsection{Fotón oscuro y vector masivo oculto}
Un fotón oscuro requiere una segunda sección de calibre \(A'\) y un acoplamiento de mezcla

\[
 \mathcal L_{\rm mix}
 =-\frac{\varepsilon_{\rm kin}}2F_{\mu\nu}F'^{\mu\nu}.
\]
En términos de fibras, \(\varepsilon_{\rm kin}\) es un entrelazador entre dos subespacios nulos. Si el segundo sector adquiere un brecha espectral Stückelberg o de frontera,

\[
 \widehat M_{A'}^2\ne0,
\]
aparece un vector masivo oscuro. HMT predice que el parámetro de mezcla debe proceder de una incidencia de memoria común; no puede fijarse con la masa que se pretende explicar.

\Needspace{7\baselineskip}
\subsection{WIMP y estado oscuro estable}
Un WIMP es una sección oscura que además cumple:

\[
 [P_{\rm weak},P_X]\ne0,\qquad
 [P_{\rm em},P_X]=[P_{\rm color},P_X]=0,
\]
y posee un modo estable

\[
 |\lambda_X|=1.
\]
La estabilidad puede estar protegida por una paridad del registro genealógico. Su masa es un autovalor del bloque oscuro acoplado débilmente. “WIMP” no constituye una nueva familia matemática: es una condición conjunta de acoplamiento y persistencia sobre una multisección.

\Needspace{7\baselineskip}
\subsection{Sección monopolar: carga magnética, condición de Dirac y compatibilidad topológica}
Un monopolo se define por una clase topológica no trivial del fibrado de calibre:

\[
 \frac1{2\pi}\int_{S^2}F=n\ne0.
\]
HMT debe realizar este entero como holonomía neta de una familia cerrada de rutas. Su masa, si el modo es dinámico y estable, pertenece al espectro del operador de defecto topológico. Una vacancia puntual o un acarreo no nulo no son por sí solos un monopolo; debe demostrarse el descenso a la clase integral.

\Needspace{7\baselineskip}
\subsection{Extensiones supersimétricas del atlas: emparejamiento de representaciones y caracteres}
Una supersimetría HMT requeriría un operador impar \(Q\) que intercambie las completaciones exterior y simétrica:

\[
 Q:\mathcal F_-\leftrightarrow\mathcal F_+,
\qquad
 Q^2=H,
\]
y que entrelace la ley de masas,

\[
 Q\widehat M_-=\widehat M_+Q.
\]
Sólo entonces los espectros bosónico y fermiónico quedarían emparejados. La existencia separada de CAR y CCR no demuestra supersimetría ni predice por sí misma selectrones, neutralinos o gluinos. Cada compañero exige un entrelazador no nulo y una regla de ruptura que determine su separación de masa.

\Needspace{7\baselineskip}
\subsection{Tetraquarks, pentaquarks, glueballs e híbridos}
La gramática de composición se prolonga sin cambiar de ley:

\[
 \mathfrak C_{12}^{(n)}:
 \Gamma_1\times\cdots\times\Gamma_n
 \times\mathcal T_n\times\partial_n
 \longrightarrow\mathcal L_{\rm comp}.
\]
Un tetraquark usa \(qq\bar q\bar q\); un pentaquark, \(qqqq\bar q\); un glueball, registros de la adjunta de color proyectados al singlete; un híbrido, rutas quark y de calibre en un mismo árbol. La clasificación física depende del proyector de color, espín, paridad y árbol de enlace, no sólo de la lista de constituyentes.

\Needspace{7\baselineskip}
\subsection{Secciones de dilatón y radión: modos escalares de volumen y compactificación}
El dilatón es un modo de la sección de escala; el radión, un modo que modifica una longitud o volumen interno. En la teoría de masa aparecen como fluctuaciones de la carta y de la memoria geométrica:

\[
 t_0\mapsto e^{\phi_d}t_0,\qquad
 \mathcal V_{\rm int}\mapsto e^{\phi_r}\mathcal V_{\rm int}.
\]
Para convertirse en partículas deben poseer término cinético, potencial, proyector persistente y polo. Mientras no los posean son grados colectivos de la carta dimensional, no especies escalares adicionales.

\Needspace{7\baselineskip}
\subsection{Exclusión de fantasmas físicos}
La medida espectral de un estado físico es positiva:

\[
 \mu_{P,a}^{\rm mass}(B)\ge0.
\]
Una dirección de norma negativa no pertenece al espacio de estados preparado. Si una parametrización introduce un “fantasma”, la realización física debe cocientar la redundancia de calibre o rechazar el acoplamiento. Esta exclusión es independiente de la prohibición taquiónica: un fantasma viola positividad de norma; un taquión persistente violaría positividad del operador de masa.

\Needspace{7\baselineskip}
\section{Inventario universal de masas y anchuras}
\Needspace{7\baselineskip}
\subsection{Correspondencia posterior entre las realizaciones HMT–MD y el censo físico externo}
El inventario de contraste contiene

\[
 1170\ \text{identidades},\qquad
 438\ \text{grupos de identidad},
\]
\[
 471\ \text{observables de masa},\qquad
 384\ \text{observables de anchura}.
\]
Su partición es

\Needspace{7\baselineskip}
\begingroup
\footnotesize
\setlength{\tabcolsep}{2.1pt}
\begin{longtable}{L{5.04cm}L{5.04cm}L{5.04cm}}
\toprule
\rowcolor{HMTNavy}\HMTtablehead{Dominio} & \HMTtablehead{Masas} & \HMTtablehead{Anchuras} \\
\midrule
\endfirsthead
\multicolumn{3}{l}{\sffamily\scriptsize\color{HMTGray}Continuación de la tabla}\\[-1pt]
\toprule
\rowcolor{HMTNavy}\HMTtablehead{Dominio} & \HMTtablehead{Masas} & \HMTtablehead{Anchuras} \\
\midrule
\endhead
\midrule
\endfoot
\bottomrule
\endlastfoot
partículas generales y leptónicas & 53 & 9 \\
\rowcolor{HMTLight}quarks & 8 & 1 \\
mesones & 243 & 225 \\
\rowcolor{HMTLight}bariones & 166 & 149 \\
registro gravitatorio de contraste & 1 & 0 \\
\rowcolor{HMTLight}\textbf{Total} & \textbf{471} & \textbf{384} \\
\end{longtable}
\endgroup
Las 471 masas no son 471 especies elementales: incluyen estados excitados, resonancias, distintas cargas y múltiples realizaciones compuestas. Las 384 anchuras son observables temporales y no deben ser leídas por el carácter de masa.

\Needspace{7\baselineskip}
\subsection{Identidad interna \(\mathcal I_X\) de una realización material y criterio de individualización frente a observaciones múltiples}
Cada estado comparado debe poseer una identidad interna

\[
 \mathcal I_X=
 (\mathcal C_X,\mathcal R_X,\Gamma_X,\mathcal L_X,
 \sigma_X,\eta_X,P_X,\mathfrak M_X,\mathcal T_X),
\]
donde:

\begin{itemize}
\item \(\mathcal C_X\) es el tipo físico de realización determinado por objeto, representación y codominio;
\item \(\mathcal R_X\) es la representación física;
\item \(\Gamma_X\) es una ruta, órbita o composición orientada;
\item \(\mathcal L_X\) es el registro genealógico;
\item \(\sigma_X\) es la firma completa;
\item \(\eta_X\) es la torre;
\item \(P_X\) es el proyector del acoplamiento;
\item \(\mathfrak M_X\) es el canal de medición;
\item \(\mathcal T_X\) es el árbol de composición cuando procede.
\end{itemize}
Dos observaciones externas pertenecen al mismo estado HMT cuando comparten esta identidad y difieren sólo por experimento, carta o estimador. No se crean dos partículas porque existan dos mediciones.

\Needspace{7\baselineskip}
\subsection{Morfismo de comparación entre invariantes internos y observables metrológicos}
La comparación requiere el diagrama

\[
 \mathcal L_M^{\rm HMT}
 \xrightarrow{\;\varsigma_{\mathcal S,\mu}\;}
 \mathbb R_{>0}u_M^{\mathcal S,\mu}
 \xleftarrow{\;\text{medición}\;}
 \mathcal O_X^{\rm ext}.
\]
Una fila científicamente completa contiene:

\[
 (m_X^{\rm HMT},m_X^{\rm ext},
 u,\mathcal S,\mu,\Sigma_{\rm exp},\Sigma_{\rm HMT},
 \text{convención de polo}).
\]
Para quarks son obligatorios esquema y escala; para resonancias, la convención de polo; para límites, el nivel de confianza; para observables correlacionados, la matriz de covarianza.

\Needspace{7\baselineskip}
\subsection{Evaluación explícita de \(17\) realizaciones representativas del atlas}
Las siguientes evaluaciones conservan la separación entre ley y dato de firma:

\Needspace{7\baselineskip}
\begingroup
\footnotesize
\setlength{\tabcolsep}{2.1pt}
\begin{longtable}{L{4.68cm}L{10.61cm}}
\toprule
\rowcolor{HMTNavy}\HMTtablehead{Estado} & \HMTtablehead{Valor del carácter en su carta} \\
\midrule
\endfirsthead
\multicolumn{2}{l}{\sffamily\scriptsize\color{HMTGray}Continuación de la tabla}\\[-1pt]
\toprule
\rowcolor{HMTNavy}\HMTtablehead{Estado} & \HMTtablehead{Valor del carácter en su carta} \\
\midrule
\endhead
\midrule
\endfoot
\bottomrule
\endlastfoot
\(\mu\) & \(105.658360\allowbreak{}294435\allowbreak{}829\ldots\ {\rm MeV}\) \\
\rowcolor{HMTLight}\(\tau\) & \(1776.860917\allowbreak{}631432\allowbreak{}3\ldots\ {\rm MeV}\) \\
\(u\) & \(2.165924\allowbreak{}536173\allowbreak{}907\ldots\ {\rm MeV}\) \\
\rowcolor{HMTLight}\(d\) & \(4.679550\allowbreak{}162948\allowbreak{}874\ldots\ {\rm MeV}\) \\
\(c\) & \(1.270500\allowbreak{}838641\allowbreak{}911\ldots\ {\rm GeV}\) \\
\rowcolor{HMTLight}\(s\) & \(92.940093\allowbreak{}907845\allowbreak{}64\ldots\ {\rm MeV}\) \\
\(t\) & \(172.597479\allowbreak{}544019\allowbreak{}1\ldots\ {\rm GeV}\) \\
\rowcolor{HMTLight}\(b\) & \(4.190119\allowbreak{}763299\allowbreak{}790\ldots\ {\rm GeV}\) \\
\(W\) & \(80.378964\allowbreak{}909704\allowbreak{}55\ldots\ {\rm GeV}\) \\
\rowcolor{HMTLight}\(Z\) & \(91.187533\allowbreak{}444313\allowbreak{}41\ldots\ {\rm GeV}\) \\
\(H\) & \(125.292466\allowbreak{}042536\allowbreak{}1\ldots\ {\rm GeV}\) \\
\rowcolor{HMTLight}\(\pi^0\) & \(134.976724\allowbreak{}327872\allowbreak{}1\ldots\ {\rm MeV}\) \\
\(\pi^\pm\) & \(139.570485\allowbreak{}171572\allowbreak{}2\ldots\ {\rm MeV}\) \\
\rowcolor{HMTLight}\(K^\pm\) & \(493.677085\allowbreak{}649530\allowbreak{}6\ldots\ {\rm MeV}\) \\
\(K^0\) & \(497.611186\allowbreak{}497750\allowbreak{}5\ldots\ {\rm MeV}\) \\
\rowcolor{HMTLight}\(p\) & \(938.271751\allowbreak{}546984\allowbreak{}9\ldots\ {\rm MeV}\) \\
\(n\) & \(939.565810\allowbreak{}472507\allowbreak{}0\ldots\ {\rm MeV}\) \\
\end{longtable}
\endgroup
El electrón beta no pertenece a esta tabla porque su ruta central, sus coeficientes basales y su lector completo están determinados por la propia construcción.

\Needspace{7\baselineskip}

\endgroup
\begingroup
\let\chapter\section
\let\section\subsection
\let\subsection\subsubsection
\let\subsubsection\paragraph
\let\clearpage\relax
\subsection{Del catálogo de realizaciones al inventario de 471 masas}
La extensión universal sigue cinco operaciones:

\[
\begin{aligned}
 \text{identidad física}
 &\longrightarrow
 \text{proyector de representación}
 \longrightarrow
 \text{registro genealógico elemental o compuesto}\\
 &\longrightarrow
 \text{operador/polo}
 \longrightarrow
 \text{carta de comparación}.
\end{aligned}
\]
Para estados elementales, el segundo mapa actúa sobre una fibra del atlas. Para mesones y bariones, actúa sobre las 432 y 1728 expresiones compuestas. Para resonancias, el cuarto mapa devuelve un polo complejo. Para sectores topológicos, devuelve primero un brecha espectral y una representación de trenza. Esta tipificación permite abarcar todo el inventario sin declarar que cada nombre externo corresponde a una celda elemental.

\Needspace{7\baselineskip}
\subsection{De las 384 anchuras al operador de supervivencia}
Cada anchura requiere

\[
 X\longmapsto Q_XUQ_X
 \longmapsto \rho(Q_XUQ_X)
 \longmapsto\Gamma_X
\]
o, equivalentemente, el polo de \(H_{\rm eff}\). El lector debe incorporar canales de decaimiento, multiplicidad, umbral y espacio de fases. La igualdad

\[
 \Gamma_{\bar X}=\Gamma_X
\]
se deduce sólo si el operador de decaimiento es covariante bajo conjugación.

\Needspace{7\baselineskip}
\section{Masa, temperatura y estadística}
\Needspace{7\baselineskip}
\subsection{Transductor térmico entre carácter de masa, temperatura y ocupación estadística}
La constante de Boltzmann es el transductor

\[
 k_B:\mathcal L_\Theta\longrightarrow\mathcal L_E.
\]
El reloj CT--108 fija la identidad

\[
 k_B\Theta_{\rm clk}\log3
 =\hbar_{\rm int}\frac{2\pi}{108t_0}.
\]
Esta ecuación vincula una unidad de información trítica, una frecuencia y una energía. No introduce la coordenada SI de \(k_B\) en la ley; determina el producto tipado una vez elegida la carta térmica.

\Needspace{7\baselineskip}
\subsection{Proyección fermiónica del atlas y distribución de Fermi–Dirac}
Para niveles

\[
 \epsilon=(0,3,6,9),\qquad
 g=(1,6,12,8),
\]
y restricciones

\[
 (N,E)=(12,66),
\]
la completación exterior produce

\[
 2\,104\,494
\]
microestados fermiónicos compatibles. Los multiplicadores efectivos son

\[
 (\beta,\mu)
 =(0.1530701135,\;4.5028651976).
\]
Esta realización ejemplifica cómo masa, energía y exclusión emergen de rutas fermiónicas sin que la estadística sea un ajuste adicional.

\Needspace{7\baselineskip}
\subsection{Proyección bosónica del atlas y distribución de Bose–Einstein}
Sobre los mismos niveles y restricciones, la completación simétrica produce

\[
 259\,436\,430
\]
microestados, con

\[
 (\beta,\mu)
 =(0.0545672796,\;-15.9256978191).
\]
La diferencia entre Fermi--Dirac y Bose--Einstein es, por tanto, una diferencia en la regla de composición de rutas, no en el valor de una masa individual.

\Needspace{7\baselineskip}
\subsection{Criterio de condensación de las secciones bosónicas y ocupación macroscópica}
Un condensado bosónico aparece cuando una componente simétrica adquiere ocupación macroscópica en el autovalor mínimo del operador. Un mar fermiónico, en cambio, llena un conjunto de autovalores distintos hasta una frontera. HMT describe ambos mediante el mismo espectro de masa/energía y dos funtores de Fock distintos. Esta separación impide identificar “muchas partículas en una ruta” con una partícula más masiva.

\Needspace{7\baselineskip}
\section{Propagación diferencial de incertidumbres y estabilidad estructural del carácter de masa}
\Needspace{7\baselineskip}
\subsection{Invariancia del grafo generativo bajo el contraste metrológico posterior}
La salida interna se congela en el orden

\[
 \mathcal R_{324}
 \to P_X
 \to\mathcal L_X
 \to\sigma_X
 \to\eta_X
 \to\widehat M_X
 \to\mu_X.
\]
Sólo después se abre la carta externa. Una permutación arbitraria de las 471 masas o de las 384 anchuras debe dejar invariantes todos los objetos anteriores.

\Needspace{7\baselineskip}
\subsection{Descomposición de la incertidumbre estructural en contribuciones experimental, de carta, ruta, torre y acoplamiento}
La incertidumbre total se descompone como

\[
 \Sigma_{\rm tot}
 =\Sigma_{\rm ext}
 +\Sigma_{\rm carta}
 +\Sigma_{\rm route}
 +\Sigma_{\rm tower}
 +\Sigma_{\rm coupling}.
\]
Aquí:

\begin{itemize}
\item \(\Sigma_{\rm ext}\) es la covarianza experimental;
\item \(\Sigma_{\rm carta}\) cuantifica esquema, escala y conversión de unidades;
\item \(\Sigma_{\rm route}\) compara todas las rutas de la órbita admisible;
\item \(\Sigma_{\rm tower}\) usa la cota de cola de la serie armónica;
\item \(\Sigma_{\rm coupling}\) mide la variación entre entrelazadores permitidos.
\end{itemize}
No se suman necesariamente como números independientes; la ecuación expresa una descomposición de fuentes que debe convertirse en una covarianza conjunta.

\Needspace{7\baselineskip}
\subsection{Derivada del observable de masa respecto de la firma y de los lectores}
Para

\[
 \log m_X=
 \log m_\star+
 \log\chi_0(\sigma_X)
 +\kappa_X\log R_{\rm act}
 +\sum_h\eta_{X,h}s_h,
\]
la variación es

\[
 d\log m_X
 =d\log m_\star
 +\langle\Theta_0,d\sigma_X\rangle
 +\langle\sigma_X,d\Theta_0\rangle
 +\log R_{\rm act}\,d\kappa_X
 +\kappa_X\,d\log R_{\rm act}
 +\sum_h(s_h\,d\eta_{X,h}+\eta_{X,h}\,ds_h).
\]
Esta fórmula permite atribuir cada decimal a ruta, constante, lector, torre o carta. Un resultado es robusto cuando perturbaciones que preservan las premisas mantienen su clase espectral y su incertidumbre calculada.

\Needspace{7\baselineskip}
\subsection{Ensayos de dependencia estructural mediante supresión controlada de hoja, orientación, acarreo, memoria, devanado y frontera}
Una prueba de dependencia elimina por separado hoja, orientación, acarreo, memoria, devanado o frontera y vuelve a calcular el objeto. El resultado esperado no es que todo cambie, sino que cambien exactamente las coordenadas dependientes de la pieza eliminada. Si suprimir memoria deja invariante un coeficiente definido por retorno, la construcción de ese coeficiente queda refutada.

\Needspace{7\baselineskip}
\subsection{Obstrucciones a la identificación de rutas, caracteres y magnitudes metrológicas}
La ley queda refutada en una aplicación concreta si ocurre cualquiera de estos hechos:

\begin{enumerate}
\item una masa externa cambia la ruta o el proyector elegido;
\item la firma publicada no se regenera desde el registro genealógico;
\item el refinamiento no converge o depende de una truncación elegida por el residual;
\item dos rutas equivalentes reciben masas distintas sin ruptura de simetría;
\item una antipartícula adquiere masa distinta sin violación declarada de conjugación;
\item un triplete de color se escinde sin ruptura de color;
\item se publica un escalar aunque la varianza espectral sea positiva;
\item una anchura se obtiene del carácter de masa sin operador temporal;
\item una carta quark omite esquema o escala;
\item una inestabilidad se interpreta como taquión persistente.
\end{enumerate}
\Needspace{7\baselineskip}

\endgroup

\section{Leptones, neutrinos, quarks y sectores de calibre}
Los leptones cargados, los seis sabores quark y las fibras neutrínicas se
realizan como multisecciones tipadas. El color es la órbita residual ternaria
\(12+12+12\); sabor, generación, esquema y escala son datos distintos. Los
neutrinos pertenecen a la completación fermiónica; la condición Majorana
requiere además autoconjugación y un término de masa compatible. Fotón y gluón
pertenecen a la carta nula exacta.

\begingroup
\let\chapter\section
\let\section\subsection
\let\subsection\subsubsection
\let\subsubsection\paragraph
\let\clearpage\relax
% Vista científica exacta materializada desde una rama condicional.
% Fuente: source/demostraciones_primarias/09_06h_atlas_familias_rev6.tex; rama: HMTIncludeAtlasPhysical.
% No modifica el contenido matemático de la rama de origen.
\chapter[Correspondencia con sectores físicos]{Correspondencia entre clases internas y sectores físicos}
\label{sec:censo-fisico-atlas-rev7}

La realización física conserva los cocientes anteriores y añade los datos
que una representación necesita sin convertirlos en nuevas celdas. Para una
sección persistente \(X\), escribimos de forma tipada
\[
 \mathcal P_X=
 \bigl[q(X),\,[X]_{\mathcal R},\,\operatorname{Rep}(X),\,
       \gamma_X,\,\tau_{12}(X),\,\sigma_X\bigr].
\]
Aquí \(q(X)\) fija la multisección, \([X]_{\mathcal R}\) la familia de ruta,
\(\operatorname{Rep}(X)\) la representación, \(\gamma_X\) la historia
persistente, \(\tau_{12}(X)\) su palabra dodecafásica y \(\sigma_X\) la firma
entera extraída después del registro. Los niveles \(G0,\ldots,G4\) siguen
siendo profundidades del atlas; no se renombran como las tres generaciones
físicas.

La información de sabor se conserva mediante tres coordenadas distintas:
\[
 \operatorname{flav}(X)=
 \bigl(\sigma_{ud}(X),g(X),\chi_{\rm fina}(X)\bigr).
\]
La primera separa las ramas leptónica, neutrino, tipo arriba y tipo abajo;
la segunda distingue la generación física; la tercera conserva la
orientación y la memoria que separan secciones dentro de una misma rama. En
los quarks, la órbita residual \(R/G/B\) se añade después de fijar el sabor y
antes de evaluar el carácter. Color, sabor y masa son, por tanto, lecturas de
tipos diferentes.

% El inventario humano sectorial se conserva en su totalidad,
% pero se monta después del cierre electrónico, en la residencia sectorial.

Este censo no agota las realizaciones de la ley. Mesones y bariones se
construyen componiendo registros enriquecidos antes de proyectar la firma;
resonancias y partículas virtuales se distinguen por su horizonte de
persistencia. Los sectores Fibonacci, \(Z_6\) y Majorana--Ising conservan
codominios topológicos propios y no se fuerzan dentro de una tabla ordinaria
de masas.

\section{Firmas de trayectoria y pares dodecafásicos coindexados por sabor}

Para el conjunto de sabores
\(\mathcal Q=\{u,d,c,s,t,b\}\), la tabla define una correspondencia finita
entre una huella geométrica y un par dodecafásico. En cada fila se cumple
\[
 W_\pm=6n_A\pm s_C.
\]
Esta tabla no es una comparación a posteriori, sino la gráfica de un mapa
finito sobre el dominio nominal certificado. Si
\(\mathscr R_{\mathcal Q}^{\rm nom}=\{\rho_u,\rho_d,\rho_c,\rho_s,\rho_t,\rho_b\}\)
es el conjunto de huellas CT--108 que siguen, definimos
\[
 H_{\rm flav}:\mathscr R_{\mathcal Q}^{\rm nom}\longrightarrow\mathbb Z^2,
 \qquad H_{\rm flav}(\rho_q)=(n_A(q),s_C(q)).
\]
Las seis imágenes son únicas, la tabla exhausta el dominio y ninguna columna
contiene una masa. Por ello \(H_{\rm flav}\) es un extractor causal cerrado en
el dominio nominal: primero lee huella, orientación y memoria, y sólo después
el carácter evalúa la masa. La órbita ternaria de color se cocienta sin cambiar
\(H_{\rm flav}\); la reversión conserva la clase de masa y transporta la
representación partícula--antipartícula. La tabla es, por tanto, la definición
ejecutable del morfismo y no una obligación futura:

\begin{longtable}{c c c c c c c}
\toprule
Sabor & Generación & Giros/cambios & Eje & \(N,E,S,O\)
& \((n_A,s_C)\) & \((W_+,W_-)\)\\
\midrule
\endfirsthead
\toprule
Sabor & Generación & Giros/cambios & Eje & \(N,E,S,O\)
& \((n_A,s_C)\) & \((W_+,W_-)\)\\
\midrule
\endhead
\(u\)&1&21/30&N--S&34,31,31,12&(11,7)&(73,59)\\
\(d\)&1&20/32&E--O&27,35,18,28&(17,8)&(110,94)\\
\(c\)&2&21/29&E--O&20,38,24,26&(61,8)&(374,358)\\
\(s\)&2&17/27&N--S&50,32,12,14&(41,-3)&(243,249)\\
\(t\)&3&21/26&N--S&46,27,16,19&(100,-1)&(599,601)\\
\(b\)&3&21/30&E--O&12,63,9,24&(71,-5)&(421,431)\\
\bottomrule
\end{longtable}

La coincidencia de dos giros o de una altura espectral no identifica
especies: la ruta cardinal, la orientación, el registro dodecafásico y la
representación permanecen en el objeto. La cadena
celda--ruta--registro cronológico enriquecido--firma--carácter se define con
independencia de \(H_{\rm flav}\). Las
cuatro orientaciones iniciales de cada celda producen
\(81\cdot4=324\) genealogías ejecutables; ese cardinal cuenta historias, no
partículas, familias ni masas.

\endgroup
\begingroup
\let\chapter\section
\let\section\subsection
\let\subsection\subsubsection
\let\subsubsection\paragraph
\let\clearpage\relax
\section{Estratificación generacional, soportes quark y realización neutrínica}
\label{sec:generaciones-quarks-neutrinos-rev2}
\index{generaciones fermiónicas!estratificación HMT--MD}
\index{quark!soportes celulares y registro}
\index{PMNS!transporte de bases}

La ley general actúa sobre multisecciones y no sobre una lista de masas. La
aparición de tres generaciones se expresa, por ello, como compatibilidad de
tres estructuras: la estratificación del atlas, los registros de sabor y el
espacio tridimensional de mezcla. En el sector leptónico cargado, las
multisecciones
\begin{equation}
 G_0^{\ell},\qquad G_2^{\ell},\qquad G_4^{\ell}
 \label{eq:tres-multisecciones-leptonicas-rev2}
\end{equation}
tienen rangos \(1,8,16\). El centro \(G_0^\ell\) contiene la sección
electrónica de rango uno; \(G_2^\ell\) y \(G_4^\ell\) conservan, en cambio,
espectros operatoriales no triviales. La organización generacional completa
es
\begin{center}
\small
\begin{tabularx}{0.96\textwidth}{@{}c c c c X@{}}
\toprule
\textbf{Generación} & \textbf{Leptón} & \textbf{Par quark} &
\textbf{Sabor neutro} & \textbf{Soporte estructural}\\
\midrule
1 & \(e\) & \((u,d)\) & \(\nu_e\) & centro \(G_0^\ell\), registro de
primera generación y primer vector de la base de interacción\\
2 & \(\mu\) & \((c,s)\) & \(\nu_\mu\) & multisección \(G_2^\ell\),
registro de segunda generación y segundo vector de la base de interacción\\
3 & \(\tau\) & \((t,b)\) & \(\nu_\tau\) & multisección \(G_4^\ell\),
registro de tercera generación y tercer vector de la base de interacción\\
\bottomrule
\end{tabularx}
\end{center}
La conjugación partícula--antipartícula invierte las coordenadas de carga y
orientación de la representación, pero conserva esta estratificación; por
consiguiente, no duplica el número de generaciones.

La dimensión tres dispone además de un portador de incidencia propio. La
representación orientada dodecafásica contiene
\begin{equation}
 V_{12}\cong V_1\oplus V_3\oplus V_{3'}\oplus V_5,
 \qquad \operatorname{rank}P_3=3,
 \label{eq:descomposicion-rango-tres-generaciones-rev2}
\end{equation}
y el entrelazador dodecafásico--Witt transporta el sector positivo según
\begin{equation}
 P_G\xleftrightarrow{\;W_{GK}\;}P_3
 \xleftrightarrow{\;U_W\;}Q_3,
 \qquad
 \operatorname{rank}P_G=\operatorname{rank}P_3
 =\operatorname{rank}Q_3=3.
 \label{eq:transporte-rango-tres-generaciones-rev2}
\end{equation}
Este portador tridimensional permite comparar bases generacionales. Las
rutas, las firmas y los operadores de masa proporcionan después el contenido
espectral que CKM y PMNS transportan.

\subsection{Posición interna y registros de los \(6\) sabores de quark}

La posición de un sabor quark es el dato compuesto
\begin{equation}
 \mathfrak q_X=
 (\text{rama de carga},\text{generación},(\mathbb Z/3\mathbb Z)_{\rm color},
  \text{multisección},\text{registro de ruta}).
 \label{eq:posicion-quark-compuesta-rev2}
\end{equation}
Las ramas de carga \(-1/3\) poseen los soportes
\begin{align}
 D_1&=\{(4,5),(6,5)\},\nonumber\\
 D_2&=\{(4,3),(4,7),(6,3),(6,7)\},\nonumber\\
 D_3&=\{(2,3),(2,5),(2,7),(8,3),(8,5),(8,7)\},\nonumber\\
 D_4&=\{(2,1),(2,9),(4,1),(4,9),(6,1),(6,9),(8,1),(8,9)\},
 \label{eq:soportes-quark-down-rev2}
\end{align}
mientras las ramas de carga \(+2/3\) ocupan
\begin{align}
 U_1&=\{(4,4),(4,6),(6,4),(6,6)\},\nonumber\\
 U_3&=\{(2,2),(2,4),(2,6),(2,8),(4,2),(4,8),
       (6,2),(6,8),(8,2),(8,4),(8,6),(8,8)\}.
 \label{eq:soportes-quark-up-rev2}
\end{align}
Cada soporte se eleva a cuatro orientaciones por celda. El cociente de color
organiza \(36=12_{\rm R}+12_{\rm G}+12_{\rm B}\) biografías orientadas y conserva una misma
masa por sabor antes de la realización de color.

El transductor nominal
\begin{equation}
 \mathcal N_{\rm ruta}(\Gamma_f)
 =((N,E,S,O),h_{369},r_5)
 \in\mathbb Z_{\ge0}^{4}\times\mathbb Z_{\ge0}^{2}
 \label{eq:transductor-nominal-quark-rev2}
\end{equation}
publica las coordenadas siguientes antes de toda realización escalar:
\begin{center}
\small
\begin{tabularx}{0.96\textwidth}{@{}c c c c c X@{}}
\toprule
\textbf{Sabor} & \textbf{Generación} & \textbf{Rama} &
\textbf{\((N,E,S,O)\)} & \textbf{\((h_{369},r_5)\)} &
\textbf{Eje dominante}\\
\midrule
\(u\)&1&\(+2/3\)&\((34,31,31,12)\)&\((51,13)\)&N--S\\
\(d\)&1&\(-1/3\)&\((27,35,18,28)\)&\((47,17)\)&E--O\\
\(c\)&2&\(+2/3\)&\((20,38,24,26)\)&\((49,16)\)&E--O\\
\(s\)&2&\(-1/3\)&\((50,32,12,14)\)&\((57,11)\)&N--S\\
\(t\)&3&\(+2/3\)&\((46,27,16,19)\)&\((57,9)\)&N--S\\
\(b\)&3&\(-1/3\)&\((12,63,9,24)\)&\((47,7)\)&E--O\\
\bottomrule
\end{tabularx}
\end{center}
Las firmas nominales asociadas son
\begin{center}
\small
\begin{tabular}{@{}c c c@{}}
\toprule
\textbf{Sabor} & \textbf{\((n_A,s_C)\)} & \textbf{\((W_+,W_-)\)}\\
\midrule
\(u\)&\((11,7)\)&\((73,59)\)\\
\(d\)&\((17,8)\)&\((110,94)\)\\
\(c\)&\((61,8)\)&\((374,358)\)\\
\(s\)&\((41,-3)\)&\((243,249)\)\\
\(t\)&\((100,-1)\)&\((599,601)\)\\
\(b\)&\((71,-5)\)&\((421,431)\)\\
\bottomrule
\end{tabular}
\end{center}
Estas coordenadas distinguen sabores que comparten rama o generación y cierran
el morfismo de sabor sin consultar ninguna masa. Para una sección física (P),
definimos
\begin{equation}
 \operatorname{sabor}(P)
 =\bigl(\sigma_{\rm ud}(P),g(P),\chi_{\rm fina}(P)\bigr),
 \qquad
 \chi_{\rm fina}(P)
 =(\text{giros},\text{switches},\text{eje},N,E,S,O,h_{369},r_5),
 \label{eq:morfismo-sabor-ruta-cerrado-rev2}
\end{equation}
y, para quarks, se conserva además la órbita ternaria
\([\gamma_P]_{\rm col}=\{\gamma_P^R,\gamma_P^G,\gamma_P^B\}\). La realización
física es el mapa ya determinado por los datos precedentes,
\begin{equation}
 \mathcal S_{\rm phys}(P,\mathfrak M):
 \bigl(\sigma_{\rm ud},g,\chi_{\rm fina},[\gamma_P]_{\rm col}\bigr)
 \longmapsto \gamma^{\rm surv}_{P,\mathfrak M}
 \longmapsto
 \bigl(\sigma(\gamma^{\rm surv}_{P,\mathfrak M}),
       B_F^D,B_F^\Omega\bigr),
 \label{eq:realizacion-fisica-sabor-ruta-cerrada-rev2}
\end{equation}
donde \(\mathfrak M\) es el acoplamiento compatible y
\(\gamma^{\rm surv}_{P,\mathfrak M}\) la ruta que persiste bajo él. Para
leptones cargados la órbita de color se reduce al objeto terminal. El carácter
de masa se evalúa entonces sobre esa misma genealogía:
\begin{equation}
 \boxed{
 \frac{m_{\mathfrak M}(P)}{m_e^{\rm HMT}}
 =\chi_{\rm mass}\!\left(
   \sigma\!\left(\gamma^{\rm surv}_{P,\mathfrak M}\right)\right).}
 \label{eq:masa-como-caracter-persistencia-rev2}
\end{equation}
Por tanto, sabor, generación y color no esperan un selector futuro: publican la
clase de ruta y su autovector antes de abrir la comparación externa. La masa no
elige retrospectivamente la ruta; sólo lee multiplicativamente la persistencia
que el acoplamiento ha estabilizado.

\subsection{Fermionicidad, autovalores de masa y transporte PMNS}

La clausura exterior del funtor de intercambio construye las relaciones CAR
y fija la fermionicidad. En el sector neutro, la etiqueta
\(\nu_e,\nu_\mu,\nu_\tau\) designa una base de interacción; los valores de
masa son autovalores de un operador neutro mezclado. Los marcos internos
\(F_\ell,F_\nu:\mathbb C^3\to\operatorname{im}P_3\) comparan esas bases, y
la matriz
\begin{equation}
 U_{\rm HMT}=F_\ell^*\Phi_LF_\nu
 \label{eq:transporte-pmns-generaciones-rev2}
\end{equation}
es un transporte unitario. PMNS cambia de base entre estados de interacción
y autovectores de masa; la escala neutra pertenece al espectro del operador,
no al cambio de base.

\subsection{Extractor neutro Z0 anterior a la proyección}
\label{sec:extractor-neutro-z0-torsion-rev2}

La construcción recuperada del sector neutro actúa sobre el registro
dodecafásico enriquecido, antes de colapsar hoja, torsión y bit
bidireccional. Sean \(E_m\), \(m\in\mathbb Z/12\mathbb Z\), los doce
eventos de nueve ticks del registro CT--108. En cada tick \(t\) se conservan
\[
 \theta(t)=\mathbf1_{\{d(t)=9\}},\qquad
 \varepsilon(t)\in\{-1,+1\},\qquad
 \kappa(t)=\mathbf1_{\{9\ {\rm en\ corona}\}},
 \qquad \omega(t)=(-1)^{\kappa(t)}.
\]
El primer factor detecta la torsión nonádica, el segundo conserva la hoja
two--ways y el tercero cambia el signo en el giro de \(180^\circ\) de la
corona. El entero torsional de cada evento y su diferencia antipodal son
\begin{equation}
 T_\nu(m)=\sum_{t\in E_m}
   \omega(t)\varepsilon(t)\theta(t),\qquad
 e_m^\nu=T_\nu(m)-T_\nu(m+6),
 \label{eq:extractor-z0-torsion-antipodal-rev2}
\end{equation}
y el trit neutro queda fijado, sin parámetros continuos, por
\begin{equation}
 \tau_m^\nu=\operatorname{sgn}(e_m^\nu)
 \in\{-1,0,+1\}.
 \label{eq:trit-neutro-torsion-antipodal-rev2}
\end{equation}

Para las cuatro órbitas de salto cuatro
\[
 I_a=\{a,a+4,a+8\}\pmod{12},\qquad a=0,1,2,3,
\]
el coeficiente de evento que entra en la firma y el observable de mayoría
de signos son, respectivamente,
\begin{equation}
 b_{120}^\nu
 =\sum_{a=0}^{3}\operatorname{sgn}
   \left(\sum_{m\in I_a}e_m^\nu\right),
 \qquad
 \nu_{120}^{\nu,\mathrm{sgn}}
 =\sum_{a=0}^{3}\operatorname{sgn}
   \left(\sum_{m\in I_a}\tau_m^\nu\right).
 \label{eq:canales-120-neutros-distintos-rev2}
\end{equation}
No se identifican estos dos enteros: el primero es la coordenada
\(b_{120}=\nu_{120}^{\rm ev}\) del carácter, mientras el segundo es un
diagnóstico de mayoría. La memoria cuadrática es
\begin{equation}
 b_\zeta^\nu=\nu_{270}^\nu
 =\sum_{m=0}^{11}\tau_m^\nu\tau_{m+3}^\nu.
 \label{eq:canal-270-neutro-rev2}
\end{equation}
Por tanto, para cada historia neutra superviviente
\(\gamma_{\nu_i}\), el registro publica antes de toda comparación la firma
\[
 v_\nu(\gamma_{\nu_i})
 =\bigl(0,0,k_{\Delta_4}^{(i)},b_{120}^{\nu_i},
        b_\zeta^{\nu_i}\bigr)
\]
y el exponente relativo
\begin{equation}
 \mathcal A_\nu^{(i)}
 =k_{\Delta_4}^{(i)}\Delta_4
  +b_{120}^{\nu_i}s_{120}
  +b_\zeta^{\nu_i}\zeta_{270},
 \qquad
 \zeta_{270}=135\Delta_4^2.
 \label{eq:caracter-relativo-neutro-z0-rev2}
\end{equation}
Aquí \(\zeta_{270}\) no se sustituye por el momento de torre
\(s_{270}=q_-^{270}-q_+^{270}\).

La conjugación global two--ways
\(\tau^\nu\mapsto-\tau^\nu\) invierte
\(b_{120}^\nu\) y conserva \(b_\zeta^\nu\). En consecuencia genera dos
ramas de espectro relativo,
\begin{equation}
 \mathcal A_{\nu,\pm}^{(i)}
 =k_{\Delta_4}^{(i)}\Delta_4
  \pm b_{120}^{\nu_i}s_{120}
  +b_\zeta^{\nu_i}\zeta_{270},
 \qquad
 \text{NO}\longleftrightarrow\text{IO},
 \label{eq:bifurcacion-neutra-no-io-rev2}
\end{equation}
que se distinguen por el signo posterior de
\(\Delta m_{31}^2\), no por una selección con datos PDG. Esta construcción
cierra el extractor discreto Z0 y la bifurcación relativa. La escala absoluta
común \(m_0\) y el núcleo físico PMNS de rango completo permanecen objetos
separados: no se los introduce dentro del selector ni se los declara
deducidos por la sola paridad two--ways.

La realización con una única fase de rango uno produce
\[
 (\theta_{12},\theta_{23},\theta_{13})
 =(19.6807^\circ,56.4970^\circ,29.6224^\circ)
\]
y no reproduce la fenomenología de mezcla. Este control selecciona como dominio adecuado un
núcleo de registro de rango completo,
\begin{equation}
 G_{ab}=\frac1{\sqrt{|L_a||N_b|}}
 \sum_{c\in L_a}\sum_{d\in N_b}
 K_{\rm registro}(c,d;\xi_c,\xi_d),
 \qquad U_{\rm int}=\operatorname{polar}(G),
 \label{eq:kernel-registro-rango-completo-rev2}
\end{equation}
cuya fuente técnica se desarrolla en el capítulo de transporte CKM y
PMNS. Así quedan separados el operador de masa, sus autovalores y el mapa de
lectura por sabor.

\endgroup
\begingroup
\let\chapter\section
\let\section\subsection
\let\subsection\subsubsection
\let\subsubsection\paragraph
\let\clearpage\relax
\section{Transporte de los operadores de masa mediante CKM}
\label{sec:transporte-ckm-operadores-rev11}

La incidencia N42 fija el sexteto \((4,5,7,6,8,3)\) y la matriz racional
\[
 M_{\rm CKM}=
 \begin{pmatrix}
 2&-5&28\\
 0&7&-25/2\\
 1&-20&-2/3
 \end{pmatrix},
 \qquad \det M_{\rm CKM}=-\frac{3857}{6}.
\]
Las coordenadas angulares se construyen por
\[
 \begin{pmatrix}\theta_{12}\\\theta_{23}\\\theta_{13}\end{pmatrix}
 =M_{\rm CKM}
 \begin{pmatrix}A\\ C^*/6\\ (180/\pi)\Delta_4\end{pmatrix},
 \qquad \delta_{\rm CP}=9A.
\]
Esta transformación pertenece al cierre matemático de HMT. En su realización
física, sean \(B_u\) y \(B_d\) los operadores bidireccionales ya construidos en las
fibras de quarks de tipo up y down. El cambio de base actúa según
\begin{equation}
 B_d^{(u)}=V_{\rm CKM}B_dV_{\rm CKM}^*,
 \qquad
 R_{\rm act}^{B_d^{(u)}}
 =V_{\rm CKM}R_{\rm act}^{B_d}V_{\rm CKM}^*.
 \label{eq:transporte-ckm-kappa-rev11}
\end{equation}
CKM transporta autovalores entre bases y conserva espectro y determinante.
El control de violación CP adopta la identidad
\[
 \left|\det[B_u,V_{\rm CKM}B_dV_{\rm CKM}^*]\right|
 =2|J|\prod_{i<j}|\kappa_i^u-\kappa_j^u|
       \prod_{a<b}|\kappa_a^d-\kappa_b^d|.
\]
La matriz de mezcla y el lector de ruta desempeñan, por tanto, funciones
complementarias: el lector determina el operador espectral y CKM determina
su representación en la base de sabor.

\section{Marcos geométricos de rango \(3\) para el sector leptónico}
\label{sec:frames-pmns-rev11}

La construcción leptónica utiliza la involución diagonal de signo
\[
 S_{\rm sign}=\operatorname{diag}(1,1,1,-1,-1,-1,1,1,1,-1,-1,-1),
 \qquad P_-=(I-S_{\rm sign})/2,
\]
y el proyector icosaédrico canónico \(P_3\) de rango tres. La compresión
\[
 C_-=P_-P_3P_-\big|_{\operatorname{im}P_-}
\]
posee espectro
\[
 \left\{\frac{5+\sqrt5}{10},\frac12,
 \frac{5-\sqrt5}{10},0,0,0\right\}.
\]
Los tres autovalores positivos son simples. Sus proyectores espectrales,
aplicados al indicador de la hexada negativa orientada y normalizados después
de proyectar mediante \(P_3\), construyen un marco ortonormal cargado
\(F_\ell\) que satisface
\[
 F_\ell^*F_\ell=I_3,
 \qquad F_\ell F_\ell^*=P_3.
\]

Para el sector neutrino, sean \(H_1,H_2,H_3\) las tres hexadas positivas
orientadas y
\[
 X=(P_3\mathbf1_{H_1},P_3\mathbf1_{H_2},P_3\mathbf1_{H_3}).
\]
Su matriz de Gram es positiva definida y cumple
\[
 \det(X^*X)=\frac{5-\sqrt5}{40}>0.
\]
Por consiguiente,
\[
 F_\nu=X(X^*X)^{-1/2},
 \qquad F_\nu^*F_\nu=I_3,
 \qquad F_\nu F_\nu^*=P_3.
\]
El transporte real
\begin{equation}
 O=F_\ell^*F_\nu\in SO(3)
 \label{eq:transporte-frames-pmns-rev11}
\end{equation}
se obtiene sin masas ni ángulos PMNS.

\HMTClaim{MD-R-PMNS-RANKONE-UNITARY}
\begin{theorem}[Unitariedad de la fase dodecafásica mínima]
\label{thm:unitaria-interna-pmns-rev11}
Sea \(b_K\in\operatorname{im}P_3\) el vector normalizado del testigo de Witt,
\(Q_K=|b_K\rangle\langle b_K|\), y
\[
 \Phi_L=I+(\zeta_{12}-1)Q_K,
 \qquad \zeta_{12}=e^{i\pi/6}.
\]
Entonces
\begin{equation}
 U_{\rm HMT}=F_\ell^*\Phi_LF_\nu
 \label{eq:unitaria-interna-pmns-rev11}
\end{equation}
es una matriz unitaria independiente de valores metrológicos objetivo. La
inversión de orientación conjuga
\(U_{\rm HMT}\).
\end{theorem}

\begin{proof}
\(Q_K\) es un proyector ortogonal, luego \(\Phi_L\) es unitaria. Los dos
marcos son isometrías con la misma imagen \(\operatorname{im}P_3\). Por
tanto
\[
 U_{\rm HMT}^*U_{\rm HMT}
 =F_\nu^*\Phi_L^*P_3\Phi_LF_\nu=I_3.
\]
El cambio \(\pi/6\mapsto-\pi/6\) produce \(\Phi_L^*\), y todas las
primitivas restantes son reales.
\end{proof}

\subsection{Control negativo de rango uno y transporte PMNS de rango completo}
\label{sec:no-go-pmns-rango-uno-rev11}
\HMTClaim{MD-R-PMNS-RANKONE-PHYSICAL-NOGO}

Congelada la construcción, el contraste en el orden declarado produce
\[
 (\theta_{12},\theta_{23},\theta_{13})
 =(19.6807^\circ,56.4970^\circ,29.6224^\circ),
 \qquad J=-0.0426621.
\]
La mejor de las 36 permutaciones discretas deja
\(\theta_{13}\simeq17.02^\circ\). La fase de rango uno genera una matriz
unitaria interna, pero conserva un estabilizador \(U(2)\) y queda descartada
como realización PMNS física.

La exclusión del portador de rango uno no deja indeterminada la realización
PMNS. Sean \(M_\ell\) y \(M_\nu\) los operadores cargado y neutro producidos
por la Ley General de Masas, y sean \(F_\ell,F_\nu\) sus marcos espectrales
ortonormales completos sobre la imagen común \(\operatorname{im}P_3\). El
transporte leptónico de rango completo queda entonces construido por
\[
 U_{\rm PMNS}:=F_\ell^*F_\nu,
 \qquad
 U_{\rm PMNS}^*U_{\rm PMNS}=I_3.
\]
Para espectros simples, esta salida está determinada salvo reajustes
diagonales de fase y el orden espectral declarado; ante degeneraciones, la
salida exacta es la clase doble de marcos correspondiente. El núcleo de
registro constituye un certificado finito adicional de esa realización, no
su generador ni una condición pendiente para su existencia. Para fibras
cargadas \(L_a\) y neutrínicas \(N_b\), su forma agregada es
\[
 G_{ab}=\frac{1}{\sqrt{|L_a||N_b|}}
 \sum_{c\in L_a}\sum_{d\in N_b}K_{\rm reg}(c,d),
 \qquad U_{\rm int}=\operatorname{polar}(G),
\]
con naturalidad bajo orientación. Cuando las entradas explícitas del registro
verifican \(\det G\ne0\), su polar proporciona un certificado alternativo del
transporte ya construido. El núcleo lee acarreo, hoja, orientación, vacancia
y holonomía; sus parámetros se congelan antes de abrir los ángulos externos.

\paragraph{Dominio y codominio de las realizaciones CKM y PMNS}
La matriz racional CKM, los marcos \(F_\ell,F_\nu\), el transporte \(O\), la
fase \(\Phi_L\), la unitariedad de \(U_{\rm HMT}\) y el transporte completo
\(U_{\rm PMNS}=F_\ell^*F_\nu\) son resultados internos exactos bajo las
primitivas declaradas. CKM constituye un transporte físico una vez fijadas
las bases de sabor y masa. El contraste falsifica únicamente la fase PMNS de
rango uno; no rebaja la construcción espectral completa. La naturalidad y el
criterio \(\det G\ne0\) pertenecen al certificado alternativo de registro y
permanecen separados de la existencia y la unitariedad del transporte PMNS.

\endgroup
\begingroup
\let\chapter\section
\let\section\subsection
\let\subsection\subsubsection
\let\subsubsection\paragraph
\let\clearpage\relax
\section{Involución de hojas en \(4\) representaciones}

La involuci\'on \(C^*\mapsto-C^*\) organiza cuatro respuestas distintas:
\begin{equation}
\label{rm:eq:modular-sheet-parities}
\begin{aligned}
S_{\rm bi}(A,-C^*)&=S_{\rm bi}(A,C^*),\\
\Gamma_{6\to8}(A,-C^*)&=-\Gamma_{6\to8}(A,C^*),\\
(\widehat\mu,\widehat\varepsilon)&\longleftrightarrow
(\widehat\varepsilon,\widehat\mu),\\
\widehat c&\longmapsto\widehat c,
\qquad D_A\longmapsto D_A.
\end{aligned}
\end{equation}
La informaci\'on bicapa es par; la orientaci\'on areal es impar; las dos
constitutivas forman un doblete, y los determinantes quedan fijos. Esta tabla
de paridades precede a cualquier interpretaci\'on como simetr\'ia f\'isica.

\section{Transformaci\'on racional involutiva en el espacio de par\'ametros CKM}

Sea
\begin{equation}
\label{rm:eq:modular-Mckm}
M_{\rm CKM}=
\begin{pmatrix}
2&-5&28\\
0&7&-25/2\\
1&-20&-2/3
\end{pmatrix},
\qquad
\det M_{\rm CKM}=-\frac{3857}{6}.
\end{equation}
La inversa existe sobre \(\mathbb Q\) y es
\begin{equation}
\label{rm:eq:modular-Mckm-inverse}
M_{\rm CKM}^{-1}=
\begin{pmatrix}
1528/3857&3380/3857&801/3857\\
75/3857&176/3857&-150/3857\\
6/551&-30/551&-12/551
\end{pmatrix}.
\end{equation}
Con
\begin{equation}
\label{rm:eq:modular-ckm-coordinates}
h=\frac{C^*}{6},
\qquad
\Delta=\frac{180}{\pi}\Delta_4,
\qquad
\boldsymbol\theta=M_{\rm CKM}(A,h,\Delta)^T,
\end{equation}
la coordenada impar entra en la direcci\'on
\begin{equation}
\label{rm:eq:modular-ch}
c_h=\begin{pmatrix}-5\\7\\-20\end{pmatrix}.
\end{equation}

Si \(c_A=(2,0,1)^T\) y
\(c_\Delta=(28,-25/2,-2/3)^T\), las tres columnas
\((c_A,c_h,c_\Delta)\) son precisamente la base transportada por
\(M_{\rm CKM}\). La segunda fila de
\cref{rm:eq:modular-Mckm-inverse} es el covector \(\ell_h\) que extrae la
coordenada impar:
\begin{equation}
\label{rm:eq:modular-dual-coordinate-check}
\ell_hc_A=0,
\qquad
\ell_hc_h=1,
\qquad
\ell_hc_\Delta=0.
\end{equation}

\begin{theorem}[Reflexi\'on CKM inducida por la hoja]
\label{rm:thm:modular-ckm-reflection}
Sea
\begin{equation}
\label{rm:eq:modular-Dh-lh}
D_h=\operatorname{diag}(1,-1,1),
\qquad
\ell_h=\frac1{3857}\begin{pmatrix}75&176&-150\end{pmatrix}.
\end{equation}
La involuci\'on inducida en el espacio de \(\boldsymbol\theta\) es
\begin{equation}
\label{rm:eq:modular-Rckm}
\boxed{
\mathcal R_{\rm CKM}
=M_{\rm CKM}D_hM_{\rm CKM}^{-1}
=I-2c_h\ell_h.}
\end{equation}
Satisface
\begin{equation}
\label{rm:eq:modular-Rckm-properties}
\mathcal R_{\rm CKM}^2=I,
\qquad
\det\mathcal R_{\rm CKM}=-1,
\qquad
\mathcal R_{\rm CKM}M_{\rm CKM}=M_{\rm CKM}D_h.
\end{equation}
\end{theorem}

\begin{proof}
La multiplicaci\'on de
\cref{rm:eq:modular-Mckm,rm:eq:modular-Mckm-inverse} da la identidad; en
particular, \cref{rm:eq:modular-dual-coordinate-check} prueba que
\(c_h\ell_h\) es el proyector de rango uno sobre \(\mathbb Qc_h\) a lo
largo de \(\operatorname{span}\{c_A,c_\Delta\}\). El producto exacto
\(\ell_hc_h=1\) implica
\[
(I-2c_h\ell_h)^2
=I-4c_h\ell_h+4c_h(\ell_hc_h)\ell_h=I.
\]
El operador fija punto por punto el hiperplano
\(\ker\ell_h=\operatorname{span}\{c_A,c_\Delta\}\) y cambia el signo de
\(c_h\). Sus autovalores son, por tanto, \(1,1,-1\); de aqu\'i
\(\det\mathcal R_{\rm CKM}=-1\) y
\(\Tr\mathcal R_{\rm CKM}=1\). Finalmente,
\[
\mathcal R_{\rm CKM}M_{\rm CKM}
=M_{\rm CKM}D_hM_{\rm CKM}^{-1}M_{\rm CKM}
=M_{\rm CKM}D_h,
\]
que prueba la identidad de entrelazamiento.
\end{proof}

La matriz expl\'icita es
\begin{equation}
\label{rm:eq:modular-Rckm-matrix}
\mathcal R_{\rm CKM}=
\begin{pmatrix}
4607/3857&1760/3857&-1500/3857\\
-150/551&199/551&300/551\\
3000/3857&7040/3857&-2143/3857
\end{pmatrix}.
\end{equation}
Aunque \(\mathcal R_{\rm CKM}\) no es una reflexi\'on ortogonal para el
producto eucl\'ideo de las tres coordenadas angulares, s\'i lo es para la
m\'etrica racional transportada
\begin{equation}
\label{rm:eq:modular-ckm-metric}
G_{\rm CKM}=(M_{\rm CKM}^{-1})^TM_{\rm CKM}^{-1},
\qquad
\mathcal R_{\rm CKM}^TG_{\rm CKM}\mathcal R_{\rm CKM}=G_{\rm CKM}.
\end{equation}
En efecto, en la base \((c_A,c_h,c_\Delta)\), la m\'etrica es la identidad
y la involuci\'on es \(D_h\). Por eso el t\'ermino ``reflexi\'on'' tiene un
contenido m\'etrico preciso y no describe s\'olo que
\(\mathcal R_{\rm CKM}^2=I\).
La fase \(\delta_{\rm CP}=9A\) permanece fija bajo esta reflexi\'on. Por
ello, \(\mathcal R_{\rm CKM}\) no es la conjugaci\'on CP f\'isica, sino el
transporte de la inversi\'on de hoja HMT al espacio de mezcla.

La inversa racional de \(M_{\rm CKM}\) recupera
\begin{align}
A&=\frac{1528\theta_{12}+3380\theta_{23}+801\theta_{13}}{3857},
\label{rm:eq:modular-A-from-ckm}\\
C^*&=\frac{6(75\theta_{12}+176\theta_{23}-150\theta_{13})}{3857}.
\label{rm:eq:modular-C-from-ckm}
\end{align}
La tercera coordenada tambi\'en se recupera sin p\'erdida:
\begin{equation}
\label{rm:eq:modular-Delta-from-ckm}
\Delta
=\frac{6\theta_{12}-30\theta_{23}-12\theta_{13}}{551}.
\end{equation}
Estas ecuaciones son cambios de coordenadas. No convierten CKM en causa
retrospectiva de \(A,C^*\) o de \(\gammaH\).

Como reconocimiento terminal, el constructor previo determina
\begin{equation}
\label{rm:eq:modular-ckm-values}
(\theta_{12},\theta_{23},\theta_{13},\delta_{\rm CP})
=(13.003313589509^\circ,
2.398056157332^\circ,
0.213977382244^\circ,
65.676173123554^\circ).
\end{equation}
La reflexi\'on racional se demuestra antes de usar estos decimales.

\begin{keyresult}
La transici\'on \(6\to8\) se formaliza primero mediante la inclusi\'on
\(\iota_{6,8}:\mathcal F_6\hookrightarrow\mathcal F_8\), de la que se
obtienen
\(90,120\) y el defecto \(-1/12\). La palabra de seis trits proporciona
despu\'es la biyecci\'on \(729\leftrightarrow729\), y
\(\mathcal J_{6\to8}\) transporta los dos modos colectivos al borde. Sobre
ese antecedente, el \`area y el Hamiltoniano modular reconstruyen exactamente
las dos familias; la purificaci\'on termocampo doble determina la entrop\'ia
de entrelazamiento, y la primera ley
finita conserva el defecto de entrop\'ia relativa. En CKM, la inversi\'on de
hoja se convierte en una reflexi\'on racional de rango impar uno.
\end{keyresult}

\begin{falsifier}
La genealog\'ia inicial falla si \(\iota_{6,8}\) no es inyectiva, si los
conteos no son \(90/120\), si \(\delta_{6,8}^{\rm inc}\neq-1/12\), si
\(\beta_{729}\) no es biyectiva o si se la presenta falsamente como
isomorfismo aditivo. El transporte colectivo falla si
\(\mathcal J_{6\to8}D_{\rm inc}\neq
D_{\rm area}\mathcal J_{6\to8}\). La reconstrucci\'on falla si el
determinante de los observables de borde no es treinta o si
\cref{rm:eq:modular-inverse-channels} no recupera las ocupaciones.
La TFD falla si su traza parcial no es \(\rho_A\). La primera ley no se
extiende a cambios finitos sin la entrop\'ia relativa de
\cref{rm:eq:modular-finite-first-law}. El tipo
\(III_{\lambda_\partial}\) queda invalidado si un canal no persiste, si no hay
producto compatible o si cambia el grupo de razones. La reflexi\'on CKM
queda refutada por cualquiera de
\(\mathcal R^2\neq I\), \(\det\mathcal R\neq-1\) o
\(\mathcal RM\neq MD_h\). Una f\'ormula
\(\boldsymbol\theta=f(\gammaH)\) sin entrelazador viola la independencia
tipada de las dos realizaciones funcionales del antecedente com\'un.
\end{falsifier}

\begin{statusbox}
Inclusi\'on \(H\subset O\), conteos \(90/120\), dos normalizaciones del
defecto \(-1/12\), biyecci\'on de seis trits \(729\leftrightarrow729\),
entrelazador colectivo, escala \(a_0\), operadores finitos, espectro areal,
TFD, defecto orientado, primera ley local y finita y reflexi\'on CKM:
\status{Teorema interno exacto con sus estructuras de partida expl\'icitas}.
Clasificaci\'on
\(III_{\lambda_\partial}\): \status{Teorema con hip\'otesis expl\'icitas de
persistencia}. Interpretaci\'on de la copia termocampo doble como exterior f\'isico:
\status{Realizaci\'on f\'isica condicionada}. La magnitud \(\gammaH\) se toma
del capítulo de Barbero--Immirzi y área mínima, donde fue derivada; aquí actúa
como dato estructural ya establecido.
\end{statusbox}

\endgroup
\begingroup
\let\chapter\section
\let\section\subsection
\let\subsection\subsubsection
\let\subsubsection\paragraph
\let\clearpage\relax
% La primera mitad de la fuente de mezcla coincide byte por byte con el
% desarrollo REV.2 ya publicado arriba. Se conserva su continuación exclusiva
% desde una extracción editor-ready y ambas procedencias permanecen intactas.
% Extracto editor-ready del contenido exclusivo de masas/12_mezcla.tex.
% Las líneas 1--104 están publicadas byte-exactamente desde
% 74_rev2_realizaciones_i.tex:254--296 y 406--468. El propietario íntegro
% permanece preservado; este extracto conserva su desarrollo exclusivo.

\subsection{Matriz PMNS como lector de mezcla de la fibra neutrínica común}
Sean \(F_\ell\) y \(F_\nu\) marcos ortonormales reales, con la misma imagen tridimensional y orientaciones elegidas de modo que \(\det(F_\ell^*F_\nu)=+1\), derivados de los operadores cargado y neutro. Entonces el transporte real es

\[
O=F_\ell^*F_\nu\in SO(3).
\]
Sin la elección correlativa de orientación sólo se concluye \(O\in O(3)\). Sea además \(b_K\) el testigo de Witt normalizado, \(Q_K=|b_K\rangle\langle b_K|\) y

\[
\Phi_L=I+(e^{i\pi/6}-1)Q_K.
\]
La fase mínima produce

\[
U_{\rm HMT}=F_\ell^*\Phi_LF_\nu.
\]
Esta compresión es unitaria sólo bajo la condición de compatibilidad

\[
\Phi_L(\operatorname{Im}F_\nu)=\operatorname{Im}F_\ell.
\]
Sin ella produce una contracción; cuando es invertible, su factor polar define el candidato unitario. Aun imponiendo la compatibilidad, la fase mínima de rango uno queda falsada como PMNS física: su mejor permutación deja \(\theta_{13}\simeq17.02^\circ\). El núcleo físico debe leer el registro completo:

\[
G_{ab}=
\frac1{\sqrt{|L_a||N_b|}}
\sum_{c\in L_a}\sum_{d\in N_b}K_{\rm reg}(c,d),
\qquad
U_{\rm int}=\operatorname{polar}(G),
\]
con \(\det G\ne0\). Acarreo, hoja, orientación, vacancia y holonomía forman parte del núcleo. El falsador de la fase de rango uno permanece impreso para evitar que la unitariedad formal se confunda con el cierre físico.

La condición Majorana de los neutrinos requiere además una autorconjugación fermiónica compatible con el operador de masa. La neutralidad, \(C_M\) y el sector topológico de Ising no bastan por separado.

\Needspace{7\baselineskip}
\subsection{Mezcla CKM y violación CP}
La gramática angular se organiza mediante

\[
M_{\rm CKM}=
\begin{pmatrix}
2&-5&28\\
0&7&-25/2\\
1&-20&-2/3
\end{pmatrix},
\qquad
\det M_{\rm CKM}=-\frac{3857}{6}.
\]
Con

\[
\begin{pmatrix}\theta_{12}\\\theta_{23}\\\theta_{13}\end{pmatrix}
=M_{\rm CKM}
\begin{pmatrix}A\\C^*/6\\{\left(180/\pi\right)\Delta_4}\end{pmatrix},
\qquad
\delta_{\rm CP}=9A,
\]
se obtiene

\[
(\theta_{12},\theta_{23},\theta_{13},\delta_{\rm CP})
=
(13.003313589509,
2.398056157332,
0.213977382244,
65.676173123554)^\circ.
\]
El invariante de Jarlskog es

\[
J=c_{12}c_{23}c_{13}^2s_{12}s_{23}s_{13}\sin\delta_{\rm CP}
=3.1189723326632974\times10^{-5}.
\]
El falsador lineal exacto es

\[
3857\delta_{\rm CP}
=9(1528\theta_{12}+3380\theta_{23}+801\theta_{13}).
\]
En la parametrización unitaria estándar,

\[
V_{\rm CKM}=
\begin{pmatrix}
c_{12}c_{13}&s_{12}c_{13}&s_{13}e^{-i\delta}\cr
-s_{12}c_{23}-c_{12}s_{23}s_{13}e^{i\delta}&
c_{12}c_{23}-s_{12}s_{23}s_{13}e^{i\delta}&s_{23}c_{13}\cr
s_{12}s_{23}-c_{12}c_{23}s_{13}e^{i\delta}&
-c_{12}s_{23}-s_{12}c_{23}s_{13}e^{i\delta}&c_{23}c_{13}
\end{pmatrix}.
\]
La evaluación interna anterior produce

\[
|V_{\rm CKM}|=
\begin{pmatrix}
0.974350259&0.225005836&0.003734601\cr
0.224873110&0.973489279&0.041841465\cr
0.008582400&0.041121745&0.999117283
\end{pmatrix}.
\]
Una comparación marginal congelada da desviaciones normalizadas \((-0.24,-0.07,-0.34,-0.31,-0.37)\) para los cuatro parámetros angulares y \(J\), respectivamente. Estas cifras no sustituyen la comparación covariante: la publicación debe identificar la edición externa y su matriz completa de covarianza antes de formar un \(\chi^2\) conjunto.

CKM transporta entre bases de quarks. Si \(B_u,B_d\) son los operadores diagonales de masa,

\[
B_d^{(u)}=V_{\rm CKM}B_dV_{\rm CKM}^*.
\]
La identidad

\[
|\det[B_u,V_{\rm CKM}B_dV_{\rm CKM}^*]|
=2|J|
\prod_{i<j}|\kappa_i^u-\kappa_j^u|
\prod_{a<b}|\kappa_a^d-\kappa_b^d|
\]
muestra que mezcla y espectro se acoplan sin que la matriz de mezcla cree los autovalores.

\Needspace{7\baselineskip}

\endgroup

\section{Composición, resonancias y sectores no ordinarios}
La masa compuesta se obtiene de los registros constituyentes, la frontera y el
operador de ligadura; no es una suma desnuda de masas. La anchura procede del
polo de supervivencia y permanece separada de la firma másica. Las secciones
virtuales, los sectores Fibonacci, Ising--Majorana y de orden seis, y las
exclusiones de taquiones y fantasmas conservan sus dominios propios.

\begingroup
\let\chapter\section
\let\section\subsection
\let\subsection\subsubsection
\let\subsubsection\paragraph
\let\clearpage\relax
\chapter{Composición de rutas, registros y firmas}
\label{chap:composicion-especies-rev6}
\index{composición de rutas}
\index{composición de registros}
\index{mesón!expresión composicional}
\index{barión!expresión composicional}

Una especie compuesta no puede entrar en la ley de masas como un nombre al
que se asigna retrospectivamente una fila de cinco enteros. Antes de la firma
debe existir un objeto que pueda componerse. El atlas del
\cref{chap:atlas-genealogia-masas} proporciona el alfabeto finito; la
extensión superviviente y el registro del
\cref{chap:extension-ruta-caracter} proporcionan la memoria que puede
componerse. La cuestión de este capítulo es construir el estrato intermedio
sin consultar masas objetivo:
\[
\begin{gathered}
 \text{configuración composicional tipada}
 \longrightarrow
 \text{registros completos de sus constituyentes}\\
 \longrightarrow
 \text{registro compuesto}
 \longrightarrow
 \ZZ^5.
\end{gathered}
\]

El orden de las flechas es parte del resultado. Las expresiones se forman
antes de asignarles una realización física; los registros se componen antes
de extraer la firma; la firma se evalúa mediante el carácter sólo después de
haber compuesto la memoria. De este modo, el dominio compuesto no queda
reducido a las especies incluidas en una tabla metrológica.

\section{Clasificación interna y gramática composicional}
\label{sec:alfabeto-composicional-rev6}

Sea
\[
 \mathcal C_{81}=I_9\times I_9,
 \qquad I_9=\{1,\ldots,9\},
\]
el atlas de celdas, y sea
\[
 \rho(r,c)=(10-r,10-c)
\]
su inversión central. Denotamos por
\(\mathcal L,\mathcal N,\mathcal D,\mathcal U\) los subconjuntos de
leptones cargados, neutrinos, quarks de tipo abajo y quarks de tipo arriba,
respectivamente, y ponemos
\[
 \mathcal Q=\mathcal D\sqcup\mathcal U.
\]
Cada celda posee un nivel combinatorio
\(\operatorname{gen}(x)\). En el sector quark se conserva además el residuo
de color
\[
 \kappa(x)=r(x)-c(x)\pmod3.
\]
Escribimos
\[
 \mathcal Q_a=\{q\in\mathcal Q:\kappa(q)=a\},
 \qquad a\in\mathbb Z/3\mathbb Z.
\]
Por la \cref{prop:color-atlas},
\[
 |\mathcal Q_0|=|\mathcal Q_1|=|\mathcal Q_2|=12.
\]

\begin{definition}[Dual formal e identidad de celda]
\label{def:dual-identidad-composicional}
Para una sección de celda orientada \(x\), su dual formal \(x^\vee\) invierte
la orientación de carga y de color, y tiene soporte \(\rho(x)\). La
notación conserva la distinción
\[
 x^\vee\ne \rho(x)
\]
cuando \(\rho(x)\) se considera como celda sin orientación. Para cada
\(x\in\mathcal C_{81}\), denotamos por
\(\operatorname{id}_x:x\to x\) la flecha identidad. El conjunto de
identidades es
\[
 \mathfrak I
 :=\{\operatorname{id}_x:x\in\mathcal C_{81}\}.
\]
\end{definition}

La identidad \(\operatorname{id}_x\) es una flecha neutra del sistema
composicional.
No es, por esa sola condición, una realización de un bosón neutro. Del mismo
modo, el dual formal incorpora orientación y no puede sustituirse por una
segunda celda desnuda.

\begin{definition}[Expresiones mesónicas y bariónicas]
\label{def:expresiones-mesonicas-barionicas}
El conjunto de expresiones mesónicas neutras en color es
\[
 \mathfrak M
 :=
 \left\{
   M(q,q')=q\otimes{(q')}^\vee:
   q,q'\in\mathcal Q,\ \kappa(q)=\kappa(q')
 \right\}.
\]
Los pares son ordenados. Su residuo total es
\(\kappa(q)-\kappa(q')=0\).

El conjunto de expresiones bariónicas con posiciones de color fijadas es
\[
 \mathfrak B
 :=
 \mathcal Q_0\times\mathcal Q_1\times\mathcal Q_2,
\]
y escribimos
\[
 B(q_0,q_1,q_2)=q_0\otimes q_1\otimes q_2.
\]
Su neutralidad residual es la identidad
\[
 0+1+2=0\pmod3.
\]
\end{definition}

Estas expresiones son objetos sintácticos tipados. En particular, no se ha
tomado todavía un cociente por permutaciones, espín, sabor, estadística,
energía de enlace o equivalencia dinámica.

\begin{definition}[Flechas cargadas compatibles]
\label{def:flechas-cargadas-compatibles}
Las flechas leptónicas orientadas son
\[
 \mathfrak A_\ell
 :=
 \left\{
 (x,y)\in
   (\mathcal L\times\mathcal N)
   \sqcup
   (\mathcal N\times\mathcal L):
 \operatorname{gen}(x)=\operatorname{gen}(y)
 \right\}.
\]
Las flechas quark orientadas son
\[
 \mathfrak A_q
 :=
 \left\{
 (x,y)\in
   (\mathcal U\times\mathcal D)
   \sqcup
   (\mathcal D\times\mathcal U):
 \begin{array}{l}
 \operatorname{gen}(x)=\operatorname{gen}(y),\\[-2pt]
 \kappa(x)=\kappa(y)
 \end{array}
 \right\}.
\]
Ponemos
\[
 \mathfrak A_{\rm ch}:=\mathfrak A_\ell\sqcup\mathfrak A_q.
\]
\end{definition}

La igualdad de nivel tipa la transición en ambos sectores; la igualdad de
color se exige además en el sector quark. La orientación distingue \(x\to y\)
de \(y\to x\).

\begin{definition}[Sistema finito de composiciones tipadas]
\label{def:gramatica-celular-compuesta}
El sistema composicional finito es la unión disjunta tipada
\[
 \mathfrak G_{\rm comp}
 :=
 \mathcal C_{81}
 \sqcup\mathfrak I
 \sqcup\mathfrak M
 \sqcup\mathfrak B
 \sqcup\mathfrak A_{\rm ch}.
\]
Las \(81\) celdas son los generadores primitivos. Las \(81\) identidades son
flechas sobre esos mismos generadores y no un segundo atlas.
\end{definition}

\section{Censo combinatorio de composiciones de rutas independiente de datos metrológicos}
\label{sec:censo-composicional-rev6}

\begin{theorem}[Censo del sistema finito de composiciones]
\label{thm:censo-gramatica-composicional-rev6}
El sistema de la \cref{def:gramatica-celular-compuesta} contiene
\[
 \boxed{\begin{gathered}
        81\ \text{identidades},\qquad
        432\ \text{expresiones mesónicas},\\
        1728\ \text{expresiones bariónicas},\qquad
        372\ \text{flechas cargadas}
 \end{gathered}}
\]
Las flechas cargadas se descomponen exactamente como
\[
 372=320_{\text{leptónicas}}+52_{\text{quark}}.
\]
El censo depende sólo del atlas, de su inversión, de los niveles y del
residuo de color. No utiliza masas, unidades ni firmas publicadas.
Los cuatro cardinales cuentan expresiones composicionales internas o flechas
admisibles del sistema; no cuentan especies físicas distintas, partículas
físicas ni masas.
\end{theorem}

\begin{proof}
La asignación \(x\mapsto\operatorname{id}_x\) es biyectiva, de modo que
\[
 |\mathfrak I|=|\mathcal C_{81}|=81.
\]
Para los mesones, la condición de neutralidad permite escoger
independientemente un par ordenado dentro de cada clase de color. Como cada
clase contiene doce celdas,
\[
 |\mathfrak M|
 =\sum_{a\in\mathbb Z/3\mathbb Z}|\mathcal Q_a|^2
 =3\cdot12^2
 =432.
\]
Para los bariones, las posiciones \(0,1,2\) están fijadas, y por tanto
\[
 |\mathfrak B|
 =|\mathcal Q_0||\mathcal Q_1||\mathcal Q_2|
 =12^3
 =1728.
\]

Las familias leptónicas cargadas tienen tamaños
\[
 |\mathcal L_{G0}|=1,\qquad
 |\mathcal L_{G2}|=8,\qquad
 |\mathcal L_{G4}|=16,
\]
mientras las familias neutrínicas tienen tamaños
\[
 |\mathcal N_{G1}|=2,\quad
 |\mathcal N_{G2}|=4,\quad
 |\mathcal N_{G3}|=6,\quad
 |\mathcal N_{G4}|=8.
\]
Sólo \(G2\) y \(G4\) son niveles comunes. Como se cuentan ambas
orientaciones,
\[
 |\mathfrak A_\ell|
 =2(8\cdot4+16\cdot8)
 =320.
\]

Para las flechas quark, en el orden de colores \((0,1,2)\), las
multiplicidades de tipo arriba y tipo abajo en los niveles comunes son
\[
\begin{array}{ccc}
 &\mathcal U&\mathcal D\\
 G1&(2,1,1)&(0,1,1)\\
 G3&(4,4,4)&(2,2,2).
\end{array}
\]
El número de pares de una orientación y del mismo color es, por ello,
\[
 (2\cdot0+1\cdot1+1\cdot1)
 +(4\cdot2+4\cdot2+4\cdot2)
 =2+24=26.
\]
Al incluir la orientación contraria se obtiene
\[
 |\mathfrak A_q|=2\cdot26=52.
\]
Finalmente,
\[
 |\mathfrak A_{\rm ch}|=320+52=372.
\]
Ninguna de estas reglas contiene un valor objetivo o un vector de masa.
\end{proof}

\begin{corollary}[Constructor composicional sin blancos]
\label{cor:constructor-composicional-sin-blancos}
El constructor que enumera conjuntamente identidades, expresiones mesónicas,
expresiones bariónicas y flechas cargadas factoriza por los campos tipados
\[
 (\text{sector},\text{familia},\text{nivel},\kappa,\rho)
\]
del atlas. Con la notación anterior, sus cuatro codominios son
\(\mathfrak I,\mathfrak M,\mathfrak B,\mathfrak A_{\rm ch}\).
En particular, el censo puede fijarse antes de abrir cualquier tabla de firmas.
\end{corollary}

\begin{proof}
Las cuatro definiciones y el cálculo del
\cref{thm:censo-gramatica-composicional-rev6} emplean únicamente esos campos.
La tabla de firmas no pertenece al dominio de ninguna de las reglas.
\end{proof}

La independencia de blancos no significa que las expresiones hayan adquirido
ya un nombre físico. Significa algo anterior y más preciso: existe un dominio
composicional construido sin usar como criterio la cantidad que después se
evaluará.

\section{Transporte de memoria desde la ruta compuesta hasta la firma integral}
\label{sec:memoria-composicion-especies-rev6}

Sea \(\gamma\in\Gamma_{108}^{\rm enr}\) una presentación enriquecida de una
extensión. Su registro leído contiene
\[
 \mathcal E(\gamma)
 =
 (n_A,e_0,\ldots,e_{11},T_\zeta,C_\zeta),
\]
además de los datos de fase, hoja, orientación y borde que tipan una composición.
Los doce eventos no son reemplazables por su total antes de componer. Para
\(0\le i<6\), definimos
\[
 p_i=e_i+e_{i+6},
 \qquad
 a_i=e_i-e_{i+6},
\]
\[
 s_C=\sum_{i=0}^{5}p_i,
 \qquad
 M_i=p_i-p_5\quad(0\le i<5).
\]
Con \(M_5=(M_0,\ldots,M_4)\) y
\(A_6=(a_0,\ldots,a_5)\), la reconstrucción es
\[
 p_5=\frac{s_C-\sum_{i=0}^{4}M_i}{6},
 \qquad
 p_i=p_5+M_i,
\]
\[
 e_i=\frac{p_i+a_i}{2},
 \qquad
 e_{i+6}=\frac{p_i-a_i}{2}.
\]
Por el \cref{thm:memoria-1-5-6}, sobre la imagen integral se obtiene la
descomposición reversible
\[
 \boxed{(s_C,M_5,A_6),\qquad 1+5+6=12.}
\]
El primer bloque conserva el total visible; el segundo, la forma relativa
dentro de la hoja simétrica; el tercero, la memoria entre las dos semicoronas.

La coordenada torsional fuerte requiere además la cocadena a nivel de paso:
\[
 k_{\Delta_4}=T_\zeta-2C_\zeta.
\]
No es una función del mero patrón de eventos nulos. Esta diferencia será
esencial al interpretar el rombo de no descenso.

\begin{definition}[Dato de composición orientada]
\label{def:dato-empalme-orientado-rev6}
Sean \(\Gamma,\Lambda\) dos registros enriquecidos. Un dato de composición es un
par \((g,B)\), donde \(g\) alinea fase, hoja, orientación y posición en la
interfaz, y \(B\) es la subtraza común, con su contribución a la acción, sus
doce eventos, sus predecesores torsionales y su corona. Sobre el dominio en
que esos datos son compatibles se define
\[
 \boxed{
 {(n,e)}_\Gamma\star_{(g,B)}{(n,e)}_\Lambda
 =
 \bigl(
 n_\Gamma+n_\Lambda-n_B,\,
 e_\Gamma+g e_\Lambda-e_B
 \bigr).}
\]
La palabra torsional y la corona se componen en el mismo nivel, antes de
calcular \(T_\zeta,C_\zeta\) y \(k_{\Delta_4}\).
\end{definition}

La fórmula es una inclusión-exclusión de historias sobre una subtraza
identificada. Está demostrada con estructura de partida explícita: los
registros, el transporte \(g\), el borde \(B\), la orientación y la palabra
de predecesores forman parte de la entrada. El sistema composicional no elige por
sí sola ninguno de esos datos.

\begin{definition}[Extractor posterior a la composición]
\label{def:extractor-posterior-empalme-rev6}
Para un registro enriquecido completo se define
\[
 L_{\rm tick}(\gamma)
 =
 \bigl(
 n_A,s_C,k_{\Delta_4},
 \nu_{120}^{\rm ev},\nu_{270}
 \bigr)\in\ZZ^5,
\]
donde, si \(\tau_m=\operatorname{sgn}(e_m)\) e
\(I_a=\{a,a+4,a+8\}\pmod {12}\),
\[
 \nu_{120}^{\rm ev}
 =
 \sum_{a=0}^{3}
 \operatorname{sgn}\left(\sum_{m\in I_a}e_m\right),
 \qquad
 \nu_{270}
 =
 \sum_{m=0}^{11}\tau_m\tau_{m+3}.
\]
La firma de una composición se define por
\[
 \operatorname{Sig}(\Gamma,\Lambda;g,B)
 :=
 L_{\rm tick}
 \bigl(\Gamma\star_{(g,B)}\Lambda\bigr).
\]
\end{definition}

Los signos y las autocorrelaciones se aplican al registro ya compuesto. Por
ello, la definición no equivale a sumar las firmas de los constituyentes.

\section{Diagrama conmutativo de composición y preservación de la fibra}
\label{sec:rombo-no-descenso-composicional-rev6}

El catálogo histórico B1 emplea una proyección de eventos que conviene
tipar por separado. Sea
\[
 \Sigma_{12}=(+,+,+,-,-,-,+,+,+,-,-,-),
 \qquad \tau_m=\operatorname{sgn}(e_m),
\]
y definamos la proxy de eventos nulos
\[
 K_0(e)
 =
 \sum_{\tau_m=0}\Sigma_{12}(m).
\]
La proyección exacta es
\[
 F_0(n_A,e)
 =
 \left(
 n_A,\,
 \sum_m e_m,\,
 K_0(e),\,
 \sum_{a=0}^{3}
   \operatorname{sgn}\sum_{m\in I_a}e_m,\,
 \sum_m\tau_m\tau_{m+3}
 \right)\in\ZZ^5.
\]
Para dos registros alineados, escribimos
\[
 (n,e)\boxplus(n',e')=(n+n',e+e').
\]
Esta operación es el caso sin subtraza común y con alineamiento identidad del
caso de composición de eventos.

\begin{proposition}[Obstrucción al descenso de la proyección B1]
\label{prop:no-descenso-proyeccion-b1-rev6}
No existe una operación binaria
\[
 \mu:\ZZ^5\times\ZZ^5\longrightarrow\ZZ^5
\]
tal que
\[
 F_0(\Gamma\boxplus\Lambda)
 =
 \mu\bigl(F_0(\Gamma),F_0(\Lambda)\bigr)
\]
para todos los registros B1.
\end{proposition}

\begin{proof}
Consideremos las rutas, escritas en la carta
\((r,c,h_0,v_0)\),
\[
 a=(1,3,-1,+1),
 \qquad
 b=(6,7,-1,+1),
 \qquad
 c=(1,2,-1,-1).
\]
Las dos primeras tienen registros de eventos distintos:
\[
\begin{aligned}
 e(a)&=(5,4,4,-4,-4,-6,\;5,4,4,-4,-4,-6),\\
 e(b)&=(4,5,4,-4,-6,-4,\;4,5,4,-4,-6,-4),
\end{aligned}
\]
pero publican la misma cinco-tupla:
\[
 F_0(a)=F_0(b)=(8,-2,0,0,-12).
\]
La ruta de prueba tiene
\[
 e(c)=(1,4,1,-2,-4,-3,\;1,4,1,-2,-4,-3)
\]
y
\[
 F_0(c)=(28,-6,0,-4,-12).
\]
Al sumar primero los eventos y reevaluar después, se obtiene
\[
 F_0(a\boxplus c)=(36,-8,0,0,-12),
\]
\[
 F_0(b\boxplus c)=(36,-8,0,-2,-12).
\]
Los dos pares de entradas de la supuesta \(\mu\) son iguales, pero las
salidas difieren en \(\nu_{120}^{\rm ev}\). Esto contradice la existencia de
\(\mu\).
\end{proof}

El censo exhaustivo del mismo dominio contiene \(324\) rutas, \(53\)
cinco-tuplas proxy y \(198\) estados de eventos distintos. Cuarenta y dos
fibras de la proyección contienen más de un estado enriquecido; cuarenta de
ellas son separadas por alguna ruta de prueba. El rombo anterior es el testigo
canónico más breve de esa pérdida de información.

\begin{remark}[Separación tipada entre el representante \(K_0\) y la torsión fuerte]
\label{rem:proxy-no-torsion-fuerte-rev6}
La tercera coordenada de \(F_0\) es \(K_0(e)\), que en B1 toma valores en
\(\{-4,-2,0,2,4\}\). No se identifica con
\(k_{\Delta_4}=T_\zeta-2C_\zeta\). De hecho, existen rutas con el mismo
\(n_A\) y el mismo vector de doce eventos a las que el lector torsional
asigna valores distintos porque conserva los predecesores de cada iteración
y la corona.

La \cref{prop:no-descenso-proyeccion-b1-rev6} demuestra exactamente que la
proyección B1 de cinco coordenadas no transporta la superposición de
registros. No afirma una imposibilidad universal para toda estructura futura
sobre \(\ZZ^5\), ni autoriza a sustituir \(K_0\) por
\(k_{\Delta_4}\). En el lector fuerte, la salida se define mediante la
\cref{def:extractor-posterior-empalme-rev6}: primero se compone la palabra
enriquecida y después se reextraen sus cinco coordenadas.
\end{remark}

\section{Teorema de composición de registros independiente de datos metrológicos}
\label{sec:teorema-composicional-sin-blancos-rev6}

\begin{theorem}[Cierre composicional sin blancos]
\label{thm:cierre-composicional-sin-blancos-rev6}
Fijados el atlas \(\mathcal C_{81}\), la dualidad formal, el residuo de color
y el registro enriquecido del \cref{chap:extension-ruta-caracter}, se cumplen
las afirmaciones siguientes.
\begin{enumerate}
 \item Existe el sistema finito de composiciones tipadas
 \(\mathfrak G_{\rm comp}\) con \(81\) identidades, \(432\) expresiones
 mesónicas, \(1728\) expresiones bariónicas y \(372\) flechas cargadas,
 descompuestas como \(320+52\).

 \item El sistema se construye sin masas objetivo. Su salida es una
 expresión, órbita o multisección tipada, no una masa ni una etiqueta física
 individualizada.

 \item Dada una expresión y dados registros constituyentes con un árbol
 explícito de composiciones compatibles \((g,B)\), la operación produce un
 registro enriquecido. Su firma se obtiene por la aplicación
 \[
 \boxed{
 \bigl(
 \text{expresión},
 \text{registros},
 \text{composiciones orientadas}
 \bigr)
 \xrightarrow{\ \mathfrak C_{12}\ }
 \text{registro compuesto}
 \xrightarrow{\ L_{\rm tick}\ }
 \ZZ^5.}
 \]
 Aquí \(\mathfrak C_{12}\) denota la iteración de
 \(\star_{(g,B)}\) con el parentizado y los bordes incluidos en los datos.

 \item La proyección B1 a cinco enteros no constituye por sí sola un espacio
 composicional: el diagrama con \(\boxplus\) no desciende a
 \(\ZZ^5\). Por tanto, reemplazar cada registro por su cinco-tupla antes del
 composición pierde información necesaria para reevaluar al menos
 \(\nu_{120}^{\rm ev}\).
\end{enumerate}
\end{theorem}

\begin{proof}
El primer punto es el
\cref{thm:censo-gramatica-composicional-rev6}. El segundo es el
\cref{cor:constructor-composicional-sin-blancos}. Para el tercero, cada
arista del árbol de composición se evalúa mediante la
\cref{def:dato-empalme-orientado-rev6}; como los bordes y transportes forman
parte del dominio, la iteración produce un registro sobre el cual está
definido \(L_{\rm tick}\). No se utiliza una firma intermedia como sustituto
del registro. El cuarto punto es la
\cref{prop:no-descenso-proyeccion-b1-rev6}.
\end{proof}

La fórmula correcta de la firma compuesta es, por tanto,
\[
 \operatorname{Sig}_{\rm comp}
 =
 L_{\rm tick}\circ\mathfrak C_{12}.
\]
No es una suma de filas de masa. Sólo después puede aplicarse el carácter
formal o su especialización global:
\[
 \text{registro compuesto}
 \xrightarrow{L_{\rm tick}}\ZZ^5
 \xrightarrow{\chi_{\rm form}}\operatorname{Fun}(\mathcal P,\RR_{>0})
 \longrightarrow
 \text{carta dimensional}.
\]

\section{Reconstrucción genealógica de mezcla, generación y persistencia de sabor}
\label{sec:significado-genealogico-composicion-rev6}

El teorema anterior completa un estrato que una lista de masas no puede
mostrar. Las celdas primitivas no se limitan a recibir nombres; generan
identidades, pares duales, triples de color y flechas orientadas mediante
reglas internas. Los objetos compuestos aparecen así dentro del dominio antes
de toda comparación metrológica.

La genealogía tampoco se reduce al árbol sintáctico. Dos expresiones con el
mismo número de constituyentes pueden componerse de maneras distintas porque
la fase, la hoja, la orientación y la subtraza común forman parte del dato.
El registro conserva esas diferencias hasta que han actuado los signos, las
correlaciones y la torsión. La masa, cuando exista una realización
dimensional, será una lectura posterior de esa memoria compuesta.

Este orden explica por qué \(\ZZ^5\) no es el dominio de composición aunque
sea el dominio del carácter. El grupo \((\ZZ^5,+)\) existe y el carácter es
multiplicativo sobre él; lo que no existe para la proyección B1 es una ley
que recupere desde las cinco-tuplas la superposición que ya olvidó la
colocación de los doce eventos. La linealidad del carácter no devuelve la
información eliminada por el extractor.

\section{Del sistema composicional interno a la realización física}
\label{sec:delimitacion-composicion-especies-rev6}

El \cref{thm:cierre-composicional-sin-blancos-rev6} construye el dominio
anterior a los nombres físicos: expresiones ordenadas, duales, triples de
color y flechas compatibles. Sus cardinales cuentan posibilidades internas
de composición. Una especie física, en cambio, es una representación persistente
con sabor, color, espín, estadística y canal determinados. Confundir ambos
niveles convertiría un censo de expresiones en un falso catálogo de
partículas.

\begin{definition}[Dato de realización física]
\label{def:dato-realizacion-fisica-rev6}
Un dato de realización para una especie \(X\) es la tupla
\[
 \mathfrak D_X=
 \bigl(
 \mathfrak t_X,\mathfrak s_X,
 (\mathcal L_{X_i})_{i=1}^{r},
 \mathcal T_X,(g_X,B_X)
 \bigr),
\]
donde:
\begin{enumerate}
 \item \(\mathfrak t_X\) fija el tipo físico ---elemental, nulo, mesónico,
 bariónico, multiquark o topológico--- y su representación;
 \item \(\mathfrak s_X\) es la multisección o expresión interna realizada;
 \item \(\mathcal L_{X_i}\) son los registros enriquecidos de los
 constituyentes;
 \item \(\mathcal T_X\) fija sabor, color, espín, estadística y
 simetrización;
 \item \((g_X,B_X)\) fija el árbol de composición, sus transportes y sus
 subtrazas comunes.
\end{enumerate}
La tupla se declara sin consultar la masa o la anchura que después se
evaluarán.
\end{definition}

\begin{proposition}[La realización cambia el registro, no la ley]
\label{prop:realizacion-compuesta-ley-unica-rev6}
Todo dato \(\mathfrak D_X\) determina un registro compuesto
\[
 \mathcal L_X
 =\mathfrak C_{12}
   \bigl((\mathcal L_{X_i})_{i=1}^{r};
         \mathcal T_X,g_X,B_X\bigr),
\]
una firma \(v_X=L_{\rm tick}(\mathcal L_X)\) y el carácter central universal
\[
 \chi_0(v_X-v_e).
\]
Si la genealogía aporta además un refinamiento espectral tipado
\(\eta_X\in\mathcal E_{\mathfrak S}\), la evaluación completa es
\[
 m_X
 =m_e^{\rm HMT}\,
   \chi_{\rm full}(v_X-v_e,\eta_X).
\]
El carácter central y el refinamiento pertenecen a niveles distintos; la
completitud de la torre no sustituye el mapa prospectivo que produzca
\(\eta_X\).
Si \(X\) es inestable, su anchura pertenece al lector temporal
\(\Gamma_X=\hbar_{\rm int}/\tau_X^{\rm life}\), no a una sexta coordenada de
la firma de masa.
\end{proposition}

\begin{proof}
El dato de realización contiene exactamente las entradas de
\(\mathfrak C_{12}\). El
\cref{thm:cierre-composicional-sin-blancos-rev6} produce el registro
enriquecido y la \cref{def:extractor-posterior-empalme-rev6} produce su
firma. El carácter central y su elevación espectral son los de
\cref{eq:caracter-masa-global-rev7,eq:ley-masa-todos-registros-rev6}. La
anchura usa una recta temporal y por ello tiene otro dominio.
\end{proof}

La separación de niveles queda así fijada:
\[
\boxed{
\begin{gathered}
\text{sistema composicional interno}
\longrightarrow\text{dato de realización}
\longrightarrow\text{registro compuesto}\\
\longrightarrow\text{firma}
\longrightarrow\text{carácter central}
\longrightarrow\text{refinamiento espectral}
\longrightarrow\text{comparación física}.
\end{gathered}}
\]
Las \(432\) expresiones mesónicas, las \(1728\) bariónicas y las \(372\)
flechas cargadas ocupan sólo la primera casilla y no constituyen un
inventario de partículas. El Modelo Estándar conserva su conjunto tipado de
especies y cada especie ocupa la cadena mediante su propio dato de
realización. La regla causal es única en todos los sectores:
\[
 \boxed{\text{componer extensiones y registros;\quad leer la firma después}.}
\]

\endgroup
\begingroup
\let\chapter\section
\let\section\subsection
\let\subsection\subsubsection
\let\subsubsection\paragraph
\let\clearpage\relax
\chapter{Banda de partícula, anti-sección y virtualidad}
\label{chap:banda-particula-virtual}
\index{partícula!banda de ruta}
\index{antipartícula!anti-sección orientada}
\index{partícula virtual}

La persistencia de una ruta produce una partícula antes de que se evalúe su
masa. Esta prioridad permite distinguir tres objetos:

\begin{enumerate}
\item la \emph{sección orientada}, que conserva la dirección de la ruta;
\item la \emph{banda}, que reúne una sección y su conjugada;
\item la \emph{persistencia}, que decide si la firma alcanza o no una rama
global.
\end{enumerate}

\begin{definition}[Palabras dodecafásicas y conjugación material]
\label{def:palabras-cm}
Sea
\[
 \mathscr T_{12}^{\rm word}:=\{-1,0,1\}^{12}.
\]
Para \(\mathbf t=(t_1,\ldots,t_{12})\) escribimos
\[
 \operatorname{rev}(\mathbf t)
 =(t_{12},\ldots,t_1),
 \qquad
 C_M(\mathbf t):=-\operatorname{rev}(\mathbf t).
\]
La inversión \(C_M\) cambia simultáneamente el sentido de lectura y la
orientación trítica.
\end{definition}

\begin{theorem}[Involución material y sector fijo]
\label{thm:involucion-cm-censo}
La aplicación \(C_M\) es una involución de
\(\mathscr T_{12}^{\rm word}\) y
\[
 \left|\operatorname{Fix}(C_M)\right|=3^6=729.
\]
\end{theorem}

\begin{proof}
La reversión y el cambio de signo conmutan y ambos tienen cuadrado igual a
la identidad; por tanto \(C_M^2=I\). Una palabra es fija si y sólo si
\[
 t_i=-t_{13-i},\qquad 1\leq i\leq6.
\]
Las seis primeras coordenadas son libres y las seis últimas quedan
determinadas por ellas. Cada coordenada libre admite tres valores, de modo
que el sector fijo tiene cardinal \(3^6\).
\end{proof}

Esta conjugación finita no se obtiene cambiando únicamente el signo visible:
invierte también la orientación del registro.

\section{La banda de ruta}

Sea \(\mathscr T_{12}^{\rm word}\) el espacio de palabras de la
\cref{def:palabras-cm}. Añadimos una orientación de sección
\(\chi\in\{-1,+1\}\) y definimos
\[
(\mathbf t,\chi)\sim(C_M\mathbf t,-\chi).
\]

\begin{definition}[Banda dodecafásica de ruta]
\label{def:banda-dodecafasica-ruta}
La banda de la sección orientada \((\mathbf t,\chi)\) es la clase
\[
[(\mathbf t,\chi)]_M
=\{(\mathbf t,\chi),(C_M\mathbf t,-\chi)\},
\]
y el espacio de bandas es
\[
\mathcal B_{12}
:=
\bigl(\mathscr T_{12}^{\rm word}\times\{-1,+1\}\bigr)/{\sim}.
\]
\end{definition}

\begin{proposition}[Libertad de la conjugación orientada]
\label{prop:accion-cm-orientada-libre}
La acción
\[
(\mathbf t,\chi)\longmapsto(C_M\mathbf t,-\chi)
\]
es libre. En consecuencia,
\[
|\mathcal B_{12}|=3^{12}.
\]
\end{proposition}

\begin{proof}
Un punto fijo exigiría \(\chi=-\chi\), imposible para
\(\chi\in\{-1,+1\}\). Todas las órbitas tienen dos elementos y el dominio
posee \(2\cdot3^{12}\).
\end{proof}

La banda no borra la orientación: la conserva como doble cubierta. La
partícula y la antipartícula son las dos secciones orientadas de la misma
banda.

\section{Masa par y carga impar}

Sea \(E_0:\mathscr T_{12}^{\rm word}\to\mathbb R\) un carácter aditivo de
ruta. Sus componentes par e impar respecto de \(C_M\) son
\[
E_+(\mathbf t)
=\frac12\bigl(E_0(\mathbf t)+E_0(C_M\mathbf t)\bigr),
\]
\[
E_-(\mathbf t)
=\frac12\bigl(E_0(\mathbf t)-E_0(C_M\mathbf t)\bigr).
\]
Entonces
\[
E_+(C_M\mathbf t)=E_+(\mathbf t),
\qquad
E_-(C_M\mathbf t)=-E_-(\mathbf t).
\]

\begin{definition}[Carácter de banda]
\label{def:caracter-banda}
Para una sección orientada definimos
\[
E_{\rm band}(\mathbf t,\chi)
=E_+(\mathbf t)+\chi E_-(\mathbf t),
\qquad
m(\mathbf t,\chi)
=m_\circ\exp\bigl(E_{\rm band}(\mathbf t,\chi)\bigr),
\]
donde \(m_\circ\) es la unidad interna de comparación de la familia.
\end{definition}

\begin{theorem}[Principio de conjugación de masa]
\label{thm:principio-conjugacion-masa-banda}
El carácter y la masa descienden al cociente de bandas:
\[
E_{\rm band}(C_M\mathbf t,-\chi)
=E_{\rm band}(\mathbf t,\chi),
\]
\[
\boxed{
m(C_M\mathbf t,-\chi)=m(\mathbf t,\chi).}
\]
Los invariantes pares pertenecen a la banda; los invariantes impares
distinguen la orientación de sus dos secciones.
\end{theorem}

\begin{proof}
Por las paridades anteriores,
\[
\begin{aligned}
E_{\rm band}(C_M\mathbf t,-\chi)
&=E_+(C_M\mathbf t)-\chi E_-(C_M\mathbf t)\\
&=E_+(\mathbf t)+\chi E_-(\mathbf t).
\end{aligned}
\]
La igualdad de masas se obtiene al aplicar la exponencial.
\end{proof}

En particular, una forma lineal con pesos simétricos es impar, mientras una
forma cuadrática compatible con la reversión es par. La carga y la
orientación cambian; la masa, construida sobre la banda, permanece.

\begin{corollary}[Electrón y positrón]
\label{cor:electron-positron-banda}
En la ecuación electrónica de banda, el par
\((n,\chi)=(-22,+1)\) y su conjugado \((+22,-1)\) producen el mismo término
\(\chi n=-22\). Por tanto, electrón y positrón poseen la misma masa y
orientaciones de carga opuestas.
\end{corollary}

\section{Cierre finito del atlas bajo anti-sección}

En las ochenta y una rutas CT--108 aparecen siete formas tríticas. El censo
exhaustivo distingue tres operaciones:
\[
\mathbf t\longmapsto-\mathbf t,\qquad
\mathbf t\longmapsto\operatorname{rev}(\mathbf t),\qquad
\mathbf t\longmapsto-\operatorname{rev}(\mathbf t).
\]

\begin{proposition}[Cierre de las formas CT--108]
\label{prop:cierre-antiseccion-ct108}
El conjunto de formas tríticas de CT--108 no es estable bajo signo puro ni
bajo reversión pura, y es estable bajo \(C_M=-\operatorname{rev}\). Sus
órbitas poseen multiplicidades
\[
\boxed{81=54+9+9+9.}
\]
La órbita de multiplicidad \(54\) es auto-conjugada; las otras tres
intercambian pares de formas.
\end{proposition}

\begin{proof}
La forma fija es
\[
\texttt{+++---+++---}
\]
y posee multiplicidad \(54\). Las tres parejas
\[
\begin{array}{c}
\texttt{+++--0+++--0}\ /\ \texttt{0++---0++---},\\
\texttt{+++-0-+++-0-}\ /\ \texttt{+0+---+0+---},\\
\texttt{+++0--+++0--}\ /\ \texttt{++0---++0---}
\end{array}
\]
poseen multiplicidad total \(9\) cada una. La aplicación
\(-\operatorname{rev}\) fija la primera y permuta los miembros de cada
pareja. El cambio de signo y la reversión por separado producen palabras
fuera de esta lista. La suma de multiplicidades es
\(54+9+9+9=81\).
\end{proof}

El resultado no identifica el sector \(54\) con otro objeto de cardinal
homónimo. Demuestra qué conjugación respeta la taxonomía concreta de las
rutas y por qué partícula y antipartícula comparten banda.

\section{Persistencia finita y partícula virtual}

Sea
\[
\cdots\xrightarrow{\rho_{M+1,M}}S_M
\xrightarrow{\rho_{M,M-1}}S_{M-1}\xrightarrow{}\cdots
\]
el sistema de estados admisibles a profundidades finitas. Para
\(s_N\in S_N\), definimos el conjunto de prolongaciones
\[
\mathcal P_M(s_N)
:=
\{s_M\in S_M:\rho_{M,N}(s_M)=s_N\},
\qquad M\geq N.
\]
Su profundidad de vida es
\[
\mathcal L(s_N)
:=\sup\{M\geq N:\mathcal P_M(s_N)\neq\varnothing\}
\in\mathbb N\cup\{\infty\}.
\]

\begin{definition}[Persistencia y virtualidad]
\label{def:persistencia-virtualidad}
Un prefijo \(s_N\) es \emph{persistente} si pertenece a una sección global
\((s_M)_{M\geq N}\) compatible. Es \emph{virtual finito} si
\(\mathcal L(s_N)<\infty\), existe una prolongación hasta
\(\mathcal L(s_N)\) y no existe ninguna a la profundidad siguiente.
\end{definition}

\begin{theorem}[Criterio de persistencia]
\label{thm:criterio-persistencia-virtual}
Supóngase que cada estado posee un número finito de prolongaciones inmediatas
y que las restricciones son coherentes. Entonces:
\[
\mathcal L(s_N)=\infty
\quad\Longleftrightarrow\quad
s_N\text{ pertenece a una sección global}.
\]
Si \(\mathcal L(s_N)=L<\infty\), el prefijo transporta ruta, memoria,
registro y firma hasta \(L\), pero no define una partícula persistente.
\end{theorem}

\begin{proof}
Una sección global proporciona una prolongación a toda profundidad, de modo
que \(\mathcal L=\infty\). Recíprocamente, las prolongaciones de \(s_N\)
forman un árbol enraizado, infinito y finitamente ramificado. Por el lema de
Kőnig contiene una rama infinita, cuyos vértices son compatibles por
construcción y definen una sección global.

Si \(\mathcal L=L<\infty\), existe un estado a profundidad \(L\) y ninguno
a \(L+1\). Los lectores definidos sobre prefijos pueden evaluarlo hasta
\(L\); la ausencia de prolongación impide convertir esa firma finita en una
sección del límite inverso.
\end{proof}

La partícula virtual posee así una definición interna. No es una partícula
``menos real'' ni una violación breve de una ley de conservación. Es una
ruta cuya memoria influye durante una profundidad finita y cuya genealogía
se extingue antes de estabilizarse. El contraste con propagadores y estados
intermedios de la física convencional se realiza después de identificar qué
lector representa esas profundidades.

\section{Cadena ruta–registro–firma–carácter–torre de la realización material}

Los resultados de este capítulo fijan la dirección completa:
\[
\boxed{
\text{extensión}\to
\text{ruta}\to
\text{registro}\to
\text{firma}\to
\text{banda orientada}\to
\text{partícula}\to
\text{masa}.}
\]
La antipartícula es la segunda sección de la banda. La virtualidad es una
persistencia finita. Ninguna de estas nociones se define a partir de una
masa observada.

\endgroup
\begingroup
\let\chapter\section
\let\section\subsection
\let\subsection\subsubsection
\let\subsubsection\paragraph
\let\clearpage\relax
% Vista científica exacta materializada desde una rama condicional.
% Fuente: source/demostraciones_primarias/31_06c_sectores_topologicos_no_omision.tex; rama: HMTIncludeCMBlock.
% No modifica el contenido matemático de la rama de origen.
\chapter[Estructura compleja \texorpdfstring{\(\widetilde J_M^{\,2}=-I\)}{J M al cuadrado = -I} y sectores topológicos]
{Estructura compleja \texorpdfstring{\(\widetilde J_M^{\,2}=-I\)}{J M al cuadrado = -I} sobre el estrato fijo, involuciones dodecafásicas y sectores topológicos}
\label{chap:sectores-topologicos-no-omision}
\index{sector topológico}\index{Fibonacci, anillo basado}
\index{Majorana}\index{Ising}\index{grupo de trenzas}

La involución \(C_M\), el anillo basado asociado a \(F_{\rm av}\), los
caracteres de trenza y las categorías apuntadas pertenecen a estratos
algebraicos distintos. Se presentan con sus dominios y controles de
identificación para que una coincidencia nominal no sustituya al morfismo que
debería relacionarlos. El censo del atlas, la jerarquía espectral infinita
cuyas ocho alturas históricas son testigos iniciales, y las completaciones
estadísticas conservan las demostraciones de
\cref{thm:censo-atlas-81,thm:descenso-necesario-fav,%
thm:principio-exclusion-pauli,thm:dos-cuatro-pi-estadistica}.

Cada resultado declara directamente su dominio, su construcción y su
conclusión; ninguna coincidencia nominal modifica esas dependencias.

\section{Estrato periódico de la involución dodecafásica}
\label{sec:involucion-cm-censo}

El dominio \(\mathscr T_{12}^{\rm word}\), la conjugación
\(C_M\mathbf t=-\operatorname{rev}(\mathbf t)\), su carácter involutivo y el
censo \(\lvert\operatorname{Fix}(C_M)\rvert=3^6=729\) tienen su residencia
primaria en las
\cref{def:palabras-cm,thm:involucion-cm-censo}, dentro de la banda de
partícula. Aquí se añade únicamente la condición periódica necesaria para el
estrato topológico. Si \(r\) denota el desplazamiento cíclico
\((r\mathbf t)_m=t_{m-1}\), ponemos
\[
 \operatorname{Per}_6
 :=\{\mathbf t\in\mathscr T_{12}^{\rm word}:r^6\mathbf t=\mathbf t\}.
\]

\begin{proposition}[Censo del estrato periódico]
\label{prop:estrato-periodico-cm-rev7}
Al imponer periodicidad de período divisor de seis,
\[
 \bigl|\operatorname{Fix}(C_M)\cap\operatorname{Per}_6\bigr|
 =3^3=27.
\]
Esta identidad es exacta.
\end{proposition}

\begin{proof}
El \cref{thm:involucion-cm-censo} da la forma general de los puntos fijos.
Si además \(r^6\mathbf t=\mathbf t\), la segunda mitad de la palabra coincide
con la primera. Escribiendo las seis primeras coordenadas, las dos
condiciones equivalen a
\[
 (t_0,t_1,t_2,t_3,t_4,t_5)
 =(a,b,c,-c,-b,-a).
\]
Los parámetros \(a,b,c\) son independientes y toman tres valores; por tanto,
la intersección tiene \(3^3\) elementos.
\end{proof}

El ejemplo general de un punto fijo, antes de exigir período seis, es
\[
 (a,b,c,d,e,f,-f,-e,-d,-c,-b,-a).
\]
Esta forma coordinada permite comprobar a mano tanto la involución como el
conteo y evita interpretar el entero \(729\) mediante otro objeto homónimo.

\begin{remark}[Censo exhaustivo de la involución]
\label{rem:procedencia-n40-cm}
La enumeración de las \(3^{12}=531\,441\) palabras no encuentra ninguna
violación de la involución y obtiene
\[
 |\operatorname{Fix}(C_M)|=729,
 \qquad
 |\operatorname{Fix}(C_M)\cap\operatorname{Per}_6|=27.
\]
La demostración primaria y la \cref{prop:estrato-periodico-cm-rev7} explican
ambos números sin depender del recorrido: los puntos fijos tienen seis
coordenadas libres, y la condición adicional de período seis reduce la
elección a tres. El censo exhaustivo confirma que no
existen órbitas excepcionales fuera de esa descripción.
\end{remark}

\begin{criterion}[Falsadores del censo \(C_M\)]
\label{crit:falsadores-cm}
Refutaría el resultado cualquiera de los hechos siguientes: una palabra para
la cual \(C_M^2\mathbf t\ne\mathbf t\); un punto fijo que no posea la forma
coordenada anterior; un total exhaustivo distinto de \(729\); o un total
distinto de \(27\) después de imponer \(r^6\mathbf t=\mathbf t\). La
coincidencia nominal de este \(729\) con otro cardinal no identifica sus
dominios.
\end{criterion}

\section{Estructura cuadrática transportada al estrato fijo}
\label{sec:estructura-cuadratica-estrato-fijo-cm}

Identificamos el alfabeto de signos con el cuerpo ternario mediante
\[
 \{-1,0,+1\}\xrightarrow{\sim}\mathbb F_3,
 \qquad
 -1\longmapsto2,\quad 0\longmapsto0,\quad +1\longmapsto1.
\]
Con esta identificación,
\(\mathscr T_{12}^{\rm word}\cong\mathbb F_3^{12}\), y tanto
\(C_M=-\operatorname{rev}\) como el desplazamiento \(r\) son operadores
\(\mathbb F_3\)-lineales. En consecuencia,
\(\operatorname{Fix}(C_M)\) y
\(\operatorname{Per}_6=\ker(r^6-I)\) son subespacios vectoriales.

La forma coordinada de los puntos fijos define el isomorfismo lineal
\[
 s_{\rm fix}:\mathbb F_3^6\xrightarrow{\sim}\operatorname{Fix}(C_M),
\]
\[
 s_{\rm fix}(a,b,c,d,e,f)
 =(a,b,c,d,e,f,-f,-e,-d,-c,-b,-a).
\]
Sea \(J_W\) el operador de Paley--Witt construido en el
\cref{thm:f9-tres-desde-paley}, con la base, la polarización y la orientación
allí declaradas.

\begin{proposition}[Estructura compleja \(\widetilde J_M^{\,2}=-I\) sobre el estrato fijo]
\label{prop:raiz-menos-uno-estrato-fijo-cm}
El transporte
\[
\widetilde J_M:=s_{\rm fix}J_Ws_{\rm fix}^{-1}
\in\operatorname{End}_{\mathbb F_3}
   \bigl(\operatorname{Fix}(C_M)\bigr)
\]
satisface
\[
 \widetilde J_M^{\,2}=-I.
\]
\end{proposition}

\begin{proof}
Por conjugación,
\[
 \widetilde J_M^{\,2}
 =s_{\rm fix}J_Ws_{\rm fix}^{-1}
  s_{\rm fix}J_Ws_{\rm fix}^{-1}
 =s_{\rm fix}J_W^2s_{\rm fix}^{-1}
 =-I.
\]
\end{proof}

El resultado está demostrado con la carta \(s_{\rm fix}\), la base, la polarización y
la orientación como estructura de partida explícita. No identifica
\(C_M\) con \(J_W\): el primero es una involución del espacio de palabras y
el segundo una estructura cuadrática sobre su estrato fijo. Tampoco demuestra
que el transporte sea natural respecto de toda la dinámica TPK.

Pongamos
\[
 W_{27}
 :=\operatorname{Fix}(C_M)\cap\operatorname{Per}_6.
\]
La inclusión
\[
 W_{27}\hookrightarrow\operatorname{Fix}(C_M)
\]
es canónica, y la \cref{prop:estrato-periodico-cm-rev7} prueba que
\(\dim_{\mathbb F_3}W_{27}=3\).

\begin{corollary}[Obstrucción cuadrática en el estrato de veintisiete]
\label{cor:obstruccion-cuadratica-w27}
El subespacio \(W_{27}\) no puede ser estable bajo
\(\widetilde J_M\).
\end{corollary}

\begin{proof}
Si fuera estable, la relación
\(\widetilde J_M^2=-I\) lo convertiría en un módulo sobre
\[
 \mathbb F_3[X]/(X^2+1)\cong\mathbb F_9.
\]
Todo espacio vectorial sobre \(\mathbb F_9\) tiene dimensión par cuando se
considera sobre \(\mathbb F_3\), en contradicción con
\(\dim_{\mathbb F_3}W_{27}=3\).
\end{proof}

La igualdad cardinal entre \(W_{27}\) y las veintisiete rectas de Schläfli
no proporciona un mapa entre ambos objetos. Tal relación exigiría construir
una aplicación de incidencia que preserve la orientación central y la acción
del estabilizador, además de una clausura dinámica compatible. No se
interpone el atlas de \(81\) celdas entre \(W_{27}\) y el estrato fijo de
\(729\) estados: el atlas de partículas pertenece a otro dominio.

\section[Niveles de Majorana]{Estratificación matemática del objeto Majorana en cuatro niveles}
\label{sec:cuatro-niveles-majorana}

La inversión \(C_M\), la paridad exterior, un fermión relativista de
Majorana y un modo cero de Majorana son objetos de tipos distintos:

\begin{enumerate}
 \item \(C_M=-\operatorname{rev}\) actúa sobre
 \(\mathscr T_{12}^{\rm word}\). Su involutividad y su censo se demuestran
 exactamente; esta construcción no proporciona por sí misma un espinor, una
 ecuación de campo ni un término de masa.

 \item El carácter de signo \(q=-1\) y la completación exterior pertenecen
 al álgebra de intercambio. Relativamente a CAR y a la graduación declarada,
 producen \(N_s^2=N_s\), espectro \(\{0,1\}\) y el retorno de orientación
 \(2\pi/4\pi\). Estas identidades se deducen de CAR y de la graduación; no se
 deducen del
 mero censo de \(C_M\).

 \item Un fermión relativista de Majorana es un objeto espinorial provisto de
 representación relativista, conjugación de carga compatible, condición de
 autorconjugación, término de masa y dinámica. Esos datos no forman parte de
 la definición de \(C_M\) ni se deducen de su censo.

 \item Un modo cero de Majorana, o equivalentemente el sector no abeliano de
 una teoría de Ising en la acepción topológica, requiere simples
 \(\mathbf 1,\psi,\sigma\) con
 \[
  \psi\otimes\psi\simeq\mathbf 1,\qquad
  \psi\otimes\sigma\simeq\sigma,\qquad
  \sigma\otimes\sigma\simeq\mathbf 1\oplus\psi,\qquad
  d_\sigma=\sqrt2,
 \]
 junto con asociador y trenzado coherentes. La brecha, el desdoblamiento y la
 plataforma pertenecen a una realización física adicional; las reglas de
 fusión anteriores no definen una masa de polo universal.
\end{enumerate}

La separación es de dominio y no de nomenclatura: involución de palabras,
paridad CAR, espinor autorconjugado y simple no abeliano son cuatro objetos
matemáticos distintos. En ausencia de un morfismo explícito entre ellos, no
se identifica ninguno por compartir el nombre «Majorana».

\endgroup
\begingroup
\let\chapter\section
\let\section\subsection
\let\subsection\subsubsection
\let\subsubsection\paragraph
\let\clearpage\relax
% Vista científica exacta materializada desde una rama condicional.
% Fuente: source/demostraciones_primarias/31_06c_sectores_topologicos_no_omision.tex; rama: HMTIncludeTopologicalBlock.
% No modifica el contenido matemático de la rama de origen.
\chapter[Incidencia y sectores topológicos]{Incidencia, Fibonacci y sectores topológicos}
\label{chap:incidencia-fibonacci-topologia}
\index{sector topológico}\index{Fibonacci, anillo basado}
\index{grupo de trenzas}

La incidencia residuo--soporte prolonga el atlas hacia estructuras de fusión
y trenza. El orden es esencial: primero se construye el anillo basado
Fibonacci y sus dimensiones; después se estudian las representaciones
algebraicas de trenzas y el centro apuntado. Cada objeto conserva su propio
dominio y no se identifica con una partícula física por compartir nombre o
cardinal.

\section{Incidencia residuo--soporte}
\label{sec:incidencia-residuo-soporte}

Fijamos las clases de residuo y los soportes en el orden
\[
 \begin{aligned}
 C_1&=\{1,4,7\},&
 C_2&=\{2,5,8\},&
 C_0&=\{3,6,9\},\\
 \mathcal D_{\rm act}&=\{1,3,5,7\},&
 \mathcal D_{\rm evt}&=\{2,4,6,8\},&
 \mathcal D_{\bar0}^{(9)}&=\{9\}.
 \end{aligned}
\]
El último símbolo significa «soporte del residuo nonádico nulo». No designa
el dominio completo \(D_9=\{1,\ldots,9\}\) ni un sector físico de masa nula.

Sean \(E_{\mathcal D}\cong\mathbb Z^3\) y
\(E_C\cong\mathbb Z^3\) los módulos libres con bases ordenadas
\[
 (\mathcal D_{\rm act},\mathcal D_{\rm evt},
   \mathcal D_{\bar0}^{(9)})
 \quad\text{y}\quad
 (C_1,C_2,C_0),
\]
respectivamente. La aplicación de incidencia
\(B_{\rm res,sup}:E_{\mathcal D}\to E_C\) cuenta intersecciones.

\begin{theorem}[Forma de Smith e imagen de la incidencia]
\label{thm:incidencia-residuo-soporte}
En las bases precedentes,
\[
 B_{\rm res,sup}
 =\bigl(|C_i\cap\mathcal D_j|\bigr)_{i,j}
 =
 \begin{pmatrix}
 2&1&0\\
 1&2&0\\
 1&1&1
 \end{pmatrix}.
\]
Además,
\[
 \det B_{\rm res,sup}=3,\qquad
 \operatorname{SNF}(B_{\rm res,sup})
 =\operatorname{diag}(1,1,3),
\]
\[
 \operatorname{im}B_{\rm res,sup}
 =\{(u,v,w)\in\mathbb Z^3:u+v\equiv0\pmod3\}.
\]
Si
\[
 P=\begin{pmatrix}0&1&0\\1&0&0\\0&0&1\end{pmatrix},
\]
el intercambio simultáneo
\(C_1\leftrightarrow C_2\) y
\(\mathcal D_{\rm act}\leftrightarrow\mathcal D_{\rm evt}\) es equivariante:
\[
 P B_{\rm res,sup}P=B_{\rm res,sup}
 \qquad\Longleftrightarrow\qquad
 PB_{\rm res,sup}=B_{\rm res,sup}P.
\]
\end{theorem}

\begin{proof}
Las nueve entradas proceden de las intersecciones
\[
\begin{array}{cccc}
 &\mathcal D_{\rm act}&\mathcal D_{\rm evt}
 &\mathcal D_{\bar0}^{(9)}\\
 C_1&\{1,7\}&\{4\}&\varnothing\\
 C_2&\{5\}&\{2,8\}&\varnothing\\
 C_0&\{3\}&\{6\}&\{9\}.
\end{array}
\]
La expansión por la tercera columna da
\[
 \det B_{\rm res,sup}
 =\det\begin{pmatrix}2&1\\1&2\end{pmatrix}=3.
\]
El máximo común divisor de las entradas es uno. El menor de orden dos
formado por las filas primera y tercera y las columnas segunda y tercera
tiene determinante uno:
\[
 \det\begin{pmatrix}1&0\\1&1\end{pmatrix}=1.
\]
Por los divisores de Smith, \(d_1=d_2=1\). El producto
\(d_1d_2d_3\) coincide con el valor absoluto del determinante, que el
cálculo anterior fija en tres. Por consiguiente,

\[
 d_1d_2d_3=3.
\]
La forma normal es, por tanto, \(\operatorname{diag}(1,1,3)\).

Cada columna satisface \(u+v\equiv0\pmod3\), así que la imagen está contenida
en el submódulo indicado. Recíprocamente, si \(u+v\equiv0\pmod3\), entonces
\[
 x=\frac{2u-v}{3},\qquad
 y=\frac{2v-u}{3},\qquad
 z=w-x-y
\]
son enteros y
\[
 B_{\rm res,sup}(x,y,z)^{\mathsf T}=(u,v,w)^{\mathsf T}.
\]
Esto prueba la igualdad de imágenes. Finalmente, la multiplicación explícita
por \(P\) intercambia las dos primeras filas y las dos primeras columnas y
deja la matriz invariante.
\end{proof}

\begin{criterion}[Control de la incidencia]
\label{crit:incidencia-residuo-soporte}
Falsaría el teorema una sola intersección mal contada, un menor que alterase
los divisores de Smith, un vector de la imagen con \(u+v\not\equiv0\pmod3\),
un vector que satisficiera la congruencia pero no admitiese la preimagen
exhibida, o la pérdida de la equivarianza \(PBP=B\).
\end{criterion}

\section[Anillo basado de Fibonacci]{Anillo basado asociado a la incidencia
  \texorpdfstring{\(F_{\rm av}\)}{F\_av}}
\label{sec:anillo-basado-fibonacci}

El \cref{thm:descenso-necesario-fav} demuestra, a partir de la órbita de
modos \((\Sigma,\Pi,\Pi)\), la matriz
\[
 F_{\rm av}=
 \begin{pmatrix}0&1\\1&1\end{pmatrix}.
\]
De esta matriz se obtiene el siguiente anillo basado.

\begin{definition}[Anillo basado de Fibonacci]
\label{def:anillo-basado-fibonacci}
Sea
\[
 K_{\rm Fib}:=\mathbb Z\mathbf 1\oplus\mathbb Z\tau
\]
el grupo abeliano libre con unidad \(\mathbf 1\) y producto bilineal fijado por
\[
 \tau\cdot\mathbf 1=\tau,\qquad
 \tau\cdot\tau=\mathbf 1+\tau.
\]
En este estrato, \(+\) y \(\cdot\) son la adición y la multiplicación del
anillo. Los símbolos categóricos \(\oplus\) y \(\otimes\) se reservan para
una categorificación elegida posteriormente.
\end{definition}

\begin{theorem}[Realización de \(F_{\rm av}\) y dimensión de Perron]
\label{thm:anillo-fibonacci-fav}
En la base ordenada \((\mathbf 1,\tau)\), la multiplicación por \(\tau\) en
\(K_{\rm Fib}\) tiene matriz \(F_{\rm av}\). Si
\(F_0=0\), \(F_1=1\) y
\(F_{n+1}=F_n+F_{n-1}\) para \(n\geq1\), entonces
\[
 \tau^n=F_{n-1}\mathbf 1+F_n\tau
 \qquad(n\geq1).
\]
El único homomorfismo positivo de dimensión que envía
\(\mathbf 1\) a uno satisface
\[
 \operatorname{FPdim}(\tau)=\varphi.
\]
La identificación de las dos hojas tipadas con la base
\((\mathbf 1,\tau)\) determina por completo este anillo y su dimensión
positiva.
\end{theorem}

\begin{proof}
Las columnas de la matriz de multiplicación por \(\tau\) son las coordenadas
de
\[
 \tau\cdot\mathbf 1=\tau
 \quad\text{y}\quad
 \tau\cdot\tau=\mathbf 1+\tau;
\]
por tanto son \((0,1)^{\mathsf T}\) y \((1,1)^{\mathsf T}\), respectivamente.
Se obtiene así \(F_{\rm av}\).

La fórmula de potencias vale para \(n=1\). Si vale para \(n\), entonces
\[
 \tau^{n+1}
 =F_{n-1}\tau+F_n(\mathbf 1+\tau)
 =F_n\mathbf 1+F_{n+1}\tau,
\]
lo que cierra la inducción. Finalmente, una dimensión multiplicativa positiva
\(d=\operatorname{FPdim}(\tau)\) obedece
\[
 d^2=1+d.
\]
Sus raíces son \(\varphi\) y \(-\varphi^{-1}\); la positividad selecciona
\(d=\varphi\).
\end{proof}

La denominación \emph{Fibonacci} queda así explicada por la propia regla de
fusión: al iterar \(\tau\), los dos canales posibles aparecen con
multiplicidades consecutivas \(F_{n-1}\) y \(F_n\). El espacio de estados de
fusión crece, por tanto, con razón asintótica \(\varphi\). Ésta es la base
combinatoria de su uso computacional: la información se codifica en canales
globales de fusión y no en una masa asignada a un punto.

\begin{criterion}[Condiciones de una categorización Fibonacci]
\label{crit:categorificacion-fibonacci}
La capa anular quedaría refutada si la matriz de multiplicación no fuese
\(F_{\rm av}\), si fallase la recurrencia de potencias o si la dimensión
positiva no fuese \(\varphi\). Una categorización trenzada cuyo anillo de
Grothendieck sea el de la \cref{def:anillo-basado-fibonacci} exige además un
funtor monoidal
\[
 \mathfrak C_{\rm Fib}:\mathsf{TPK}
 \dashrightarrow\mathcal C_{\rm Fib},
\]
la verificación de pentágono y hexágono en su imagen y una regla interna que
seleccione calibre y quiralidad.
Estas condiciones no son consecuencias de la identidad anular por sí sola.
\end{criterion}

\section[Caracteres de trenza]{Seis caracteres y cuatro torsiones de trenza}
\label{sec:rescate-trenzas-ex02}

Sea \(B_n\), \(n\geq2\), el grupo de trenzas de Artin con generadores
\(\sigma_1,\ldots,\sigma_{n-1}\), relaciones
\[
 \sigma_i\sigma_{i+1}\sigma_i
 =\sigma_{i+1}\sigma_i\sigma_{i+1},
 \qquad
 \sigma_i\sigma_j=\sigma_j\sigma_i\quad(|i-j|\geq2),
\]
y sea \(w:B_n\to\mathbb Z\) la suma de exponentes. Los lectores
\[
 w_2:=w\bmod2,\qquad w_3:=w\bmod3
\]
se reúnen, por el teorema chino del resto, en
\[
 w_6:B_n\longrightarrow
 \mathbb Z/2\mathbb Z\times\mathbb Z/3\mathbb Z
 \simeq\mathbb Z/6\mathbb Z.
\]

\begin{theorem}[Abelianización, seis caracteres y torsiones de permutación]
\label{thm:trenzas-caracteres-z6}
Para \(n\geq2\),
\[
 B_n^{\rm ab}\simeq\mathbb Z,
\]
mediante \(w\). En consecuencia, \(w_6\) produce seis caracteres unitarios
\[
 \chi_k(b)
 :=\exp\!\left(\frac{2\pi\mathrm i k}{6}w(b)\right),
 \qquad k\in\mathbb Z/6\mathbb Z.
\]

Sea \(V\ne0\) un espacio de Hilbert, sea
\(P_i:V^{\otimes n}\to V^{\otimes n}\) el intercambio de los factores
\(i,i+1\), y sea \(q\in\mu_6\). La regla
\[
 \rho_q(\sigma_i):=qP_i,
 \qquad
 \rho_q(b)=q^{w(b)}\pi_n(b),
\]
donde \(\pi_n\) es la representación de permutación, define una
representación de \(B_n\). Esta representación desciende al grupo simétrico
si y sólo si
\[
 q^2=1.
\]
Para los cuatro valores \(q\in\mu_6\setminus\{\pm1\}\),
\[
 \operatorname{Inv}(\rho_q)=0.
\]
Esos cuatro valores definen torsiones escalares genuinamente trenzadas, pero
no constituyen por sí solos cuatro espacios de Fock ni cuatro sectores
anyónicos físicos.
\end{theorem}

\begin{proof}
En todo cociente abeliano, la relación de trenza da
\[
 2[\sigma_i]+[\sigma_{i+1}]
 =[\sigma_i]+2[\sigma_{i+1}],
\]
luego todos los generadores tienen la misma clase. No queda otra relación
sobre esa clase, pues \(w(\sigma_i)=1\); por tanto la abelianización es
\(\mathbb Z\).

Los operadores \(P_i\) satisfacen las relaciones de trenza y
\(P_i^2=I\). Multiplicar cada generador por el mismo escalar \(q\) conserva
las dos relaciones, de modo que \(\rho_q\) es una representación. Para
factorizar a través de \(S_n\) debe satisfacer la relación adicional
\(\sigma_i^2=1\), y
\[
 \rho_q(\sigma_i)^2=q^2I.
\]
Esto prueba el criterio \(q^2=1\).

Si \(v\) fuera invariante, entonces
\(qP_i v=v\), es decir, \(P_i v=q^{-1}v\), para todo \(i\). Como cada
involución \(P_i\) sólo tiene autovalores \(\pm1\), no existe tal vector no
nulo cuando \(q\ne\pm1\). Por consiguiente,
\(\operatorname{Inv}(\rho_q)=0\) para los cuatro valores restantes.
\end{proof}

Cuando \(\dim V>1\) y \(n\geq3\), los operadores \(P_1\) y \(P_2\) no
conmutan; por ello la representación completa no se vuelve abeliana por
proceder de un carácter escalar. El carácter \(\chi_k\) y la representación
\(\rho_q\) son objetos distintos: el primero es unidimensional; la segunda
retiene la acción de permutación.

\section[Centro modular y realizaciones topológicas]{Centro modular apuntado y condiciones de realización Fibonacci e Ising}
\label{sec:centro-apuntado-no-go}

Sea
\[
 A=(\mathbb Z/9\mathbb Z)^2
\]
y sea \(\operatorname{Vec}_A\) la categoría de espacios vectoriales
\(A\)-graduados con asociador trivial. Fijar un bicharácter no degenerado
identifica \(\widehat A\) con \(A\), pero no altera los conteos siguientes.

\begin{theorem}[Centro apuntado y obstrucción dimensional]
\label{thm:centro-apuntado-no-go}
El centro de Drinfeld
\[
 \mathcal Z(\operatorname{Vec}_A)
\]
es una categoría modular apuntada cuyos simples se parametrizan por
\((a,\chi)\in A\times\widehat A\). Por tanto posee
\[
 |A|\,|\widehat A|=81^2=6561
\]
simples, todos de dimensión cuántica uno, y dimensión cuántica total
\[
 \mathcal D
 =\sqrt{\sum_{x\ {\rm simple}}d_x^2}
 =\sqrt{6561}=81.
\]
En esta categoría no puede realizarse, preservando simples y dimensiones de
Frobenius--Perron, ni el anillo Fibonacci ni el sector Ising: éstos requieren
respectivamente un simple de dimensión \(\varphi\) y un simple de dimensión
\(\sqrt2\).
\end{theorem}

\begin{proof}
Para un grupo abeliano finito y asociador trivial, cada objeto simple del
centro queda determinado por un grado \(a\in A\) y un carácter
\(\chi\in\widehat A\) que especifica su media trenza. Tanto el objeto
graduado como el carácter son unidimensionales; en consecuencia, todos los
simples son invertibles y tienen dimensión cuántica uno. Como
\(|A|=|\widehat A|=9^2=81\), el número de pares es \(81^2\), y la fórmula de
\(\mathcal D\) da \(81\).

Toda subcategoría de fusión de una categoría apuntada vuelve a ser apuntada;
en particular, sus simples tienen dimensión uno. Esto contradice
\(\operatorname{FPdim}(\tau)=\varphi\) para Fibonacci y
\(d_\sigma=\sqrt2\) para Ising. La obstrucción es dimensional y no depende
de una semejanza de nombres o cardinales.
\end{proof}

\begin{criterion}[Condiciones para una realización física de los sectores]
\label{crit:promocion-sectores-topologicos}
Los resultados de
\cref{thm:trenzas-caracteres-z6,thm:centro-apuntado-no-go} quedarían
algebraicamente falsados por un fallo de las relaciones de Artin, del
criterio \(q^2=1\), del censo \(6561\), de la apuntabilidad o de
\(\mathcal D=81\). Una realización física exige, además, un
Hamiltoniano o protocolo impulsado, una brecha espectral, operadores de
trenza realizables, protección frente a perturbaciones y un mapa desde las
multisecciones enriquecidas. Ninguno de esos datos se deduce únicamente de
los dos teoremas algebraicos anteriores.
\end{criterion}

\endgroup
\begingroup
\let\chapter\section
\let\section\subsection
\let\subsection\subsubsection
\let\subsubsection\paragraph
\let\clearpage\relax
% Las formulaciones iniciales de esta fuente coinciden byte por byte con
% las residencias REV.2 de realizaciones y ontología. Se publica aquí únicamente
% su desarrollo exclusivo posterior, sin alterar ninguna fuente de procedencia.
% Extracto editor-ready del contenido exclusivo de masas/08_composicion.tex.
% Las líneas 1--52 son byte-exactas a 75_rev2_ontologia_particula.tex:231--282
% y las líneas 53--396 son byte-exactas a 74_rev2_realizaciones_i.tex:641--1052
% en dos tramos contiguos. Sus propietarios íntegros permanecen preservados;
% esta residencia publica una sola vez el contenido exacto y conserva desde
% la línea 397 todo el desarrollo exclusivo del propietario de composición.

\subsection{Objeto compuesto como sección persistente de un producto cocíclico de rutas}
Una partícula compuesta es una sección del producto fibrado de biografías constituyentes con una frontera de ligadura. Los registros genealógicos se componen antes de aplicar el extractor de firma. Si \(\Gamma\) y \(\Lambda\) se enlazan mediante la frontera \(B\) y la acción \(g\),

\[
(n,e)_\Gamma\star_{(g,B)}(n,e)_\Lambda
=
(n_\Gamma+n_\Lambda-n_B,
e_\Gamma+ge_\Lambda-e_B).
\]
La firma compuesta es

\[
\operatorname{Sig}_{\rm comp}
=L_{\rm tick}\circ\mathfrak C_{12}.
\]
No coincide en general con la suma de las firmas individuales, porque el registro genealógico de enlace porta energía, acarreo y memoria. Esta es la residencia de la energía de ligadura.

\Needspace{7\baselineskip}
\subsection{Resonancia como sección transitoria y anchura como tasa de desintegración}
Una resonancia es una sección metastable cuyo resolvente posee un polo de energía complejo. Si \(z_X=M_Xc_{\rm int}^{2}-i\Gamma_X/2\), la masa es \(M_X=c_{\rm int}^{-2}\operatorname{Re}z_X\); la anchura mide su vida media:

\[
\Gamma_X=\frac{\hbar_{\rm int}}{\tau_X^{\rm life}}.
\]
La anchura no es una sexta coordenada de \(\sigma\). Procede del operador de supervivencia, de los canales de decaimiento y de la continuación analítica del resolvente. El inventario de 384 anchuras exige por ello un lector propio, separado del inventario de 471 masas.

\Needspace{7\baselineskip}
\subsection{Sección virtual como extensión intermedia no asintótica del atlas}
Una sección virtual es una extensión finita que participa en una composición o propagador pero no define una rama global estable en el codominio físico. Si \(L(s)\) es su horizonte,

\[
L(s)<\infty,
\qquad
\operatorname{Ext}_{L(s)}(s)\ne\varnothing,
\qquad
\operatorname{Ext}_{L(s)+1}(s)=\varnothing.
\]
La virtualidad es una obstrucción terminal, no una etiqueta documental ni una masa imaginaria. Una resonancia puede ser inestable sin ser virtual en este sentido; una línea interna puede ser virtual sin constituir una nueva especie.

\Needspace{7\baselineskip}
\subsection{Sección de Majorana como punto fijo de la conjugación partícula–antipartícula}
Cuatro construcciones deben mantenerse separadas. La involución de palabras \(C_M\) define la antisección; un espinor de Majorana satisface una condición de realidad; un modo cero Majorana es una excitación localizada de energía cero en un Hamiltoniano topológico; el sector de Ising posee reglas de fusión

\[
\sigma\otimes\sigma=1\oplus\psi.
\]
Una partícula física es Majorana cuando su representación fermiónica admite una autorconjugación compatible con el Hamiltoniano, la localización y el producto interno. La neutralidad eléctrica y la igualdad de una palabra bajo \(C_M\) no bastan. En particular, decidir la naturaleza Majorana de los neutrinos requiere un entrelazador entre el operador neutro, la base de masas y la estructura real fermiónica.

\Needspace{7\baselineskip}
\subsection{Sección anyónica y representación no trivial del grupo de trenzas}
Un anyón es una sección bidimensional cuya representación del grupo de trenzas no se reduce a la permutación bosónica o fermiónica. En el sector Fibonacci,

\[
\tau\otimes\tau=1\oplus\tau,
\qquad
\operatorname{FPdim}(\tau)=\varphi.
\]
En el sector abeliano \(\mathbb Z/6\mathbb Z\), el cociente de devanado posee seis caracteres unidimensionales. Para obtener una partícula física se necesitan matrices \(F\) y \(R\), pentágono, hexágono, Hamiltoniano con brecha y operadores de trenzado. El anillo de fusión exacto determina el codominio topológico, pero no fija por sí solo una masa universal.

\Needspace{7\baselineskip}
\subsection{Clases taquiónica y fantasma como obstrucciones de positividad y estabilidad}
Una partícula persistente posee evolución unitaria, energía acotada inferiormente y masa espectral real en la recta física. Un taquión exigiría un polo con \(m^2<0\) que permaneciera como estado asintótico estable. En la arquitectura HMT--MD, una dirección que abandona el cono de realizaciones positivas o bien retorna a la frontera nula, o bien pierde supervivencia, o bien se manifiesta como inestabilidad de un fondo que debe condensar. No define una sección persistente superlumínica. La imposibilidad se formula como incompatibilidad entre persistencia global, positividad espectral y \(m^2<0\), no como prohibición de parámetros taquiónicos transitorios en una expansión perturbativa.

Un fantasma es un modo de norma negativa. La positividad del producto interno y la conservación unitaria excluyen su realización física. Los fantasmas de Faddeev--Popov pertenecen al complejo de calibre y no al espacio físico de estados; esta distinción debe conservarse.

\Needspace{7\baselineskip}
\section{Composición, resonancias y anchuras}
\Needspace{7\baselineskip}
\subsection{Composición no aditiva de masas mediante la ley cocíclica del enlace}
La Ley General de Masas actúa sobre biografías genealógicas. Cuando varios constituyentes forman un estado ligado, el objeto de entrada no es la lista de sus masas aisladas, sino un registro compuesto que conserva el orden de acoplamiento, la frontera común, el acarreo, la orientación y la memoria. Sean

\[
\Gamma_1,\ldots,\Gamma_n
\]
rutas constituyentes, \(\mathcal T\) un árbol temporal de composición y \(\partial\) la frontera de enlace. El operador dodecafásico de composición se escribe

\[
\mathfrak C_{12}^{(n)}:
(\Gamma_1,\ldots,\Gamma_n;\mathcal T,\partial)
\longmapsto \mathcal L_{\rm comp}.
\]
Sólo después se aplican el lector del registro, el extractor de firma, el carácter y las torres:

\[
\mathcal L_{\rm comp}
\longmapsto \sigma_{\rm comp}
\longmapsto \chi_{\rm comp}
\longmapsto \widehat M_{\rm comp}.
\]
Esta precedencia es esencial. Si se evaluasen primero las masas aisladas y se sumasen después, desaparecerían precisamente los datos que distinguen un estado ligado de un umbral de constituyentes libres. La energía de ligadura no es un término agregado al final: es la imagen dimensional de una diferencia de registro producida por el enlace.

\Needspace{7\baselineskip}
\subsection{Ley cocíclica del enlace}
Para dos registros \(L_1,L_2\), la composición tiene la forma

\[
L_1\star L_2
=L_1+L_2+\mathfrak b(L_1,L_2),
\]
donde \(\mathfrak b\) registra frontera, acarreo y memoria compartidos. La independencia respecto de la parentización exige la identidad cocíclica

\[
\mathfrak b(L_1,L_2)
+\mathfrak b(L_1\star L_2,L_3)
=
\mathfrak b(L_2,L_3)
+\mathfrak b(L_1,L_2\star L_3).
\]
Si esta igualdad falla, las dos maneras de construir un mismo árbol triple producen registros genealógicos distintos y no existe una especie compuesta bien definida. Si se cumple, el cociclo no desaparece: garantiza que su contribución es una propiedad del estado y no un artefacto de la notación.

Una presentación más explícita conserva además la acción de borde. Sea \(e\in\mathbb Z^{12}\) el vector ordenado de acontecimientos y sea \(n\) el contador de acción o de acontecimientos de un registro. Si \((n,e)_\Gamma\) y \((n,e)_\Lambda\) son registros, si \(B\) es el subregistro común con datos \((n_B,e_B)\), y si \(g\in\operatorname{Aut}(\mathbb Z^{12})\) transporta la orientación del segundo registro a la del primero preservando el lector,

\[
(n,e)_\Gamma\star_{(g,B)}(n,e)_\Lambda
=
(n_\Gamma+n_\Lambda-n_B,
e_\Gamma+g e_\Lambda-e_B).
\]
El término \(B\) evita contar dos veces la frontera identificada; la acción \(g\) transporta el segundo registro a la orientación del primero. La firma compuesta es, por tanto,

\[
\operatorname{Sig}_{\rm comp}
=L_{\rm tick}\!\left(
\mathfrak C_{12}^{(n)}
(\Gamma_1,\ldots,\Gamma_n;\mathcal T,\partial)
\right),
\]
y no la suma automática de firmas elementales.

\Needspace{7\baselineskip}
\subsection{Clausura cocíclica de la composición mesónica quark–antiquark}
Una expresión mesónica orientada tiene la forma

\[
M(q,q')=q\otimes(q')^\vee.
\]
El dominio bruto contiene \(6\times6\times3\times4=432\) expresiones al considerar sabor, color compatible y orientación. Esta cardinalidad no es el número de mesones físicos. Antes de individualizar una especie se aplican:

\begin{enumerate}
\item el proyector de singlete de color;
\item el proyector de carga e isospín;
\item la representación de espín y paridad;
\item el árbol temporal de enlace;
\item la condición radial u orbital;
\item el canal de preparación y decaimiento cuando el estado es resonante.
\end{enumerate}
El descriptor físico completo es

\[
\mathcal I_M=
(q\bar q',J^{PC},I,I_3,S,C,B,\mathcal T,\partial,
\Gamma_M,\mathcal L_M,\sigma_M,\eta_M).
\]
Dos estados con los mismos sabores pueden diferir en \(J^{PC}\), excitación radial, momento orbital, hoja o canal, y por ello poseer polos distintos. La masa no se deriva del mero nombre de los constituyentes.

Las evaluaciones publicadas para \(\pi^0,\pi^\pm,K^\pm,K^0\) son exactas una vez fijados los registros históricos declarados. La afirmación más fuerte ---la generación prospectiva de cada registro nominal desde su árbol de enlace--- requiere que \(\mathfrak C_{12}\) reconstruya esos registros sin consultar el valor que será comparado. Ambas capas se conservan juntas: el cálculo condicionado no se borra, y tampoco se presenta como selector de identidad.

\Needspace{7\baselineskip}
\subsection{Clausura trilineal de la composición bariónica y neutralización del carácter de color}
Una expresión bariónica orientada usa tres constituyentes y tres colores. El dominio bruto del censo adoptado contiene

\[
\lvert\mathcal B_{\rm RGB}^{\rm raw}\rvert=1728
\]
expresiones. La proyección física debe imponer de forma simultánea:

\[
P_{\rm singlete\ de\ color},
\qquad
P_{J^P},
\qquad
P_{I,I_3,S,C,B},
\qquad
P_{\rm Pauli}.
\]
El árbol triple conserva qué par se acopló primero y cómo atravesó la frontera el tercer constituyente. La identidad cocíclica garantiza que las parentizaciones equivalentes representan el mismo estado cuando el entrelazador físico las identifica; los canales inequivalentes permanecen como estados distintos.

Las firmas históricas

\[
\sigma_p=(59,0,9,1,0),
\qquad
\sigma_n=(59,1,-34,0,1)
\]
producen las evaluaciones publicadas de protón y neutrón al aplicar el carácter completo. Son controles exactos del lector para esos registros. La explicación física de la diferencia neutrón--protón exige además derivar, antes del contraste, el signo del isospín, la orientación de vacancias, el borde de enlace y la contribución electromagnética en una carta común.

\Needspace{7\baselineskip}
\subsection{Realizaciones compuestas exóticas y condiciones de estabilidad del enlace}
La misma gramática admite más constituyentes sin modificar la ley:

\[
\mathfrak C_{12}^{(n)}:
\prod_{j=1}^{n}\Gamma_j\times\mathcal T_n\times\partial_n
\longrightarrow\mathcal L_{\rm comp}.
\]
Un tetraquark corresponde a \(qq\bar q\bar q\); un pentaquark, a \(qqqq\bar q\); un glueball, a rutas de la representación adjunta de color proyectadas al singlete; un híbrido, a rutas fermiónicas y de calibre reunidas en un mismo árbol. La lista de constituyentes no determina por sí sola la especie. Una clasificación física completa debe incluir color, espín, paridad, conjugación, isospín, topología del árbol y polo del operador acoplado.

Esta formulación discrimina también dos descripciones rivales de un mismo candidato ---molécula hadrónica y estado compacto, por ejemplo---: son dos árboles o dos proyectores sobre un mismo sector de constituyentes. Se vuelven equivalentes únicamente si existe un unitario que entrelace sus operadores, sus canales y sus polos.

\Needspace{7\baselineskip}
\subsection{Estado estable, estado ligado y resonancia}
Conviene separar tres nociones.

Un estado estable es un autovector normalizable del Hamiltoniano físico con energía real y sin canal de decaimiento abierto. Un estado ligado es un autovector por debajo del umbral pertinente; puede ser estable o decaer por otro canal. Una resonancia es un polo de la continuación del resolvente o de la matriz de dispersión. No es un autovector normalizable del Hamiltoniano autoadjunto sobre el eje real.

Para un polo de energía

\[
z_X=E_X-\frac{i}{2}\Gamma_X
=M_Xc_{\rm int}^{2}-\frac{i}{2}\Gamma_X,
\qquad \Gamma_X>0,
\]
la energía de polo es \(E_X=\operatorname{Re}z_X\), la masa de polo es \(M_X=c_{\rm int}^{-2}\operatorname{Re}z_X\) y la anchura energética es \(\Gamma_X=-2\operatorname{Im}z_X\). La Ley General de Masas suministra la rama real condicionada por la identidad y el registro; la continuación abierta de canales determina el desplazamiento complejo y la anchura. En una carta de unidades naturales \(c_{\rm int}=1\), las coordenadas numéricas de energía y masa coinciden, pero sus tipos no se identifican antes de elegir esa carta.

\Needspace{7\baselineskip}
\subsection{Corrección del signo en la lectura discreta del polo}
Sea \(t_0>0\) un paso temporal y supóngase que el modo de un paso tiene autovalor

\[
\lambda=r e^{-i\theta},
\qquad 0<r<1,
\]
con la convención

\[
\lambda=\exp\!\left(-\frac{i}{\hbar_{\rm int}}z t_0\right).
\]
Entonces

\[
z=\frac{\hbar_{\rm int}}{t_0}
\bigl(\theta+i\log r\bigr)
=E-\frac{i}{2}\Gamma
=Mc_{\rm int}^{2}-\frac{i}{2}\Gamma,
\]
de modo que

\[
E=\frac{\hbar_{\rm int}}{t_0}\theta,
\qquad
M=\frac{\hbar_{\rm int}}{t_0c_{\rm int}^{2}}\theta,
\qquad
\Gamma=-\frac{2\hbar_{\rm int}}{t_0}\log r>0.
\]
El signo positivo delante de \(i\log r\) es obligado: como \(\log r<0\), la parte imaginaria de \(z\) es negativa. La forma \(\theta-i\log r\) colocaría el polo en el semiplano superior y describiría crecimiento, no decaimiento.

La probabilidad de supervivencia tras \(N\) pasos es

\[
|\lambda|^{2N}=r^{2N}
=\exp\!\left(-\frac{\Gamma}{\hbar_{\rm int}}Nt_0\right).
\]
Por ello el tiempo de vida asociado es

\[
\tau_X^{\rm life}=\frac{\hbar_{\rm int}}{\Gamma_X}.
\]
Esta identidad demuestra además por qué la anchura no debe introducirse como sexta coordenada de la firma \(\mathbb Z^5\): nace del módulo del modo abierto, no del carácter positivo que organiza la masa.

\Needspace{7\baselineskip}
\subsection{Operador efectivo de canales}
Sea \(P\) el subespacio preparado y \(Q=I-P\) el conjunto de canales eliminados. El complemento se incorpora mediante el operador de Feshbach

\[
H_{\rm eff}(z)
=PHP+PHQ(z-QHQ)^{-1}QHP.
\]
Los polos satisfacen

\[
\det\bigl(z-H_{\rm eff}(z)\bigr)=0.
\]
El primer término conserva el operador interno preparado. El segundo es la autoenergía de los canales: desplaza la parte real y produce una parte imaginaria cuando se abre el continuo. La masa de carácter y la anchura de supervivencia confluyen en el mismo polo, pero proceden de mapas distintos.

Una formulación equivalente usa un proyector de supervivencia distinto del complemento de Feshbach. Sea

\[
S_X=S_X^*=S_X^2
\]
el proyector sobre el subespacio superviviente del estado preparado. El operador comprimido de un paso es

\[
T_X=S_XUS_X.
\]
Supóngase que \(T_X\) posee un autovalor dominante simple \(\lambda_X\), que el estado preparado tiene solapamiento no nulo con su subespacio espectral y que el resto del espectro está separado en módulo. Si \(r_X=|\lambda_X|=\rho(T_X)<1\), entonces

\[
\Gamma_X=-\frac{2\hbar_{\rm int}}{t_0}\log r_X.
\]
Si el operador se normaliza desde el principio sobre probabilidades en vez de amplitudes, desaparece el factor dos. La publicación debe declarar cuál de las dos convenciones usa; mezclarlas produce una anchura errónea por un factor dos.

\Needspace{7\baselineskip}
\subsection{Conjugación de partícula y antipartícula}
Sea \(C_M\) la involución que invierte orientación y hoja, y sea \(U_C\) su realización unitaria sobre el espacio físico. La igualdad de masas de polo y anchuras requiere covariancia tanto del operador como de los canales:

\[
U_CH_{\rm eff}^{X}(z)U_C^{-1}
=H_{\rm eff}^{\bar X}(z).
\]
Entonces los conjuntos de polos retardados coinciden después de transportar el sector y la prescripción causal, y

\[
M_{\bar X}=M_X,
\qquad
\Gamma_{\bar X}=\Gamma_X.
\]
La equivalencia unitaria desnuda no convierte el polo retardado en su conjugado complejo; transporta el mismo problema espectral al sector antipartícula. La paridad de la firma másica no basta para demostrar la segunda igualdad: también los canales y sus medidas deben ser covariantes. Si la simetría se representa antiunitariamente, debe transportarse además la prescripción retardada; la conjugación desnuda \(z\mapsto\bar z\) intercambia polos retardados y avanzados y no es, por sí sola, la afirmación física de igualdad de vidas medias.

\Needspace{7\baselineskip}
\subsection{Inventario tipado de realizaciones y observables de contraste}
El corte externo preservado contiene

\[
471\ \text{observables de masa},
\qquad
384\ \text{observables de anchura}.
\]
La partición de masas es

\Needspace{7\baselineskip}
\begingroup
\footnotesize
\setlength{\tabcolsep}{1.9pt}
\begin{longtable}{L{5.04cm}L{5.04cm}L{5.04cm}}
\toprule
\rowcolor{HMTNavy}\HMTtablehead{Dominio físico de llegada} & \HMTtablehead{Masas} & \HMTtablehead{Anchuras} \\
\midrule
\endfirsthead
\multicolumn{3}{l}{\sffamily\scriptsize\color{HMTGray}Continuación de la tabla}\\[-1pt]
\toprule
\rowcolor{HMTNavy}\HMTtablehead{Dominio físico de llegada} & \HMTtablehead{Masas} & \HMTtablehead{Anchuras} \\
\midrule
\endhead
\midrule
\endfoot
\bottomrule
\endlastfoot
partículas elementales, leptónicas y extensiones & 53 & 9 \\
\rowcolor{HMTLight}quarks y extensiones quark & 8 & 1 \\
mesones & 243 & 225 \\
\rowcolor{HMTLight}bariones & 166 & 149 \\
registro gravitatorio de contraste & 1 & 0 \\
\rowcolor{HMTLight}\textbf{Total} & \textbf{471} & \textbf{384} \\
\end{longtable}
\endgroup
Estas filas son observables, no especies elementales. Varias filas pueden referirse a diferentes cargas de un mismo multiplete, estimadores de una misma resonancia, masas Breit--Wigner, masas de polo o estados excitados. La individualización interna precede siempre a la fila externa.

El inventario conserva además una propiedad metodológica fuerte: ninguna de sus 855 magnitudes fue autorizada como selector de ruta. Los 471 registros de masa y los 384 de anchura sólo se abren después de congelar la identidad HMT, el proyector, el operador y la carta de comparación.

\iffalse
\Needspace{7\baselineskip}
\subsection{Malla histórica extendida de pares y resonancias}
Una formulación histórica asoció a la malla angular extendida la expresión

\[
\Theta(n,k,t)
=nA_{\rm rad}-k\frac{C^*_{\rm rad}}6+t\Delta_4.
\]
Con las definiciones vigentes de \(A\), \(C^*\), \(\Delta_4\) y de la carta exponencial, esta expresión aislada no reproduce las cifras archivadas. Se conserva como hipótesis histórica de parametrización, no como su generador ejecutable. Conservamos las doce evaluaciones importadas porque forman parte de la genealogía numérica de la ley. Cinco nombres ---\(\eta,\rho^0,\phi(1020),J/\psi\) y \(\Upsilon(1S)\)--- amplían materialmente las etiquetas visibles respecto de las diecisiete evaluaciones rectoras:

\Needspace{7\baselineskip}
\begingroup
\fontsize{7.2}{8.2}\selectfont
\setlength{\tabcolsep}{1.9pt}
\begin{longtable}{L{1.25cm}L{0.78cm}L{0.78cm}L{0.78cm}L{3.43cm}L{3.90cm}L{4.06cm}}
\toprule
\rowcolor{HMTNavy}\HMTtablehead{Estado} & \HMTtablehead{\(n\)} & \HMTtablehead{\(k\)} & \HMTtablehead{\(t\)} & \HMTtablehead{Referencia posterior (MeV/\allowbreak{}c\({}^{2}\))} & \HMTtablehead{Evaluación HMT histórica (MeV/\allowbreak{}c\({}^{2}\))} & \HMTtablehead{Diferencia (ppm)} \\
\midrule
\endfirsthead
\multicolumn{7}{l}{\sffamily\scriptsize\color{HMTGray}Continuación de la tabla}\\[-1pt]
\toprule
\rowcolor{HMTNavy}\HMTtablehead{Estado} & \HMTtablehead{\(n\)} & \HMTtablehead{\(k\)} & \HMTtablehead{\(t\)} & \HMTtablehead{Referencia posterior (MeV/\allowbreak{}c\({}^{2}\))} & \HMTtablehead{Evaluación HMT histórica (MeV/\allowbreak{}c\({}^{2}\))} & \HMTtablehead{Diferencia (ppm)} \\
\midrule
\endhead
\midrule
\endfoot
\bottomrule
\endlastfoot
\(\mu\) & 64 & 146 & 1 & 105.658374 & 105.656529 & -17.461938\allowbreak{}2274 \\
\rowcolor{HMTLight}\(\tau\) & 64 & 349 & -1 & 1776.86 & 1776.835 & -14.069763\allowbreak{}5154 \\
\(\pi^\pm\) & 37 & 125 & 0 & 139.57039 & 139.570582 & 1.375649\allowbreak{}94982 \\
\rowcolor{HMTLight}\(K^\pm\) & 52 & 177 & 1 & 493.677 & 493.679 & 4.051231\allowbreak{}87833 \\
\(K^0\) & 53 & 178 & 1 & 497.611 & 497.616 & 10.048009\allowbreak{}3889 \\
\rowcolor{HMTLight}\(p\) & 55 & 195 & 1 & 938.272081 & 938.27244 & 0.382618\allowbreak{}226919 \\
\(n\) & 55 & 197 & 0 & 939.565413 & 939.538 & -29.176254\allowbreak{}9161 \\
\rowcolor{HMTLight}\(\eta\) & 46 & 158 & 0 & 547.862 & 547.836 & -47.457206\allowbreak{}3768 \\
\(\rho^0\) & 50 & 176 & 0 & 775.26 & 775.267 & 9.029228\allowbreak{}90385 \\
\rowcolor{HMTLight}\(\phi(1020)\) & 54 & 190 & 0 & 1019.461 & 1019.475 & 13.732747\allowbreak{}0104 \\
\(J/\allowbreak{}\psi\) & 64 & 237 & 0 & 3096.916 & 3096.979 & 20.342818\allowbreak{}4684 \\
\rowcolor{HMTLight}\(\Upsilon(1S)\) & 89 & 334 & 0 & 9460.3 & 9460.34 & 4.228195\allowbreak{}72318 \\
\end{longtable}
\endgroup
Las doce cifras archivadas son evaluaciones históricas importadas; no se atribuyen a \(\Theta\) hasta localizar la carta que las reproduzca. Además, la malla no construye el morfismo

\[
\text{estado compuesto}
\longrightarrow\text{ruta y registro genealógico}
\longrightarrow(n,k,t).
\]
Además conserva la referencia externa en la misma tabla histórica. Por ello su fuerza es evidencia numérica condicionada, no selección física anterior al contraste. Publicarla amplía el censo visible sin convertir cinco resonancias en cinco predicciones anteriores al contraste ni desplazar las evaluaciones E3 vigentes.
\fi

\Needspace{7\baselineskip}
\subsection{Normalización de identidades nominales en registros comparables}
El empalme prospectivo completo tiene la forma

\[
\begin{aligned}
\text{identidad física}
&\longrightarrow
\text{sector y representación}
\longrightarrow
\text{proyector de fibra o árbol compuesto}\\
&\longrightarrow
\text{ruta enriquecida}
\longrightarrow
\text{registro genealógico}
\longrightarrow
\text{firma y torres}\\
&\longrightarrow
\text{operador preparado}
\longrightarrow
\text{polo o autovalor}
\longrightarrow
\text{carta dimensional}\\
&\longrightarrow
\text{comparación externa}.
\end{aligned}
\]
Para un estado estable, la salida operatoria es un autovalor real. Para una resonancia, es un polo complejo. Para un quark, el mapa incluye esquema y escala. Para un límite experimental, incluye el nivel de confianza. Para un observable correlacionado, incluye la matriz de covarianza. Estas diferencias son estructurales: determinan qué magnitudes pueden compararse.

\Needspace{7\baselineskip}
\subsection{Condiciones de clausura individual de las realizaciones materiales}
Una fila del inventario queda individualizada desde HMT--MD cuando se han construido, sin consultar su magnitud externa:

\begin{enumerate}
\item el sector y la representación;
\item la ruta, multisección o árbol de composición;
\item el registro genealógico y la firma;
\item el operador de masa y, si procede, el operador de canales;
\item el proyector que selecciona el estado nombrado;
\item la convención de masa, anchura, esquema y escala;
\item la carta dimensional;
\item el falsador interno.
\end{enumerate}
La ley universal ya proporciona la gramática común. El trabajo por estado no consiste en inventar una fórmula diferente, sino en construir el entrelazador que transporta una identidad física hasta el operador apropiado.

\Needspace{7\baselineskip}
\subsection{obstrucciones y condiciones de no degeneración específicos}
La composición queda refutada si dos parentizaciones físicamente equivalentes dan firmas o espectros distintos. El selector nominal queda refutado si cambia al ocultar la columna externa de masa. El lector de anchura queda refutado si produce \(\Gamma<0\) para una contracción, confunde amplitud con probabilidad o asigna anchura a una sección estable sin canal abierto. La carta de resonancia queda refutada si no declara masa de polo o Breit--Wigner. La comparación quark queda indeterminada si omite esquema o escala. La igualdad partícula--antipartícula de anchuras queda refutada si el operador de canales no es covariante bajo conjugación.

\Needspace{7\baselineskip}
\subsection{Transporte de las clases materiales al corredor excepcional mediante rutas, firmas y representaciones}
La gramática de composición, los dominios mesónico y bariónico, el operador de Feshbach, el censo 471/384 y el principio de comparación se conservan en el desarrollo integral antecedente, secciones 16 y 18. Los registros individualizados se presentan íntegramente en los inventarios exhaustivos de masas y anchuras del apéndice. La corrección del signo del polo se obtiene directamente de \(\lambda=\exp(-izt_0/\hbar_{\rm int})\) y fortalece la formulación sin alterar ninguno de esos registros.

\Needspace{7\baselineskip}
\section{Sectores topológicos, extensiones hipotéticas y gravedad}
\Needspace{7\baselineskip}
\subsection{Separación tipada entre sectores topológicos, extensiones hipotéticas y realización gravitatoria}
La arquitectura distingue con precisión:

\begin{enumerate}
\item una identidad algebraica interna;
\item un codominio de fusión o trenzado;
\item una partícula física con Hamiltoniano, brecha, localización y canal de
medida.

\end{enumerate}
Una regla de fusión exacta no es todavía una masa. Un modo de borde no es todavía una especie. Una involución de rutas no es todavía una condición de Majorana. El paso físico exige un funtor que preserve producto, involución, orientación, energía y composición de caminos.

\Needspace{7\baselineskip}
\subsection{Sector de Fibonacci: anillo basado, fusión y representación de trenzas}
La incidencia areal--volumétrica está gobernada por

\[
F_{\rm av}=
\begin{pmatrix}
0&1\\[2pt]1&1
\end{pmatrix}.
\]
Su polinomio característico es

\[
x^2-x-1,
\]
y su autovalor dominante es \(\varphi\). En la base \(\{\mathbf1,\tau\}\), la multiplicación por \(\tau\) se representa por \(F_{\rm av}\); por tanto

\[
\tau\otimes\tau\cong\mathbf1\oplus\tau,
\qquad
\operatorname{FPdim}(\tau)=\varphi.
\]
Esta es una realización exacta del anillo basado de Fibonacci. Para obtener una categoría anyónica física se deben construir además el asociador \(F\), el trenzado \(R\), las ecuaciones de pentágono y hexágono, la quiralidad, un Hamiltoniano local con brecha y el mapa de preparación y lectura. La salida de masa, cuando existe, es una brecha de energía convertida por una velocidad efectiva:

\[
m_{\rm brecha}=\frac{\Delta E}{c_{\rm eff}^{2}}.
\]
No existe una masa universal del anyón Fibonacci deducible únicamente de \(\varphi\).

\Needspace{7\baselineskip}
\subsection{Separación entre el sector topológico y el centro apuntado del atlas}
El centro trenzado del sector apuntado basado en \(\mathbb Z/9\mathbb Z\) posee 81 objetos simples invertibles, cada uno de dimensión cuántica uno. Un simple Fibonacci exigiría dimensión \(\varphi\), y un simple Ising exigiría \(\sqrt2\). Por tanto ninguno de los dos aparece como objeto simple interno de ese centro apuntado.

Este control no elimina las realizaciones topológicas. Determina su lugar: deben obtenerse mediante extensión, condensación, cociente, defecto o funtor monoidal hacia un codominio no apuntado. Presentar el anillo Fibonacci como si ya fuese toda la categoría borraría precisamente el morfismo físico que se debe construir.

\Needspace{7\baselineskip}
\subsection{Sector Ising y modos de Majorana}
La categoría de Ising tiene objetos simples

\[
\mathbf1,\quad\psi,\quad\sigma
\]
y reglas

\[
\psi\otimes\psi=\mathbf1,
\qquad
\psi\otimes\sigma=\sigma,
\qquad
\sigma\otimes\sigma=\mathbf1\oplus\psi,
\qquad
d_\sigma=\sqrt2.
\]
Cuatro objetos distintos suelen recibir nombres próximos y deben conservarse separados:

\begin{itemize}
\item la involución de rutas \(C_M=-\operatorname{rev}\);
\item la paridad de la completación CAR;
\item un espinor relativista autorconjugado;
\item un modo cero topológico de Majorana en una fase Ising.
\end{itemize}
Un fermión relativista es de Majorana cuando existe una conjugación real \(\mathcal C\) compatible con su representación y dinámica,

\[
\psi=\mathcal C\bar\psi^{T},
\]
y el término de masa respeta esa condición. Un modo cero de Majorana se describe mediante un operador autoadjunto \(\gamma=\gamma^\dagger\) que satisface

\[
\{\gamma_i,\gamma_j\}=2\delta_{ij}
\]
y permanece separado del continuo por una brecha. La igualdad de una palabra con su imagen por \(C_M\) no demuestra ninguna de estas dos realizaciones.

Para los neutrinos, la alternativa Dirac/Majorana sólo se decide después de construir el operador neutro, su base de masas, el entrelazador de mezcla y la estructura real fermiónica. La neutralidad eléctrica es necesaria para una autorconjugación no trivial, pero no suficiente.

\Needspace{7\baselineskip}
\subsection{Trenzado abeliano de orden \(6\)}
Sea

\[
\operatorname{wind}_6:B_n\longrightarrow\mathbb Z/6\mathbb Z
\]
el lector de devanado y

\[
\chi_k(b)=\exp\!\left(
\frac{2\pi i k}{6}\operatorname{wind}_6(b)\right).
\]
Las realizaciones con fase \(q\) descienden a meras permutaciones sólo cuando

\[
q^2=1.
\]
Los cuatro caracteres restantes distinguen clases algebraicas del lector de devanado. El resultado exacto es la clasificación de los seis caracteres unidimensionales del cociente de devanado; no es todavía una clasificación de seis fases materiales. Para afirmar que conservan trenzado físico se requieren operadores de trenza y una representación del grupo de trenzas. Su existencia no implica cuatro nuevas partículas: una especie anyónica exige un Hamiltoniano, operadores de trenza, brecha, estabilidad y un acoplamiento a una multisección física. Su promoción a fases materiales depende de esa realización.

\Needspace{7\baselineskip}
\subsection{Positividad y exclusión del taquión persistente}
La recta de acción posee orientación positiva,

\[
\hbar_{\rm int}>0,
\]
el carácter de masa toma valores positivos,

\[
\chi_{\rm full}(\sigma,\eta)>0,
\]
y el refinamiento bidireccional se implementa por conjugación. Supóngase explícitamente

\[
\widehat M^{(0)}\ge0,
\qquad
\widehat K=\widehat K^*,
\qquad
R_{\rm act}>0.
\]
Entonces

\[
\widehat M^{\beta}
=R_{\rm act}^{\widehat K/2}
\widehat M^{(0)}
R_{\rm act}^{\widehat K/2}\ge0.
\]
En consecuencia,

\[
\operatorname{Spec}(\widehat M^\beta)\subset[0,\infty).
\]
Por tanto el operador de masa al cuadrado

\[
\widehat{M^2}_{\beta}:=(\widehat M^\beta)^2
\]
es positivo y satisface

\[
\operatorname{Spec}(\widehat{M^2}_{\beta})\subset[0,\infty).
\]
Una partícula asintótica persistente que satisfaga las premisas de positividad y unitariedad no puede tener \(m^2<0\). Si la linealización de un fondo produce frecuencia imaginaria, el objeto es una dirección inestable: el fondo debe relajarse o condensar. Si la ruta deja de prolongarse, el objeto es una sección virtual con obstrucción terminal. Ninguno de los dos casos define una partícula estable superlumínica.

El alcance de la proposición está fijado por sus premisas. Un parámetro "taquiónico" en un potencial perturbativo puede señalar una transición de fase; lo excluido es su promoción a especie persistente de norma positiva.

\Needspace{7\baselineskip}
\subsection{Exclusión independiente de fantasmas}
Un fantasma físico sería un modo de norma negativa. Sea \(\mathcal H_{\rm phys}=\ker G/\operatorname{im}G_0\) el cociente de calibre dotado de producto interno positivo, y sea \(\widehat M_{\rm phys}\) autoadjunto sobre él. Entonces la medida espectral de un estado preparado satisface

\[
\mu_{P,a}^{\rm mass}(B)\ge0.
\]
Una dirección negativa indica que falló la positividad del cociente o que el modo no pertenece al espacio físico; no desaparece por una afirmación nominal. Los fantasmas de Faddeev--Popov son variables del complejo de calibre, no partículas asintóticas. Esta exclusión es distinta de la taquiónica: un fantasma viola positividad de norma; una dirección taquiónica viola estabilidad del fondo.

\Needspace{7\baselineskip}
\subsection{Fotón, sesgo unidireccional y puente Proca--Stückelberg}
Los canales armónicos producen el sesgo

\[
\varepsilon=q_-^{90}-q_+^{90}.
\]
Una realización cinemática bidireccional puede escribirse

\[
c_+=\frac{c_{\rm int}}{1+\varepsilon},
\qquad
c_-=\frac{c_{\rm int}}{1-\varepsilon},
\]
para la cual la medida de ida y vuelta conserva exactamente

\[
\frac{2}{c_+^{-1}+c_-^{-1}}=c_{\rm int}.
\]
El sesgo permanece en la lectura unidireccional y cancela en la lectura recíproca. Esta identidad pertenece a la arquitectura de dos orientaciones; no demuestra por sí sola una masa fotónica.

Si el cierre de frontera produce una frecuencia de brecha \(\mu_T\in\mathcal L_T^{-1}\) en el registro,

\[
\omega^2=c_{\rm int}^{2}k^2+\mu_T^2,
\]
un entrelazador dimensional puede transportarlo a una realización Proca--Stückelberg con masa efectiva

\[
m_\gamma=\frac{\hbar_{\rm int}\mu_T}{c_{\rm int}^{2}}.
\]
Para una propagación de frecuencia \(\omega\), la corrección temporal de primer orden satisface

\[
\frac{\Delta T}{T}
\simeq
\frac12
\left(
\frac{m_\gamma c_{\rm int}^{2}}
     {\hbar_{\rm int}\omega}
\right)^2.
\]
La cadena requerida es

\[
\text{ida/retorno TPK}
\longrightarrow
\text{memoria}
\longrightarrow
\mu_T
\longrightarrow
m_\gamma.
\]
No se identifica automáticamente \(\varepsilon\), el cociente de acción \(R_{\rm act}\) y \(\mu_T\): viven en codominios distintos. La carta nula del fotón ordinario permanece exacta mientras no se construya la brecha y su entrelazador físico.

\Needspace{7\baselineskip}
\subsection{Clasificación operatoria de los sectores oscuros por dominio, lector y codominio físico}
Un sector oscuro no se define por una masa arbitraria, sino por el núcleo de sus acoplamientos. Si \(P_X\) es su proyector,

\[
P_{\rm em}P_X=0,
\qquad
P_{\rm color}P_X=0,
\]
mientras puede ocurrir

\[
P_{\rm grav}P_X\ne0
\]
o existir un acoplamiento débil. La ley entrega primero un operador de fibra; el canal determina después si su espectro contiene un modo estable, metaestable o resonante.

Un fotón oscuro requiere una segunda conexión \(A'\), su curvatura \(F'\) y un entrelazador cinético

\[
\mathcal L_{\rm mix}
=-\frac{\varepsilon_{\rm kin}}2F_{\mu\nu}F'^{\mu\nu}.
\]
Si el segundo sector adquiere masa por Stückelberg, Higgs oculto o frontera, su polo es distinto del fotón nulo. La magnitud \(\varepsilon_{\rm kin}\) debe derivarse de una incidencia común antes de comparar observables.

Sean \(P_{\rm weak},P_{\rm em},P_{\rm color}\) los proyectores sobre los canales de representación débil, electromagnético y cromático, y \(P_X\) el soporte de la sección oscura. Un WIMP exige solapamiento débil no nulo y ausencia de soporte electromagnético y cromático:

\[
P_{\rm weak}P_X\ne0,
\qquad
P_{\rm em}P_X=P_{\rm color}P_X=0.
\]
El primer producto expresa incidencia de subespacios; un conmutador no nulo no sería un criterio de acoplamiento, pues dos proyectores pueden conmutar y, aun así, compartir soporte. La intensidad de la interacción requiere después un operador de vértice o un Hamiltoniano de acoplamiento.

Su estabilidad requiere un autovalor unitario del transporte o una paridad protectora. "WIMP" es una condición de representación y persistencia, no una masa predeterminada.

\Needspace{7\baselineskip}
\subsection{Clasificación de la sección neutrínica estéril por fibra invariante, mezcla y carga de calibre nula}
Un neutrino estéril es un autovector del bloque neutro que satisface

\[
P_{\rm weak}n_s=0,
\qquad
P_{\rm mass}n_s=n_s.
\]
No es sinónimo de la cuarta profundidad estructural del soporte neutrínico. Su existencia requiere un núcleo débil no nulo, un autovalor del operador de masa en ese núcleo y un transporte de mezcla con los modos activos. El criterio es operatorio y puede falsarse antes de introducir una escala externa.

\Needspace{7\baselineskip}
\subsection{Clasificación de las secciones axiónica y pseudoescalares por carácter topológico y acoplamiento de borde}
Un candidato axiónico debe ser un modo pseudoescalar con término cinético normalizado,

\[
CP\,a\,CP^{-1}=-a,
\qquad
\mathcal L_{\rm kin}=\frac{Z_a}{2}\,\partial_\mu a\,\partial^\mu a,
\qquad Z_a>0,
\]
acoplado a la densidad topológica de color. Su masa procede de la curvatura del potencial efectivo,

\[
m_a^2=\frac1{Z_a}
\left.
\frac{\partial^2V_{\rm eff}}{\partial a^2}
\right|_{a=a_0}.
\]
El canal \(C^*\mapsto-C^*\) aporta una orientación impar, pero no basta para construir \(V_{\rm eff}\), el término cinético ni el proyector físico. La especie sólo queda definida cuando esos mapas están presentes.

\Needspace{7\baselineskip}
\subsection{Realizaciones monopolares y defectos topológicos: clases de carga y estabilidad}
Sea \(F\) la curvatura de una conexión \(U(1)\) en la normalización integral elegida. Un monopolo magnético se caracteriza por la clase

\[
\frac1{2\pi}\int_{S^2}F=n\ne0.
\]
HMT debe realizar \(n\) como holonomía neta de una familia cerrada de rutas y demostrar que sobrevive al refinamiento. Una vacancia puntual o un acarreo aislado no son un monopolo. Si el defecto es dinámico y estable, su masa es un autovalor del operador linealizado alrededor de esa clase.

\Needspace{7\baselineskip}
\subsection{Clasificación de las extensiones supersimétricas mediante emparejamiento de representaciones y caracteres}
La coexistencia de completaciones fermiónica y bosónica no demuestra supersimetría. En cuatro dimensiones se necesita una familia de cargas impares que satisfaga

\[
Q_\alpha:\mathcal F_-\longleftrightarrow\mathcal F_+,
\qquad
\{Q_\alpha,\bar Q_{\dot\beta}\}
=2\sigma^\mu_{\alpha\dot\beta}P_\mu,
\]
que entrelace los operadores de masa:

\[
Q_\alpha\widehat M_-=\widehat M_+Q_\alpha.
\]
La identidad \(Q^2=H\) se reserva para una reducción hermítica unidimensional explícita; no sustituye al álgebra relativista anterior.

Sólo entonces aparecen pares espectrales. La ruptura debe derivar la separación de masa sin usar el valor del compañero buscado. Selectrones, neutralinos, gluinos u otros nombres permanecen como realizaciones posibles, no como filas predichas por la mera existencia de CAR y CCR.

\Needspace{7\baselineskip}
\subsection{Clasificación de las secciones de dilatón y radión mediante modos escalares de volumen y compactificación}
El dilatón es una fluctuación de la sección de escala; el radión, de una longitud o volumen interno. Sus acciones elementales son

\[
t_0\mapsto e^{\phi_d}t_0,
\qquad
\mathcal V_{\rm int}\mapsto e^{\phi_r}\mathcal V_{\rm int}.
\]
Para constituir partículas necesitan término cinético, potencial, proyector persistente y polo. Antes de ello son grados colectivos de la carta dimensional. Esta distinción evita convertir una libertad de normalización en una nueva especie escalar.

\Needspace{7\baselineskip}
\subsection{Compuestos y energía de enlace}
Un estado compuesto \(Y\) con constituyentes \(X_1,\ldots,X_k\) se construye en el orden

\[
(\Gamma_{X_1},\ldots,\Gamma_{X_k})
\longrightarrow
\mathscr L_{\rm bind}
\longrightarrow
\sigma_Y
\longrightarrow
\eta_Y
\longrightarrow
\widehat M_Y.
\]
El registro de enlace \(\mathscr L_{\rm bind}\) contiene el orden de las rutas, la frontera común, el acarreo transferido, el canal de color y la holonomía de ligadura. No es lícito reemplazarlo por \(\sum_i m_{X_i}\). Mesones y bariones requieren lectores distintos porque sus números de constituyentes, representaciones y simetrías de intercambio son diferentes.

Las seis evaluaciones compuestas publicadas prueban la evaluación del carácter al fijar la firma. El cierre estado por estado exige además que el registro de enlace regenere esa firma. Esta condición se aplica del mismo modo a núcleos, átomos, moléculas y candidatos tetraquark o pentaquark, con sus respectivos Hamiltonianos de frontera.

\Needspace{7\baselineskip}
\subsection{Anchura como 2.º observable}
Las 384 anchuras no se obtienen añadiendo una sexta coordenada a \(\sigma\). Para cada estado inestable se construyen los canales \(\mathcal H_{\rm in}\oplus\mathcal H_{\rm out}\), el operador efectivo \(H_{\rm eff}(z)\) y el polo

\[
z_X=M_Xc_{\rm int}^{2}-\frac{i}{2}\Gamma_X.
\]
La parte real pertenece a la ley de masa; la parte imaginaria, al lector de decaimiento. La relación \(\tau_X=\hbar_{\rm int}/\Gamma_X\) se interpreta después de fijar la carta temporal. Así, una teoría exhaustiva de 471 masas y 384 anchuras contiene dos lectores coordinados y no una sola columna con 855 números heterogéneos.

\Needspace{7\baselineskip}

\endgroup
\begingroup
\let\chapter\section
\let\section\subsection
\let\subsection\subsubsection
\let\subsubsection\paragraph
\let\clearpage\relax
% La subsección inicial de gluones coincide byte por byte con la residencia
% REV.2 anterior; la continuación conjunta fotón--gluón se conserva íntegra.
% Extracto editor-ready del contenido exclusivo de masas/16_hidrodinamica_y_gauge.tex.
% Las líneas 1--14 están publicadas byte-exactamente desde
% 74_rev2_realizaciones_i.tex:544--557. El propietario íntegro permanece
% preservado; este extracto conserva la continuación exclusiva.

\subsection{Comparación de las realizaciones nulas electromagnética y cromodinámica}
El fotón y los ocho gluones pertenecen a cartas nulas distintas. Para el fotón,

\[
m_\gamma=0
\]
en la rama electromagnética. Para los gluones,

\[
m_g=0
\]
en la representación adjunta de color. No se obtiene la nulidad tomando el límite de un carácter positivo; el dominio se bifurca antes. La realización fotónica exige Maxwell, invariancia de calibre y dos polarizaciones transversales. La realización gluónica exige la conexión de color, sus ocho generadores y la restricción al espacio físico de calibre.

La propagación bidireccional puede escribirse

\[
c_\pm=\frac{c}{1\pm\epsilon},
\qquad
\epsilon=q_-^{90}-q_+^{90},
\]
con media armónica

\[
\frac{2}{1/c_++1/c_-}=c.
\]
La velocidad bidireccional conserva \(c\); la orientación unidireccional conserva el sesgo. Una brecha adicional del registro genealógico puede realizarse mediante Proca--Stückelberg, pero no debe identificarse automáticamente con \(\epsilon\) ni con el desdoblamiento de la acción.

\Needspace{7\baselineskip}

\endgroup
\begingroup
\let\chapter\section
\let\section\subsection
\let\subsection\subsubsection
\let\subsubsection\paragraph
\let\clearpage\relax
\section{Corredor excepcional, teoría conforme, cuerdas y clasificación de masas}
\Needspace{7\baselineskip}
\subsection{Preimagen común y coordinación de las \(2\) lecturas del corredor excepcional}
La relación entre masas y simetrías excepcionales no comienza en una acción del Monstruo sobre una tabla de partículas. Comienza en el mismo estado enriquecido del TPK. Sea

\[
x\in\mathcal X_{\rm TPK}^{\rm enr}.
\]
La lectura arquimediana contrae cilindros compatibles y conserva el valor real con su genealogía:

\[
\mathcal P_{\rm arq}(x)\in\mathbb R.
\]
La lectura excepcional conserva incidencias de frontera y produce una clase en el límite proyectivo:

\[
\mathcal P_{\rm exc}(x)=Q_\infty(x).
\]
La naturalidad exige, para todo refinamiento \(n\ge m\),

\[
r_{n,m}\circ Q_n=Q_m\circ p_{n,m}.
\]
La aplicación conjunta

\[
\mathcal P_{\rm joint}:x\longmapsto
\bigl(\mathcal P_{\rm arq}(x),\mathcal P_{\rm exc}(x)\bigr)
\]
es la forma positiva de la doble proyección. No niega la conexión Monstruo--masa: la tipa. La rama de masa y la rama excepcional son imágenes de una misma preimagen con orientación, gradación, acarreo, memoria y holonomía comunes. Tampoco afirma que el Monstruo actúe literalmente sobre las cifras decimales o que un coeficiente de Moonshine sea por sí mismo una masa.

\Needspace{7\baselineskip}
\subsection{Construcción del corredor excepcional}
La incidencia ternaria prolonga el estado por la cadena

\[
C_3
\longrightarrow A_2^{12}
\longrightarrow N(A_2^{12})
\longrightarrow \Lambda_{24}
\longrightarrow V^\natural
\longrightarrow \mathbb M,
\]
donde \(\Lambda_{24}\) es la red de Leech, \(V^\natural\) el álgebra de vértices del Moonshine y \(\mathbb M\) el grupo Monstruo. El paso no consiste en yuxtaponer nombres clásicos: cada flecha debe preservar la incidencia, la gradación y la compatibilidad de refinamiento heredadas del estado.

La traza ternaria se escribe

\[
t_3(q)=q^{-1}
\prod_{n\ge1}(1+q^n+q^{2n})^{-12}.
\]
El primer defecto no trivial satisface

\[
D_{\rm APP}=[q]t_3=v_\alpha^2=54,
\]
y el primer coeficiente del carácter del Monstruo conserva la descomposición

\[
196884=54+10\cdot3^9.
\]
Estas identidades transportan gradación y multiplicidad. No asignan una masa en MeV. Su función en la teoría de masas es clasificar el espacio sobre el que actúa el operador antes de escalarizarlo.

\Needspace{7\baselineskip}
\subsection{Emparejamiento fibra a fibra con la ley de masas}
La rama material del mismo estado produce

\[
x\longmapsto
\bigl(F_X,\sigma_X,\eta_X,\widehat M_X\bigr),
\]
donde \(F_X\) es punto, órbita, fibra o multisección; \(\sigma_X\in\mathbb Z^5\) es la firma; \(\eta_X\in\mathcal E_{90,120}\) es el refinamiento; y \(\widehat M_X\) el operador de masa. La lectura conjunta es entonces

\[
\mathfrak J_{\rm mass,exc}(x)
=\bigl(F_X,\sigma_X,\eta_X,\widehat M_X,Q_\infty(x)\bigr).
\]
Dos estados con el mismo valor escalar de masa pueden permanecer distintos por su componente excepcional; dos estados con el mismo carácter excepcional pueden poseer operadores de masa diferentes. La clasificación correcta es el emparejamiento de fibras, no la igualdad aislada de números.

Si una representación excepcional \(G\) actúa en una fibra física \(\mathcal H_X\), la compatibilidad con la masa exige

\[
U_g\widehat M_XU_g^{-1}=\widehat M_X,
\qquad g\in G.
\]
La descomposición

\[
\mathcal H_X=\bigoplus_{\lambda}
\mathcal H_{X,\lambda}
\]
separa multipletes excepcionales. Sobre un sumando irreducible, la conmutación fuerza que la compresión de \(\widehat M_X\) sea escalar. Ésta es una vía legítima de escalarización por simetría. El valor del escalar sigue procediendo del carácter, las torres, la sección dimensional y el acoplamiento; la dimensión de la representación sólo determina la degeneración.

\Needspace{7\baselineskip}
\subsection{Fock y álgebra de vértices}
La prolongación de una ruta a un espacio de muchas partículas posee dos completaciones principales. La exterior produce CAR, Pauli y estadística fermiónica; la simétrica produce CCR y estadística bosónica. Una álgebra de vértices añade producto operatorio, vacío, traslación y gradación:

\[
V=\bigoplus_{n\ge0}V_n,
\qquad
L_0|_{V_n}=nI.
\]
El carácter

\[
\operatorname{ch}_V(q)
=\operatorname{Tr}_V q^{L_0-c/24}
\]
cuenta estados por grado. Una multiplicidad en \(V_n\) no es una masa. Para convertir la gradación en energía se necesita un Hamiltoniano y una sección de longitud o tiempo. En una realización conforme sobre un cilindro de circunferencia \(L\),

\[
H_{\rm CFT}
=\frac{2\pi\hbar_{\rm int}v}{L}
\left(L_0+\bar L_0-\frac{c}{12}\right).
\]
La masa o brecha física aparece sólo después de fijar \(v\), \(L\), las condiciones de borde y el mapa desde la ruta TPK a los módulos conformes. La función del corredor excepcional consiste en organizar las degeneraciones y las reglas de composición de esos módulos.

\Needspace{7\baselineskip}
\subsection{Espectro de cuerda como realización dimensional posterior}
Sea \(\alpha'\in\mathcal L_L^2\) la longitud cuadrática que representa la tensión inversa en la carta naturalizada de cuerda declarada. En una compactificación circular de radio \(R\), los momentos izquierdo y derecho son

\[
p_L=\frac{m}{R}+\frac{wR}{\alpha'},
\qquad
p_R=\frac{m}{R}-\frac{wR}{\alpha'},
\]
con \(m,w\in\mathbb Z\). En unidades naturales de la realización, el espectro cerrado satisface

\[
M^2=p_L^2+\frac{4}{\alpha'}(N_L-a_L)
=p_R^2+\frac{4}{\alpha'}(N_R-a_R),
\]
junto con la condición de nivel

\[
N_L-N_R=a_L-a_R-mw.
\]
La involución

\[
(m,w,R)\longleftrightarrow
\left(w,m,\frac{\alpha'}{R}\right)
\]
conserva el espectro: es la realización de dualidad T. HMT aporta antes la orientación bidireccional, la memoria y el emparejamiento de hojas; la teoría de cuerda añade tensión, radio, osciladores, nivel y condiciones de contorno. No se obtiene una masa de cuerda reemplazando \(N_L,N_R\) por alturas de la torre \(\langle90,120\rangle\). Toda identificación entre ambas gradaciones requiere un entrelazador que conserve suma, orientación y carácter.

\Needspace{7\baselineskip}
\subsection{Pantallas reticulares de teoría M: gradación, compactificación y transporte de la memoria residual}
Las pantallas internas

\[
12\longrightarrow11\longrightarrow10,
\qquad
26\longrightarrow24,
\]
conservan el descenso dodecafásico, la eliminación de una dirección redundante y la pantalla transversal. En una realización de teoría M, los campos

\[
g_{MN},\qquad C_{MNP},\qquad\Psi_M
\]
deben obtenerse mediante un entrelazador que preserve calibre, supersimetría, restricciones y producto interno. Las pantallas numéricas no sustituyen ese mapa.

La clasificación de masas se transporta a esta rama mediante

\[
\text{ruta}
\longrightarrow\text{registro genealógico}
\longrightarrow\text{firma}
\longrightarrow\text{carácter}
\longrightarrow\text{operador}
\longrightarrow\text{espectro compactificado}.
\]
Los modos de Kaluza--Klein, devanado, membrana o excitación oscilatoria son estados distintos aunque compartan una masa. El gravitón de una realización de cuerda cerrada, si aparece, es un modo colectivo sin masa del sector geométrico; no se convierte por ello en semilla elemental del TPK ni invalida la formulación de la gravedad como holonomía y torsión.

\Needspace{7\baselineskip}
\subsection{Realizaciones ordinarias e hipotéticas inducidas por el corredor excepcional}
La rama excepcional permite distinguir cuatro fuentes de multiplicidad:

\begin{enumerate}
\item multiplicidad celular o de ruta en el atlas;
\item degeneración de una representación excepcional;
\item ocupación de Fock;
\item nivel conforme, oscilatorio o compactificado.
\end{enumerate}
Confundirlas produce falsos compañeros, falsas generaciones o masas duplicadas. Un estado gemelo, supersimétrico, oscuro o exótico sólo constituye una nueva especie cuando su quíntuplo

\[
(F_X,Q_\infty(x),\rho_X,\widehat M_X,P_X)
\]
difiere físicamente y el Hamiltoniano posee un polo o una sección persistente adicional. Una coincidencia de dimensión, carácter o masa no basta.

En sentido inverso, una degeneración protegida por el corredor excepcional es una predicción estructural: perturbaciones que preserven la simetría no pueden escindir el multiplete. Una escisión observada exige identificar la ruptura de simetría, el acoplamiento o la carta que la produce.

\Needspace{7\baselineskip}
\subsection{Condiciones de positividad, cierre y estabilidad de las realizaciones topológicas}
La conexión entre el corredor excepcional y las masas queda falsada en una realización concreta si ocurre cualquiera de los hechos siguientes:

\begin{itemize}
\item falla la naturalidad \(r_{n,m}Q_n=Q_mp_{n,m}\);
\item la etiqueta excepcional cambia al permutar las masas externas;
\item una multiplicidad de Moonshine se usa como valor de masa sin Hamiltoniano;
\item el operador de masa no conmuta con la simetría que se afirma preservada;
\item se identifica una altura de torre con un nivel de cuerda sin entrelazador;
\item falla la condición de nivel o la invariancia por dualidad T;
\item una pantalla dimensional se declara teoría M sin construir sus campos;
\item se usa la existencia de un modo de espín dos para redefinir la gravedad
interna anterior.

\end{itemize}
Estos controles mantienen positiva la conexión: existe un estado común y una clasificación conjunta; lo que se prohíbe es sustituir las flechas por una analogía numérica.

\Needspace{7\baselineskip}
\subsection{Construcción matemática de las clases materiales a partir de rutas, firmas y representaciones}
La doble realización, la naturalidad proyectiva, la cadena \(A_2^{12}\)--Leech--\(V^\natural\)--Monstruo, la traza \(t_3\), las identidades \(D_{\rm APP}=54\) y \(196884=54+10\cdot3^9\), y las pantallas dimensionales pertenecen a la arquitectura matemática APP--TRIT--TPK desarrollada previamente. La formulación de \(\mathfrak J_{\rm mass,exc}\), la separación de las cuatro multiplicidades y los controles anteriores hacen explícita su función dentro de una teoría autónoma de masas. Las fórmulas conforme y de cuerda son realizaciones dimensionales posteriores con sus parámetros declarados; no se emplean para generar retroactivamente APP, TRIT, TPK, el carácter ni las torres.

\endgroup

\section{Resultados operatoriales vigentes y contraste metrológico posterior}
La comparación sólo se abre después de congelar ruta, firma, operador,
esquema, escala y unidad. La tabla vigente dispone en columnas contiguas la
salida HMT, el valor externo, la incertidumbre, el residual y el estatuto. Los
valores quark no admiten un residual público si esquema y escala no coinciden.

\begingroup
\let\chapter\section
\let\section\subsection
\let\subsection\subsubsection
\let\subsubsection\paragraph
\let\clearpage\relax
\section{Operadores de masa en HMT--MD: escala absoluta, atlas de rutas y realización por sectores}
\label{sec:tabla-completa-masas-hmt-sm-revfinal}

La masa aparece en Mecánica Dimensional como la realización de dos ramas
matemáticas que confluyen sin confundirse. La primera conserva la genealogía del
estado: una extensión superviviente determina una ruta observable; la ruta escribe
un registro cronológico dodecafásico; el registro produce una firma entera, y la firma
alimenta el carácter y la jerarquía armónica. La segunda construye la recta
dimensional: el lector determinantal fija la escala de acción, el reloj interno la
transporta a energía y la carta cinemática la lleva a masa. El operador de especie
actúa entonces por homotecias positivas sobre esa recta. La coordenada experimental
se abre solamente después de congelar ambas ramas.

El dominio interno completo posee cuatro censos distintos:
\begin{equation}
 81=25_{\ell^-}+20_{\nu}+20_d+16_u,\qquad
 56=41\!\times\!1+10\!\times\!2+2\!\times\!3+1\!\times\!4+2\!\times\!5,
 \label{eq:censos-81-56-masas-rectoras}
\end{equation}
\begin{equation}
 13\ \text{multisecciones},\qquad
 324=81\times4\ \text{rutas orientadas}.
 \label{eq:censos-13-324-masas-rectoras}
\end{equation}
Las realizaciones físicas comparadas más adelante constituyen una
clasificación de realizaciones; no son otro nombre para las 81 celdas, las 56
familias, las 13 multisecciones o las 324 historias orientadas. Sobre este dominio
los dos lectores de ruta producen catorce perfiles \(K_D\), nueve perfiles
\(K_\Omega\), cuarenta pares conjuntos y dieciocho coincidencias. El sector nulo
permanece separado del carácter positivo de masa.

\subsection{De la incompatibilidad APP a la firma de masa}

La diferencia primordial entre las hojas de suma y producto aparece en la
descomposición espectral del bloque central de APP:
\begin{equation}
 B_c=7I+2R=9P_++5P_-,\qquad
 \operatorname{Spec}(B_c)=\{9,5\},\qquad
 \Delta_{\rm APP}=9-5=4.
 \label{eq:salto-app-cuatro-ley-masas}
\end{equation}
Los enteros que reaparecen después conservan fuentes distintos. El coeficiente
local de intercambio es \(2\) en \(B_c=7I+2R\); la compresión modular publica
\(51-45=6\), y su despliegue sobre nueve bloques produce \(9\cdot6=54\). Estas
igualdades pertenecen a niveles tipados de una misma genealogía y no constituyen
una cadena de definiciones del salto \(4\). En la sección electrónica, \(54\) es
además \(D_1(\Gamma_e)\) y la periodicidad cinemática de orden \(54\) del
endomorfismo temporal CT108: mide una componente de la historia ya transportada
por TPK, no la separación espectral primordial.

Para cada ruta enriquecida \(\Gamma\), el registro dodecafásico determina
\begin{equation}
 \sigma(\Gamma)
 =(n_A,s_C,k_{\Delta_4},b_{120},b_\zeta)\in\mathbb Z^5,
 \qquad \zeta_{270}=135\Delta_4^2.
 \label{eq:firma-z5-masas-rectoras}
\end{equation}
El carácter central y su refinamiento armónico son
\begin{align}
 \log\chi_0(\sigma)
 &=n_A A_{\rm rad}+\frac{s_C}{6}C^*_{\rm rad}
   +k_{\Delta_4}\Delta_4+b_{120}s_{120}+b_\zeta\zeta_{270},
 \label{eq:caracter-central-masas-rectoras}\\
 \log\chi_{\rm full}(\sigma,\eta)
 &=n_A A_{\rm rad}+\frac{s_C}{6}C^*_{\rm rad}
   +k_{\Delta_4}\Delta_4+b_\zeta\zeta_{270}
   +(b_{120}+\eta_{120})s_{120}
   +\sum_{\substack{h\in\mathfrak S\\h\ne120}}\eta_hs_h.
 \label{eq:caracter-completo-masas-rectoras}
\end{align}
La división \(s_C/6\) procede de la integralidad dodecafásica. El par etiquetado
\((b_{120},\eta_{120})\) conserva su genealogía incluso cuando su evaluación se
agrupa en un único coeficiente.

Los momentos que intervienen en
\eqref{eq:caracter-completo-masas-rectoras} pertenecen a una jerarquía infinita:
\begin{equation}
 q_\pm=\exp\!\left[-\frac{\pi}{180}(A\pm C^*)\right],
 \qquad s_h=q_-^h-q_+^h,\qquad c_h=q_-^h+q_+^h,
 \label{eq:momentos-torres-masas-rectoras}
\end{equation}
\begin{equation}
 \mathfrak S=\langle90,120\rangle
 =\{90,120\}\cup\{30r:r\ge6\}.
 \label{eq:semigrupo-global-torres-masas-rectoras}
\end{equation}
Tanto \(s_n\) como \(c_n\) satisfacen
\begin{equation}
 X_{n+2}=c_1X_{n+1}-(q_-q_+)X_n.
 \label{eq:recurrencia-torres-masas-rectoras}
\end{equation}
Las alturas \(90,120,180,210,240,270,360,420\) son testigos genealógicos
particularmente visibles; el dominio de la ley es el semigrupo completo. En
particular, \(360\) conserva dos historias antes del cociente, cuatro pasos de
\(90\) y tres de \(120\), aun cuando ambas produzcan el mismo momento escalar.

\subsection{La sección absoluta de masa y la década \texorpdfstring{\(10^{-34}\)}{10⁻³⁴}}

El bloque local de Hadamard
\[
 H_4=
 \begin{pmatrix}
  1&1&1&1\\
  1&1&-1&-1\\
  1&-1&1&-1\\
  1&-1&-1&1
 \end{pmatrix},
 \qquad
 \operatorname{SNF}(H_4)=\operatorname{diag}(1,2,2,4),
\]
tiene cokernel de orden \(16\). Su lector determinantal produce
\begin{equation}
 \Sigma_H=\varphi\alpha_{\rm HMT}^{16},
 \qquad
 10^{-34}<\Sigma_H<10^{-33},
 \qquad
 \left\lfloor\log_{10}\Sigma_H\right\rfloor=-34.
 \label{eq:orden-menos-34-tabla-masas}
\end{equation}
La conversión a \(\mathrm{J\,s}\) es una carta posterior: la potencia
\(10^{-34}\) nace en la estructura discreta y fija la escala de una única recta
de acción. Sus secciones anterior y posterior al retorno son
\begin{equation}
 \hbar_{\rm pre}=H_5\,10^{-34}\mathcal U_S,\qquad
 \hbar_{\rm ret}=\eta_{\rm ret}\,10^{-34}\mathcal U_S,
 \qquad
 R_{\rm act}=\frac{H_5}{\eta_{\rm ret}}>0.
 \label{eq:accion-absoluta-y-razon-tabla-masas}
\end{equation}
El desdoblamiento posee una carta aditiva,
\(\delta_{2w}=H_5-\eta_{\rm ret}\), y la carta multiplicativa adimensional
\(R_{\rm act}\). Declarada una duración interna \(t_0\), la periodicidad de
108 iteraciones del endomorfismo temporal CT108 construye
\begin{align}
 \omega_0&=\frac{2\pi}{108t_0},
 &\varepsilon_{\rm clk}&=\hbar_{\rm ret}\omega_0,
 \label{eq:cuanto-reloj-masa-rectora}\\
 \mathbf m_\star^{\rm pre}
 &=\left(H_5\,10^{-34}\mathcal U_S\right)
   \frac{2\pi}{108t_0}c_{\rm int}^{-2},
 &\mathbf m_\star^{\rm ret}
 &=\left(\eta_{\rm ret}\,10^{-34}\mathcal U_S\right)
   \frac{2\pi}{108t_0}c_{\rm int}^{-2}\in\mathcal L_M.
 \label{eq:seccion-reloj-masa-rectora}
\end{align}
Ambas secciones satisfacen
\(\mathbf m_\star^{\rm pre}/\mathbf m_\star^{\rm ret}=R_{\rm act}\).
La década \(10^{-34}\), el reloj común y \(c_{\rm int}^{-2}\) fijan la escala
absoluta. En una razón de masas esa sección común cancela exactamente; el
refinamiento relativo pertenece a la ruta y no a una segunda introducción de la
potencia decimal.

\subsection{Lectores bidireccionales de trayectoria y ley operatorial sobre las \(13\) multisecciones}

Los lectores que actúan sobre una ruta \(\Gamma\) son
\begin{equation}
 K_D(\Gamma)
 =\left(D_{27}+D_{36},\,D_1,\,\frac{D_4-D_3}{2}\right),
 \qquad
 K_\Omega(\Gamma)=(\Omega,54,P_6).
 \label{eq:lectores-kd-komega-masas-rectoras}
\end{equation}
Cada lector induce sobre una multisección \(F\) un operador diagonal
\(\widehat K_F^r\), \(r\in\{D,\Omega\}\). El índice \(s\) designa una sección o
bloque físico dentro de \(F\), y \(\widehat B_{F,s}^r\) es la restricción
autoadjunta de \(\widehat K_F^r\) a ese bloque. Si
\(\widehat{\mathcal A}^{(0)}_{F,s}\) es la restricción del operador que reúne la
firma, el carácter completo y la jerarquía global de torres, la ley general toma
la forma
\begin{equation}
 \boxed{
 \widehat{\mathbf M}^{\,\beta,r}_{F,s}
 =\mathbf m_\star^{\rm ret}\otimes
  R_{\rm act}^{\widehat B_{F,s}^r/2}
  \widehat{\mathcal A}^{(0)}_{F,s}
  R_{\rm act}^{\widehat B_{F,s}^r/2},
 \qquad r\in\{D,\Omega\}.}
 \label{eq:ley-operatorial-general-masas-rectoras}
\end{equation}
La escritura simétrica preserva positividad sin imponer la conmutación de
\(\widehat{\mathcal A}^{(0)}_{F,s}\) y \(\widehat B_{F,s}^r\). En la base
simultáneamente diagonal de los certificados vigentes coincide con el producto
\(\mathbf m_\star^{\rm ret}\otimes\widehat{\mathcal A}^{(0)}_F
\exp((\log R_{\rm act})\widehat K_F^r)\).
Esta ecuación gobierna las 81 celdas, las 56 familias, las 13 multisecciones y las
324 rutas orientadas. Fuera de una fibra de rango uno, la salida primaria es el
operador completo y su espectro. La realización escalar usa el morfismo cerrado
\(\mathcal S_{\rm phys}:(\sigma_{\rm ud},g,\chi_{\rm fina},[\gamma]_{\rm col})
\mapsto\gamma^{\rm surv}_{P,\mathfrak M}\) de
\cref{eq:realizacion-fisica-sabor-ruta-cerrada-rev2}; esa selección es
independiente de masas objetivo y antecede a cualquier comparación metrológica.

\subsection{Reducción de rango \(1\) en la construcción del electrón central}

La sección electrónica es el único punto fijo central del atlas:
\begin{equation}
 (r,c)=(5,5),\qquad
 \Gamma_e=(5,5,++),\qquad
 \tau_e=\mathtt{+++---+++---}.
 \label{eq:ruta-electronica-central-masas-rectoras}
\end{equation}
Sobre ella coinciden ambos lectores,
\begin{equation}
 K_D(\Gamma_e)=K_\Omega(\Gamma_e)=(80,54,6),\qquad
 \kappa_e=80-54\alpha_{\rm HMT}+6\Delta_4.
 \label{eq:kappa-electron-rector}
\end{equation}
La reducción de rango uno de
\eqref{eq:ley-operatorial-general-masas-rectoras} entrega
\begin{equation}
 \mathbf m_e^\beta
 =\mathbf m_\star^{\rm ret}\,
   \mu_e^{(0)}R_{\rm act}^{\,80-54\alpha_{\rm HMT}+6\Delta_4},
 \label{eq:electron-rank1-ley-masas}
\end{equation}
cuya coordenada en la carta comparativa es
\begin{equation}
 \boxed{m_e^\beta
 =0.5109989506900008086082571648\ldots\ \mathrm{MeV}/c^2.}
 \label{eq:masa-electron-tabla-completa-revfinal}
\end{equation}
El valor basal
\(0.5109989224880660725\ldots\ \mathrm{MeV}/c^2\) conserva su función como
etapa intermedia; la salida beta de
\eqref{eq:masa-electron-tabla-completa-revfinal} gobierna la comparación final.

\subsection{Clasificación física y comparación cuantitativa por sectores}

La exposición se organiza en cuatro estratos coordinados. El primero publica el
resultado interno y su codominio; el segundo declara el morfismo de realización
física; el tercero muestra la mejor coordenada numérica ya disponible y conserva
sus hipótesis y su dominio; el cuarto introduce el dato externo como contraste posterior.
Cada columna conserva así su tipo matemático: salida HMT, morfismo de
realización, coordenada calculada y contraste externo. La tabla publica
simultáneamente las evaluaciones disponibles y los mapas de construcción que las
producen.

El electrón publica su clausura beta escalar. El fotón y los gluones publican
secciones nulas exactas en sus cartas de masa, junto con extensiones gauge no
triviales. Los otros leptones cargados y los seis quarks publican simultáneamente
sus rutas nominales, caracteres escalares, operadores diagonales y espectros sobre
sus fibras; los neutrinos publican el operador neutro y su mezcla. \(W\), \(Z\) y
\(H\) poseen rutas CT--108 individualizadas, huellas completas y caracteres
\((n_A,s_C)=(94,-1),(95,-1),(97,9)\). El par \(K_D/K_\Omega\) pertenece al
refinamiento de las trece multisecciones del atlas y no es una premisa de la
realización bosónica por ruta nominal. Compuestos, sectores topológicos y estados
virtuales mantienen respectivamente sus codominios de ligadura, fusión, trenzado
y obstrucción.

\Needspace{15\baselineskip}
\subsubsection{Correspondencia tipada entre objetos internos y codominios físicos}


\endgroup
\begingroup
\let\chapter\section
\let\section\subsection
\let\subsection\subsubsection
\let\subsubsection\paragraph
\let\clearpage\relax

\Needspace{12\baselineskip}
\subsubsection{Operadores de las \(13\) multisecciones}

Las cifras de familias consignadas en cada ficha son cardinalidades locales del
soporte de esa multisección. No se suman entre fichas: el cociente de las
56 familias de ruta y el cociente de las 13 multisecciones son dos particiones
distintas del atlas de 81 celdas, relacionadas por los mapas de proyección y no
por una descomposición disjunta común.

\begingroup
\small
\setlength{\parskip}{.12\baselineskip}
\setlist[itemize]{leftmargin=1.65em,itemsep=0pt,topsep=.1em,
  parsep=0pt,partopsep=0pt}
\let\texttt\textnormal
Las fichas siguientes publican el operador en su codominio nativo. Cada espectro se reproduce íntegramente; ninguna lista de autovalores se sustituye por una coordenada de masa.
\Needspace{10\baselineskip}
\noindent\textbf{sector leptónico cargado, profundidad G0 (centro).} \(\mathcal B_{\ell,\mathrm{G0}}\).\par
\noindent\textbf{Dominio de fibra.} Rango 1; celda 41; sin órbita cromática; 1 familia de ruta.
\noindent\textbf{Espectro completo de \(K_D\).}
\begin{itemize}
\item \texttt{\detokenize{(80,54,6)x1}}
\end{itemize}
\noindent\textbf{Espectro completo de \(K_\Omega\).}
\begin{itemize}
\item \texttt{\detokenize{(80,54,6)x1}}
\end{itemize}
\noindent\textbf{Conclusión en la fibra.} Operador diagonal exacto sobre la fibra \(B_1\); la ruta nominal y el morfismo declarado publican el carácter escalar sin sustituir el espectro de la fibra.
\medskip
\Needspace{10\baselineskip}
\noindent\textbf{sector leptónico cargado, profundidad G2.} \(\mathcal B_{\ell,\mathrm{G2}}\).\par
\noindent\textbf{Dominio de fibra.} Rango 8; celdas 21,23,25,39,43,57,59,61; sin órbita cromática; 8 familias de ruta.
\noindent\textbf{Espectro completo de \(K_D\).}
\begin{itemize}
\item \texttt{\detokenize{(0,24,0)x1}}
\item \texttt{\detokenize{(40,78,27)x2}}
\item \texttt{\detokenize{(80,54,7)x1}}
\item \texttt{\detokenize{(80,66,2)x1}}
\item \texttt{\detokenize{(80,66,3)x1}}
\item \texttt{\detokenize{(100,66,0)x1}}
\item \texttt{\detokenize{(100,66,2)x1}}
\end{itemize}
\noindent\textbf{Espectro completo de \(K_\Omega\).}
\begin{itemize}
\item \texttt{\detokenize{(80,54,6)x8}}
\end{itemize}
\noindent\textbf{Conclusión en la fibra.} Operador diagonal exacto sobre la fibra \(B_1\); la ruta nominal y el morfismo declarado publican el carácter escalar sin sustituir el espectro de la fibra.
\medskip
\Needspace{22\baselineskip}
\noindent\textbf{sector leptónico cargado, profundidad G4.} \(\mathcal B_{\ell,\mathrm{G4}}\).\par
\noindent\textbf{Dominio de fibra.} Rango 16; celdas 1,3,5,7,9,19,27,37,45,55,63,73,75,77,79,81; sin órbita cromática; 16 familias de ruta.
\noindent\textbf{Espectro completo de \(K_D\).}
\begin{itemize}
\item \texttt{\detokenize{(0,24,0)x1}}
\item \texttt{\detokenize{(40,24,2)x2}}
\item \texttt{\detokenize{(40,54,6)x2}}
\item \texttt{\detokenize{(40,84,30)x2}}
\item \texttt{\detokenize{(60,78,25)x3}}
\item \texttt{\detokenize{(80,66,2)x2}}
\item \texttt{\detokenize{(80,66,3)x3}}
\item \texttt{\detokenize{(80,66,4)x1}}
\end{itemize}
\noindent\textbf{Espectro completo de \(K_\Omega\).}
\begin{itemize}
\item \texttt{\detokenize{(24,54,2)x1}}
\item \texttt{\detokenize{(56,54,5)x4}}
\item \texttt{\detokenize{(80,54,6)x11}}
\end{itemize}
\noindent\textbf{Conclusión en la fibra.} Operador diagonal exacto sobre la fibra \(B_1\); la ruta nominal y el morfismo declarado publican el carácter escalar sin sustituir el espectro de la fibra.
\medskip
\Needspace{10\baselineskip}
\noindent\textbf{sector neutrínico, profundidad G1.} \(\mathcal B_{\nu,\mathrm{G1}}\).\par
\noindent\textbf{Dominio de fibra.} Rango 2; celdas 40,42; sin órbita cromática; 2 familias de ruta.
\noindent\textbf{Espectro completo de \(K_D\).}
\begin{itemize}
\item \texttt{\detokenize{(40,24,2)x1}}
\item \texttt{\detokenize{(40,78,27)x1}}
\end{itemize}
\noindent\textbf{Espectro completo de \(K_\Omega\).}
\begin{itemize}
\item \texttt{\detokenize{(40,54,4)x1}}
\item \texttt{\detokenize{(80,54,6)x1}}
\end{itemize}
\noindent\textbf{Conclusión en la fibra.} Operador diagonal exacto sobre la fibra \(B_1\); la ruta nominal y el morfismo declarado publican el carácter escalar sin sustituir el espectro de la fibra.
\medskip
\Needspace{10\baselineskip}
\noindent\textbf{sector neutrínico, profundidad G2.} \(\mathcal B_{\nu,\mathrm{G2}}\).\par
\noindent\textbf{Dominio de fibra.} Rango 4; celdas 22,24,58,60; sin órbita cromática; 4 familias de ruta.
\noindent\textbf{Espectro completo de \(K_D\).}
\begin{itemize}
\item \texttt{\detokenize{(0,24,0)x1}}
\item \texttt{\detokenize{(40,84,30)x1}}
\item \texttt{\detokenize{(60,78,25)x1}}
\item \texttt{\detokenize{(80,66,3)x1}}
\end{itemize}
\noindent\textbf{Espectro completo de \(K_\Omega\).}
\begin{itemize}
\item \texttt{\detokenize{(4,54,2)x1}}
\item \texttt{\detokenize{(80,54,6)x3}}
\end{itemize}
\noindent\textbf{Conclusión en la fibra.} Operador diagonal exacto sobre la fibra \(B_1\); la ruta nominal y el morfismo declarado publican el carácter escalar sin sustituir el espectro de la fibra.
\medskip
\Needspace{18\baselineskip}
\noindent\textbf{sector neutrínico, profundidad G3.} \(\mathcal B_{\nu,\mathrm{G3}}\).\par
\noindent\textbf{Dominio de fibra.} Rango 6; celdas 20,26,38,44,56,62; sin órbita cromática; 6 familias de ruta.
\noindent\textbf{Espectro completo de \(K_D\).}
\begin{itemize}
\item \texttt{\detokenize{(40,24,2)x2}}
\item \texttt{\detokenize{(40,78,27)x1}}
\item \texttt{\detokenize{(60,84,28)x1}}
\item \texttt{\detokenize{(80,54,6)x1}}
\item \texttt{\detokenize{(80,66,3)x1}}
\end{itemize}
\noindent\textbf{Espectro completo de \(K_\Omega\).}
\begin{itemize}
\item \texttt{\detokenize{(36,54,4)x1}}
\item \texttt{\detokenize{(56,54,5)x1}}
\item \texttt{\detokenize{(80,54,6)x4}}
\end{itemize}
\noindent\textbf{Conclusión en la fibra.} Operador diagonal exacto sobre la fibra \(B_1\); la ruta nominal y el morfismo declarado publican el carácter escalar sin sustituir el espectro de la fibra.
\medskip
\Needspace{10\baselineskip}
\noindent\textbf{sector neutrínico, profundidad G4.} \(\mathcal B_{\nu,\mathrm{G4}}\).\par
\noindent\textbf{Dominio de fibra.} Rango 8; celdas 2,4,6,8,74,76,78,80; sin órbita cromática; 8 familias de ruta.
\noindent\textbf{Espectro completo de \(K_D\).}
\begin{itemize}
\item \texttt{\detokenize{(0,24,0)x1}}
\item \texttt{\detokenize{(40,24,2)x1}}
\item \texttt{\detokenize{(40,78,27)x2}}
\item \texttt{\detokenize{(40,84,30)x1}}
\item \texttt{\detokenize{(80,54,7)x1}}
\item \texttt{\detokenize{(80,66,2)x1}}
\item \texttt{\detokenize{(100,66,0)x1}}
\end{itemize}
\noindent\textbf{Espectro completo de \(K_\Omega\).}
\begin{itemize}
\item \texttt{\detokenize{(36,54,4)x1}}
\item \texttt{\detokenize{(56,54,5)x1}}
\item \texttt{\detokenize{(80,54,6)x6}}
\end{itemize}
\noindent\textbf{Conclusión en la fibra.} Operador diagonal exacto sobre la fibra \(B_1\); la ruta nominal y el morfismo declarado publican el carácter escalar sin sustituir el espectro de la fibra.
\medskip
\Needspace{10\baselineskip}
\noindent\textbf{rama quark de tipo inferior, profundidad G1.} \(\mathcal B_{d,\mathrm{G1}}\).\par
\noindent\textbf{Dominio de fibra.} Rango 2; celdas 32,50; B | G; 2 familias de ruta.
\noindent\textbf{Espectro completo de \(K_D\).}
\begin{itemize}
\item \texttt{\detokenize{(40,24,2)x1}}
\item \texttt{\detokenize{(40,78,27)x1}}
\end{itemize}
\noindent\textbf{Espectro completo de \(K_\Omega\).}
\begin{itemize}
\item \texttt{\detokenize{(40,54,4)x1}}
\item \texttt{\detokenize{(80,54,6)x1}}
\end{itemize}
\noindent\textbf{Conclusión en la fibra.} Operador diagonal exacto sobre la fibra \(B_1\); la ruta nominal y el morfismo declarado publican el carácter escalar sin sustituir el espectro de la fibra.
\medskip
\Needspace{10\baselineskip}
\noindent\textbf{rama quark de tipo inferior, profundidad G2.} \(\mathcal B_{d,\mathrm{G2}}\).\par
\noindent\textbf{Dominio de fibra.} Rango 4; celdas 30,34,48,52; B | G | R; 4 familias de ruta.
\noindent\textbf{Espectro completo de \(K_D\).}
\begin{itemize}
\item \texttt{\detokenize{(0,24,0)x1}}
\item \texttt{\detokenize{(40,84,30)x1}}
\item \texttt{\detokenize{(60,78,25)x1}}
\item \texttt{\detokenize{(80,66,2)x1}}
\end{itemize}
\noindent\textbf{Espectro completo de \(K_\Omega\).}
\begin{itemize}
\item \texttt{\detokenize{(4,54,2)x1}}
\item \texttt{\detokenize{(80,54,6)x3}}
\end{itemize}
\noindent\textbf{Conclusión en la fibra.} Operador diagonal exacto sobre la fibra \(B_1\); la ruta nominal y el morfismo declarado publican el carácter escalar sin sustituir el espectro de la fibra.
\medskip
\Needspace{10\baselineskip}
\noindent\textbf{rama quark de tipo inferior, profundidad G3.} \(\mathcal B_{d,\mathrm{G3}}\).\par
\noindent\textbf{Dominio de fibra.} Rango 6; celdas 12,14,16,66,68,70; B | G | R; 6 familias de ruta.
\noindent\textbf{Espectro completo de \(K_D\).}
\begin{itemize}
\item \texttt{\detokenize{(40,24,2)x2}}
\item \texttt{\detokenize{(40,78,27)x1}}
\item \texttt{\detokenize{(60,84,28)x1}}
\item \texttt{\detokenize{(80,54,6)x1}}
\item \texttt{\detokenize{(80,66,4)x1}}
\end{itemize}
\noindent\textbf{Espectro completo de \(K_\Omega\).}
\begin{itemize}
\item \texttt{\detokenize{(36,54,4)x1}}
\item \texttt{\detokenize{(80,54,6)x5}}
\end{itemize}
\noindent\textbf{Conclusión en la fibra.} Operador diagonal exacto sobre la fibra \(B_1\); la ruta nominal y el morfismo declarado publican el carácter escalar sin sustituir el espectro de la fibra.
\medskip
\Needspace{10\baselineskip}
\noindent\textbf{rama quark de tipo inferior, profundidad G4.} \(\mathcal B_{d,\mathrm{G4}}\).\par
\noindent\textbf{Dominio de fibra.} Rango 8; celdas 10,18,28,36,46,54,64,72; B | G | R; 8 familias de ruta.
\noindent\textbf{Espectro completo de \(K_D\).}
\begin{itemize}
\item \texttt{\detokenize{(0,24,0)x1}}
\item \texttt{\detokenize{(40,24,2)x1}}
\item \texttt{\detokenize{(40,78,27)x2}}
\item \texttt{\detokenize{(40,84,30)x1}}
\item \texttt{\detokenize{(80,54,7)x1}}
\item \texttt{\detokenize{(80,66,3)x1}}
\item \texttt{\detokenize{(100,66,2)x1}}
\end{itemize}
\noindent\textbf{Espectro completo de \(K_\Omega\).}
\begin{itemize}
\item \texttt{\detokenize{(36,54,4)x1}}
\item \texttt{\detokenize{(80,54,6)x7}}
\end{itemize}
\noindent\textbf{Conclusión en la fibra.} Operador diagonal exacto sobre la fibra \(B_1\); la ruta nominal y el morfismo declarado publican el carácter escalar sin sustituir el espectro de la fibra.
\medskip
\Needspace{10\baselineskip}
\noindent\textbf{rama quark de tipo superior, profundidad G1.} \(\mathcal B_{u,\mathrm{G1}}\).\par
\noindent\textbf{Dominio de fibra.} Rango 4; celdas 31,33,49,51; B | G | R; 3 familias de ruta.
\noindent\textbf{Espectro completo de \(K_D\).}
\begin{itemize}
\item \texttt{\detokenize{(40,54,6)x1}}
\item \texttt{\detokenize{(60,78,25)x1}}
\item \texttt{\detokenize{(80,66,2)x1}}
\item \texttt{\detokenize{(80,66,3)x1}}
\end{itemize}
\noindent\textbf{Espectro completo de \(K_\Omega\).}
\begin{itemize}
\item \texttt{\detokenize{(80,54,6)x4}}
\end{itemize}
\noindent\textbf{Conclusión en la fibra.} Operador diagonal exacto sobre la fibra \(B_1\); la ruta nominal y el morfismo declarado publican el carácter escalar sin sustituir el espectro de la fibra.
\medskip
\Needspace{29\baselineskip}
\noindent\textbf{rama quark de tipo superior, profundidad G3.} \(\mathcal B_{u,\mathrm{G3}}\).\par
\noindent\textbf{Dominio de fibra.} Rango 12; celdas 11,13,15,17,29,35,47,53,65,67,69,71; B | G | R; 12 familias de ruta.
\noindent\textbf{Espectro completo de \(K_D\).}
\begin{itemize}
\item \texttt{\detokenize{(40,24,2)x2}}
\item \texttt{\detokenize{(40,78,27)x2}}
\item \texttt{\detokenize{(60,84,28)x1}}
\item \texttt{\detokenize{(80,54,6)x3}}
\item \texttt{\detokenize{(80,66,3)x1}}
\item \texttt{\detokenize{(80,66,4)x1}}
\item \texttt{\detokenize{(100,66,0)x1}}
\item \texttt{\detokenize{(100,66,2)x1}}
\end{itemize}
\noindent\textbf{Espectro completo de \(K_\Omega\).}
\begin{itemize}
\item \texttt{\detokenize{(56,54,5)x3}}
\item \texttt{\detokenize{(80,54,6)x9}}
\end{itemize}
\noindent\textbf{Conclusión en la fibra.} Operador diagonal exacto sobre la fibra \(B_1\); la ruta nominal y el morfismo declarado publican el carácter escalar sin sustituir el espectro de la fibra.
\medskip
\endgroup


\Needspace{15\baselineskip}
\subsubsection{Morfismos de realización física}

\begingroup
\setlength{\tabcolsep}{4.5pt}
\renewcommand{\arraystretch}{1.28}
\begin{longtable}{@{}>{\raggedright\arraybackslash}p{.14\textwidth}>{\raggedright\arraybackslash}p{.33\textwidth}>{\raggedright\arraybackslash}p{.25\textwidth}>{\raggedright\arraybackslash}p{.22\textwidth}@{}}
\caption{Realización física de los leptones cargados.}\label{tab:masas-realizacion-leptones}\\
\toprule
\textbf{Clase} & \textbf{Objeto HMT} & \textbf{Mapa de realización} & \textbf{Conclusión y condiciones} \\
\midrule
\endfirsthead
\toprule
\textbf{Clase} & \textbf{Objeto HMT} & \textbf{Mapa de realización} & \textbf{Conclusión y condiciones} \\
\midrule
\endhead
\midrule
\multicolumn{4}{r}{\footnotesize Continúa en la página siguiente.}\\
\endfoot
\bottomrule
\endlastfoot
$e^-/e^+$ & Celda central $(5,5)$, ruta $\Gamma_e=(5,5,++)$ y $K_D(\Gamma_e)=K_\Omega(\Gamma_e)=(80,54,6)$. & Reducción de rango uno cerrada. & Clausura beta de rango uno; selector independiente de magnitudes metrológicas empleadas como blanco. \\
$\mu^-/\mu^+$ & familia leptónica cargada; niveles G2; 8 celdas; $K_D$: 7 perfiles; $K_\Omega$: 1 perfil & Generación y firma fina $\to$ ruta persistente $\to\chi_{\rm mass}\to\mathcal L_M$. & Ruta nominal y carácter escalar cerrados; el par operatorial conserva el refinamiento completo de la fibra. \\
$\tau^-/\tau^+$ & familia leptónica cargada; niveles G4; 16 celdas; $K_D$: 8 perfiles; $K_\Omega$: 3 perfiles & Generación y firma fina $\to$ ruta persistente $\to\chi_{\rm mass}\to\mathcal L_M$. & Ruta nominal y carácter escalar cerrados; el par operatorial conserva el refinamiento completo de la fibra. \\
\end{longtable}
\endgroup

\begingroup
\setlength{\tabcolsep}{4.5pt}
\renewcommand{\arraystretch}{1.28}
\begin{longtable}{@{}>{\raggedright\arraybackslash}p{.14\textwidth}>{\raggedright\arraybackslash}p{.33\textwidth}>{\raggedright\arraybackslash}p{.25\textwidth}>{\raggedright\arraybackslash}p{.22\textwidth}@{}}
\caption{Realización física de los neutrinos y transporte de bases.}\label{tab:masas-realizacion-neutrinos}\\
\toprule
\textbf{Clase} & \textbf{Objeto HMT} & \textbf{Mapa de realización} & \textbf{Conclusión y condiciones} \\
\midrule
\endfirsthead
\toprule
\textbf{Clase} & \textbf{Objeto HMT} & \textbf{Mapa de realización} & \textbf{Conclusión y condiciones} \\
\midrule
\endhead
\midrule
\multicolumn{4}{r}{\footnotesize Continúa en la página siguiente.}\\
\endfoot
\bottomrule
\endlastfoot
$\nu_e/\bar\nu_e$ & familia neutra; niveles G1, G2, G3, G4; 20 celdas; $K_D$: 11 perfiles; $K_\Omega$: 5 perfiles & Registro neutro de rango pleno $\to$ autovectores $\to$ base PMNS. & La estructura exterior/CAR fija la fermionicidad; PMNS transporta bases y no crea autovalores de masa. \\
$\nu_\mu/\bar\nu_\mu$ & familia neutra; niveles G1, G2, G3, G4; 20 celdas; $K_D$: 11 perfiles; $K_\Omega$: 5 perfiles & Registro neutro de rango pleno $\to$ autovectores $\to$ base PMNS. & La estructura exterior/CAR fija la fermionicidad; PMNS transporta bases y no crea autovalores de masa. \\
$\nu_\tau/\bar\nu_\tau$ & familia neutra; niveles G1, G2, G3, G4; 20 celdas; $K_D$: 11 perfiles; $K_\Omega$: 5 perfiles & Registro neutro de rango pleno $\to$ autovectores $\to$ base PMNS. & La estructura exterior/CAR fija la fermionicidad; PMNS transporta bases y no crea autovalores de masa. \\
\end{longtable}
\endgroup

\begingroup
\setlength{\tabcolsep}{4.5pt}
\renewcommand{\arraystretch}{1.28}
\begin{longtable}{@{}>{\raggedright\arraybackslash}p{.14\textwidth}>{\raggedright\arraybackslash}p{.33\textwidth}>{\raggedright\arraybackslash}p{.25\textwidth}>{\raggedright\arraybackslash}p{.22\textwidth}@{}}
\caption{Realización física de los quarks: generación, sabor y color.}\label{tab:masas-realizacion-quarks}\\
\toprule
\textbf{Clase} & \textbf{Objeto HMT} & \textbf{Mapa de realización} & \textbf{Conclusión y condiciones} \\
\midrule
\endfirsthead
\toprule
\textbf{Clase} & \textbf{Objeto HMT} & \textbf{Mapa de realización} & \textbf{Conclusión y condiciones} \\
\midrule
\endhead
\midrule
\multicolumn{4}{r}{\footnotesize Continúa en la página siguiente.}\\
\endfoot
\bottomrule
\endlastfoot
$u/\bar u$ & rama quark \emph{up}; niveles G1 y G3; color $\mathbb Z/3$; 16 celdas; $K_D$: 11 perfiles; $K_\Omega$: 2 perfiles & Sabor--color--generación $\to$ ruta persistente $\to\chi_{\rm mass}\to\mathcal L_M$. & Sabor, color y generación determinan la clase de ruta; carácter escalar y par operatorial cerrados. \\
$d/\bar d$ & rama quark \emph{down}; niveles G1, G2, G3, G4; color $\mathbb Z/3$; 20 celdas; $K_D$: 12 perfiles; $K_\Omega$: 4 perfiles & Sabor--color--generación $\to$ ruta persistente $\to\chi_{\rm mass}\to\mathcal L_M$. & Sabor, color y generación determinan la clase de ruta; carácter escalar y par operatorial cerrados. \\
$c/\bar c$ & rama quark \emph{up}; niveles G1 y G3; color $\mathbb Z/3$; 16 celdas; $K_D$: 11 perfiles; $K_\Omega$: 2 perfiles & Sabor--color--generación $\to$ ruta persistente $\to\chi_{\rm mass}\to\mathcal L_M$. & Sabor, color y generación determinan la clase de ruta; carácter escalar y par operatorial cerrados. \\
$s/\bar s$ & rama quark \emph{down}; niveles G1, G2, G3, G4; color $\mathbb Z/3$; 20 celdas; $K_D$: 12 perfiles; $K_\Omega$: 4 perfiles & Sabor--color--generación $\to$ ruta persistente $\to\chi_{\rm mass}\to\mathcal L_M$. & Sabor, color y generación determinan la clase de ruta; carácter escalar y par operatorial cerrados. \\
$t/\bar t$ & rama quark \emph{up}; niveles G1 y G3; color $\mathbb Z/3$; 16 celdas; $K_D$: 11 perfiles; $K_\Omega$: 2 perfiles & Sabor--color--generación $\to$ ruta persistente $\to\chi_{\rm mass}\to\mathcal L_M$. & Sabor, color y generación determinan la clase de ruta; carácter escalar y par operatorial cerrados. \\
$b/\bar b$ & rama quark \emph{down}; niveles G1, G2, G3, G4; color $\mathbb Z/3$; 20 celdas; $K_D$: 12 perfiles; $K_\Omega$: 4 perfiles & Sabor--color--generación $\to$ ruta persistente $\to\chi_{\rm mass}\to\mathcal L_M$. & Sabor, color y generación determinan la clase de ruta; carácter escalar y par operatorial cerrados. \\
\end{longtable}
\endgroup

\Needspace{12\baselineskip}
\paragraph{Localización estructural de los \(6\) sabores de quark}
\begingroup
\setlength{\parskip}{.3\baselineskip}
La posición de un sabor quark se especifica mediante cuatro datos simultáneos: rama, generación, soporte celular y familia de rutas orientadas. Las tres generaciones de una misma rama comparten el soporte del atlas; el clasificador fino de generación conserva la distinción física que todavía debe realizarse como autovector y sección de masa.
\begingroup
\setlength{\tabcolsep}{4.5pt}
\renewcommand{\arraystretch}{1.28}
\begin{longtable}{@{}>{\raggedright\arraybackslash}p{.12\textwidth}>{\raggedright\arraybackslash}p{.17\textwidth}>{\raggedright\arraybackslash}p{.19\textwidth}>{\raggedright\arraybackslash}p{.46\textwidth}@{}}
\caption{Clasificación generacional de los seis sabores quark.}\label{tab:masas-posicion-quarks-seis-sabores}\\
\toprule
\textbf{Sabor} & \textbf{Generación} & \textbf{Rama} & \textbf{Clasificador de rama y generación} \\
\midrule
\endfirsthead
\toprule
\textbf{Sabor} & \textbf{Generación} & \textbf{Rama} & \textbf{Clasificador de rama y generación} \\
\midrule
\endhead
\midrule
\multicolumn{4}{r}{\footnotesize Continúa en la página siguiente.}\\
\endfoot
\bottomrule
\endlastfoot
$u/\bar u$ & primera & superior & $(\operatorname{up},1,\chi_{\rm fina})$; par nominal de ruta $(21,30)$ sobre el eje $\mathrm{N\!-\!S}$ \\
$d/\bar d$ & primera & inferior & $(\operatorname{down},1,\chi_{\rm fina})$; par nominal de ruta $(20,32)$ sobre el eje $\mathrm{E\!-\!O}$ \\
$c/\bar c$ & segunda & superior & $(\operatorname{up},2,\chi_{\rm fina})$; par nominal de ruta $(21,29)$ sobre el eje $\mathrm{E\!-\!O}$ \\
$s/\bar s$ & segunda & inferior & $(\operatorname{down},2,\chi_{\rm fina})$; par nominal de ruta $(17,27)$ sobre el eje $\mathrm{N\!-\!S}$ \\
$t/\bar t$ & tercera & superior & $(\operatorname{up},3,\chi_{\rm fina})$; par nominal de ruta $(21,26)$ sobre el eje $\mathrm{N\!-\!S}$ \\
$b/\bar b$ & tercera & inferior & $(\operatorname{down},3,\chi_{\rm fina})$; par nominal de ruta $(21,30)$ sobre el eje $\mathrm{E\!-\!O}$ \\
\end{longtable}
\endgroup

\Needspace{18\baselineskip}
\noindent\textbf{Soporte estructural de la rama superior.}\par
\noindent\textit{Multisecciones.} \(\mathcal B_{u,\mathrm{G1}}\), \(\mathcal B_{u,\mathrm{G3}}\).\par
\noindent\textit{Profundidades y cardinales.} $\mathrm{G1}$, $\mathrm{G3}$; 16 celdas y 14 familias de ruta distintas.\par
\noindent\textit{Órbita cromática y orientaciones.} $\mathcal C=\{B,G,R\}$ y $\mathcal O=\{(-1,-1),(-1,1),(1,-1),(1,1)\}$.\par
\noindent\textit{Celdas y familias orientadas $B_1$.} Cada cadena $k\leftrightarrow(r,c)\leftrightarrow\Gamma_{r,c}^{(h,v)}$ conserva el índice celular, su coordenada en la carta $9\times9$ y los dos signos de orientación del registro de ruta enriquecido.\par
\noindent Fila $r=2$: $11\leftrightarrow(2,2)\leftrightarrow\Gamma_{2,2}^{(-,+)}$, $13\leftrightarrow(2,4)\leftrightarrow\Gamma_{2,4}^{(+,-)}$, $15\leftrightarrow(2,6)\leftrightarrow\Gamma_{2,6}^{(-,+)}$, $17\leftrightarrow(2,8)\leftrightarrow\Gamma_{2,8}^{(-,-)}$.\par
\noindent Fila $r=4$: $29\leftrightarrow(4,2)\leftrightarrow\Gamma_{4,2}^{(+,-)}$, $31\leftrightarrow(4,4)\leftrightarrow\Gamma_{4,4}^{(-,+)}$, $33\leftrightarrow(4,6)\leftrightarrow\Gamma_{4,6}^{(+,+)}$, $35\leftrightarrow(4,8)\leftrightarrow\Gamma_{4,8}^{(-,-)}$.\par
\noindent Fila $r=6$: $47\leftrightarrow(6,2)\leftrightarrow\Gamma_{6,2}^{(-,+)}$, $49\leftrightarrow(6,4)\leftrightarrow\Gamma_{6,4}^{(-,-)}$, $51\leftrightarrow(6,6)\leftrightarrow\Gamma_{6,6}^{(+,+)}$, $53\leftrightarrow(6,8)\leftrightarrow\Gamma_{6,8}^{(+,+)}$.\par
\noindent Fila $r=8$: $65\leftrightarrow(8,2)\leftrightarrow\Gamma_{8,2}^{(+,+)}$, $67\leftrightarrow(8,4)\leftrightarrow\Gamma_{8,4}^{(+,+)}$, $69\leftrightarrow(8,6)\leftrightarrow\Gamma_{8,6}^{(-,-)}$, $71\leftrightarrow(8,8)\leftrightarrow\Gamma_{8,8}^{(+,+)}$.\par
\noindent\textit{Morfismo de realización cerrado.} Para $g\in\{1,2,3\}$, la rama, la generación, la órbita de color y la firma fina determinan
\[\chi_{\rm fina}(\operatorname{up},g,\mathcal C)\longrightarrow\Gamma_{\rm enr}\longrightarrow\bigl(\operatorname{quot}_{\mathcal C},S(B_D,B_\Omega)\bigr).\]
Esta flecha selecciona el autovector físico y publica después la coordenada escalar en $\mathcal L_M$ sin usar el valor de contraste.\par
\medskip
\Needspace{18\baselineskip}
\noindent\textbf{Soporte estructural de la rama inferior.}\par
\noindent\textit{Multisecciones.} \(\mathcal B_{d,\mathrm{G1}}\), \(\mathcal B_{d,\mathrm{G2}}\), \(\mathcal B_{d,\mathrm{G3}}\), \(\mathcal B_{d,\mathrm{G4}}\).\par
\noindent\textit{Profundidades y cardinales.} $\mathrm{G1}$, $\mathrm{G2}$, $\mathrm{G3}$, $\mathrm{G4}$; 20 celdas y 18 familias de ruta distintas.\par
\noindent\textit{Órbita cromática y orientaciones.} $\mathcal C=\{B,G,R\}$ y $\mathcal O=\{(-1,-1),(-1,1),(1,-1),(1,1)\}$.\par
\noindent\textit{Celdas y familias orientadas $B_1$.} Cada cadena $k\leftrightarrow(r,c)\leftrightarrow\Gamma_{r,c}^{(h,v)}$ conserva el índice celular, su coordenada en la carta $9\times9$ y los dos signos de orientación del registro de ruta enriquecido.\par
\noindent Fila $r=2$: $10\leftrightarrow(2,1)\leftrightarrow\Gamma_{2,1}^{(+,+)}$, $12\leftrightarrow(2,3)\leftrightarrow\Gamma_{2,3}^{(-,-)}$, $14\leftrightarrow(2,5)\leftrightarrow\Gamma_{2,5}^{(+,+)}$, $16\leftrightarrow(2,7)\leftrightarrow\Gamma_{2,7}^{(-,+)}$, $18\leftrightarrow(2,9)\leftrightarrow\Gamma_{2,9}^{(-,+)}$.\par
\noindent Fila $r=4$: $28\leftrightarrow(4,1)\leftrightarrow\Gamma_{4,1}^{(-,-)}$, $30\leftrightarrow(4,3)\leftrightarrow\Gamma_{4,3}^{(+,+)}$, $32\leftrightarrow(4,5)\leftrightarrow\Gamma_{4,5}^{(+,-)}$, $34\leftrightarrow(4,7)\leftrightarrow\Gamma_{4,7}^{(-,+)}$, $36\leftrightarrow(4,9)\leftrightarrow\Gamma_{4,9}^{(-,+)}$.\par
\noindent Fila $r=6$: $46\leftrightarrow(6,1)\leftrightarrow\Gamma_{6,1}^{(-,-)}$, $48\leftrightarrow(6,3)\leftrightarrow\Gamma_{6,3}^{(-,+)}$, $50\leftrightarrow(6,5)\leftrightarrow\Gamma_{6,5}^{(-,+)}$, $52\leftrightarrow(6,7)\leftrightarrow\Gamma_{6,7}^{(-,+)}$, $54\leftrightarrow(6,9)\leftrightarrow\Gamma_{6,9}^{(+,-)}$.\par
\noindent Fila $r=8$: $64\leftrightarrow(8,1)\leftrightarrow\Gamma_{8,1}^{(+,-)}$, $66\leftrightarrow(8,3)\leftrightarrow\Gamma_{8,3}^{(-,+)}$, $68\leftrightarrow(8,5)\leftrightarrow\Gamma_{8,5}^{(-,+)}$, $70\leftrightarrow(8,7)\leftrightarrow\Gamma_{8,7}^{(+,-)}$, $72\leftrightarrow(8,9)\leftrightarrow\Gamma_{8,9}^{(-,+)}$.\par
\noindent\textit{Morfismo de realización cerrado.} Para $g\in\{1,2,3\}$, la rama, la generación, la órbita de color y la firma fina determinan
\[\chi_{\rm fina}(\operatorname{down},g,\mathcal C)\longrightarrow\Gamma_{\rm enr}\longrightarrow\bigl(\operatorname{quot}_{\mathcal C},S(B_D,B_\Omega)\bigr).\]
Esta flecha selecciona el autovector físico y publica después la coordenada escalar en $\mathcal L_M$ sin usar el valor de contraste.\par
\medskip
\endgroup

\begingroup
\setlength{\tabcolsep}{4.5pt}
\renewcommand{\arraystretch}{1.28}
\begin{longtable}{@{}>{\raggedright\arraybackslash}p{.14\textwidth}>{\raggedright\arraybackslash}p{.33\textwidth}>{\raggedright\arraybackslash}p{.25\textwidth}>{\raggedright\arraybackslash}p{.22\textwidth}@{}}
\caption{Realización física de los sectores gauge y escalar.}\label{tab:masas-realizacion-gauge-escalar}\\
\toprule
\textbf{Clase} & \textbf{Objeto HMT} & \textbf{Mapa de realización} & \textbf{Conclusión y condiciones} \\
\midrule
\endfirsthead
\toprule
\textbf{Clase} & \textbf{Objeto HMT} & \textbf{Mapa de realización} & \textbf{Conclusión y condiciones} \\
\midrule
\endhead
\midrule
\multicolumn{4}{r}{\footnotesize Continúa en la página siguiente.}\\
\endfoot
\bottomrule
\endlastfoot
fotón $\gamma$ & Carta nula sectorial y representación gauge correspondiente. & Inclusión exacta de la sección nula en la recta de masa. & Nulidad exacta en su carta sectorial. \\
gluones $g$ & Carta nula sectorial y representación gauge correspondiente. & Inclusión exacta de la sección nula en la recta de masa. & Nulidad exacta en su carta sectorial. \\
$W^\pm$ & ruta CT--108 W: 15/37, eje E-O, $(N,E,S,O)=(23,44,11,30)$, hoja/residuo $(41,8)$, $(n_A,s_C)=(94,-1)$, $(W_+,W_-)=(563,565)$; carácter escalar exacto & Ruta CT--108 nominal $\to(n_A,s_C)\to(W_+,W_-)\to\chi_{\rm mass}\to\mathcal L_M$. & Ruta gauge CT--108, firma y carácter escalar determinados antes del contraste. \\
$Z^0$ & ruta CT--108 Z: 20/23, eje E-O, $(N,E,S,O)=(33,45,10,20)$, hoja/residuo $(57,7)$, $(n_A,s_C)=(95,-1)$, $(W_+,W_-)=(569,571)$; carácter escalar exacto & Ruta CT--108 nominal $\to(n_A,s_C)\to(W_+,W_-)\to\chi_{\rm mass}\to\mathcal L_M$. & Ruta gauge CT--108, firma y carácter escalar determinados antes del contraste. \\
Higgs $H$ & ruta CT--108 H: 24/29, eje N-S, $(N,E,S,O)=(40,34,22,12)$, hoja/residuo $(47,10)$, $(n_A,s_C)=(97,9)$, $(W_+,W_-)=(591,573)$; carácter escalar exacto & Ruta CT--108 nominal $\to(n_A,s_C)\to(W_+,W_-)\to\chi_{\rm mass}\to\mathcal L_M$. & Ruta escalar CT--108, firma y carácter determinados antes del contraste. \\
\end{longtable}
\endgroup

\begingroup
\setlength{\tabcolsep}{4.5pt}
\renewcommand{\arraystretch}{1.28}
\begin{longtable}{@{}>{\raggedright\arraybackslash}p{.14\textwidth}>{\raggedright\arraybackslash}p{.33\textwidth}>{\raggedright\arraybackslash}p{.25\textwidth}>{\raggedright\arraybackslash}p{.22\textwidth}@{}}
\caption{Realización de compuestos, sectores topológicos y estados virtuales.}\label{tab:masas-realizacion-no-escalar}\\
\toprule
\textbf{Clase} & \textbf{Objeto HMT} & \textbf{Mapa de realización} & \textbf{Conclusión y condiciones} \\
\midrule
\endfirsthead
\toprule
\textbf{Clase} & \textbf{Objeto HMT} & \textbf{Mapa de realización} & \textbf{Conclusión y condiciones} \\
\midrule
\endhead
\midrule
\multicolumn{4}{r}{\footnotesize Continúa en la página siguiente.}\\
\endfoot
\bottomrule
\endlastfoot
mesones & Registro doble, frontera de enlace y firma compuesta. & Registro compuesto y ligadura $\to$ carta escalar. & Composición de registros; la ligadura precede a la carta escalar. \\
bariones & Registro triple, acarreo y firma compuesta. & Registro compuesto y ligadura $\to$ carta escalar. & Composición de registros; la ligadura precede a la carta escalar. \\
Fibonacci $(1,\tau)$ & Anillo de fusión Fibonacci: $\tau\otimes\tau=1+\tau$. & Fusión y trenzado $\to$ observable físico del sector. & Codominio exacto de fusión y trenzado; el contraste usa sus invariantes propios. \\
Ising $(1,\psi,\sigma)$; realización Majorana separada & Anillo de fusión de Ising: $(1,\psi,\sigma)$. & Fusión y trenzado $\to$ observable físico del sector. & Codominio exacto de fusión y trenzado; el contraste usa sus invariantes propios. \\
sectores $\mathbb Z/6$ & Carácter de trenzado sobre $\mathbb Z/6$. & Fusión y trenzado $\to$ observable físico del sector. & Codominio exacto de fusión y trenzado; el contraste usa sus invariantes propios. \\
secciones virtuales & Sistema inverso de extensiones y horizonte de persistencia. & Prolongación registrada $\to$ horizonte físico de persistencia. & Objeto exacto de prolongación y obstrucción; su salida es el horizonte de persistencia. \\
\end{longtable}
\endgroup


\Needspace{15\baselineskip}
\subsubsection{Comparabilidad física en los codominios nativos}

La comparación externa se define sobre el codominio nativo de cada clase. En las
clases con sección escalar cerrada se contrastan magnitudes homogéneas, con unidad,
esquema y escala declarados. En las fibras operatoriales se contrasta el espectro
interno y se conserva de forma explícita el morfismo de realización cerrado que
selecciona la ruta sin usar el valor externo. Los registros mesónicos y bariónicos sólo admiten comparación estado por
estado después de incorporar la frontera de ligadura. En el sector Fibonacci están
cerrados el anillo basado y \(\operatorname{FPdim}(\tau)=\varphi\); las matrices
\(F/R\), el pentágono, el hexágono, el Hamiltoniano y la brecha pertenecen a la
realización posterior. Ising y \(\mathbb Z/6\) conservan por separado sus invariantes
de fusión y trenzado, y la virtualidad se compara mediante persistencia y estructura
analítica. En ninguno de estos seis
casos existe una masa escalar universal de clase; asignarla sustituiría el objeto
demostrado por una cifra sin morfismo.

\endgroup
\begingroup
\let\chapter\section
\let\section\subsection
\let\subsection\subsubsection
\let\subsubsection\paragraph
\let\clearpage\relax

\Needspace{15\baselineskip}
\subsubsection{Evaluaciones internas de firmas explícitas y contraste posterior}

La tabla siguiente publica de forma legible las diecisiete coordenadas HMT--MD ya
calculadas ---once elementales y seis compuestas--- y las confronta con la
edición 2026 del conjunto de datos del Particle Data Group, con fecha de corte
editorial de 15 de enero de 2026 \cite{PDG2026}. No constituyen otra ontología
distinta de la capa operatoria anterior. El símbolo \(\dagger\)
identifica una evaluación obtenida al aplicar el mapa de construcción a una firma
explícita. En las once filas elementales la firma procede de la ruta nominal
determinada sin masas; en las seis compuestas procede del registro de ligadura
declarado. Cada cifra depende de esa firma, del lector bidireccional y de la carta
dimensional común; el conjunto de datos externo interviene sólo en la comparación
y no selecciona sabor, generación, color ni ruta.

\enlargethispage{2\baselineskip}
Para leptones cargados y bosones masivos se muestra la diferencia relativa entre
las dos coordenadas centrales y, cuando el conjunto de datos aporta una incertidumbre
bilateral, su cociente por esa incertidumbre. En los quarks se publica la
diferencia aritmética entre valores centrales, pero se reserva el término
\emph{residual metrológico}: éste exige que la masa interna y la referencia externa
usen la misma definición, el mismo esquema de renormalización y la misma escala.
Las masas ligeras se leen en la convención \(\overline{\mathrm{MS}}\) a
\(2\,\mathrm{GeV}\); las de \(c\) y \(b\), en
\(\mu=\overline m_q\); la masa superior procede de reconstrucción cinemática.

\begingroup
\setlength{\tabcolsep}{4.5pt}
\renewcommand{\arraystretch}{1.28}
\begin{longtable}{@{}>{\raggedright\arraybackslash}p{.12\textwidth}>{\raggedright\arraybackslash}p{.27\textwidth}>{\raggedright\arraybackslash}p{.27\textwidth}>{\raggedright\arraybackslash}p{.28\textwidth}@{}}
\caption{Evaluaciones internas de firmas explícitas y contraste metrológico posterior.}\label{tab:masas-evaluaciones-firmas-explicitas}\\
\toprule
\textbf{Estado} & \textbf{Coordenada de la firma} & \textbf{Referencia física externa} & \textbf{Alcance} \\
\midrule
\endfirsthead
\toprule
\textbf{Estado} & \textbf{Coordenada de la firma} & \textbf{Referencia física externa} & \textbf{Alcance} \\
\midrule
\endhead
\midrule
\multicolumn{4}{r}{\footnotesize Continúa en la página siguiente.}\\
\endfoot
\bottomrule
\endlastfoot
% HMT-MASS-EXPLICIT-SIGNATURE:SM-LC-2
$\mu^-/\mu^+$ & $105.658360\allowbreak{}294435\allowbreak{}829303\allowbreak{}4\;\mathrm{MeV}/c^2$\textsuperscript{\(\dagger\)}; torre E3, firma (42,-3,8,0,2) & $105.6583755\pm0.0000023\;\mathrm{MeV}/c^2$; masa leptónica de reposo; promedio PDG y conversión metrológica declarada & Diferencia relativa: $-0.143912\allowbreak{}530\;\mathrm{ppm}$\textsuperscript{\(\dagger\)}; $-6.611114\,\sigma$. La cifra HMT es la evaluación exacta de la firma declarada. \\
% HMT-MASS-EXPLICIT-SIGNATURE:SM-LC-3
$\tau^-/\tau^+$ & $1776.860917\allowbreak{}631432\allowbreak{}295595\;\mathrm{MeV}/c^2$\textsuperscript{\(\dagger\)}; torre E3, firma (64,0,25,-1,6) & $1776.93\pm0.09\;\mathrm{MeV}/c^2$; promedio directo PDG & Diferencia relativa: $-38.877371\allowbreak{}966\;\mathrm{ppm}$\textsuperscript{\(\dagger\)}; $-0.813547\,\sigma$. La cifra HMT es la evaluación exacta de la firma declarada. \\
% HMT-MASS-EXPLICIT-SIGNATURE:SM-Q-U1
$u/\bar u$ & $2.165924\allowbreak{}536173\allowbreak{}907\;\mathrm{MeV}/c^2$\textsuperscript{\(\dagger\)}; carta angular, firma nominal de ruta (11,7) & $2.16\pm0.07\;\mathrm{MeV}/c^2$; masa $\overline{\mathrm{MS}}$ a $\mu=2\,\mathrm{GeV}$; 90\% C.L. & Diferencia entre valores centrales: $2742.840821\allowbreak{}253\;\mathrm{ppm}$\textsuperscript{\(\dagger\)}. No es un residual metrológico hasta fijar el mismo esquema y la misma escala. \\
% HMT-MASS-EXPLICIT-SIGNATURE:SM-Q-D1
$d/\bar d$ & $4.679550\allowbreak{}162948\allowbreak{}874\;\mathrm{MeV}/c^2$\textsuperscript{\(\dagger\)}; carta angular, firma nominal de ruta (17,8) & $4.70\pm0.07\;\mathrm{MeV}/c^2$; masa $\overline{\mathrm{MS}}$ a $\mu=2\,\mathrm{GeV}$; 90\% C.L. & Diferencia entre valores centrales: $-4351.029159\allowbreak{}814\;\mathrm{ppm}$\textsuperscript{\(\dagger\)}. No es un residual metrológico hasta fijar el mismo esquema y la misma escala. \\
% HMT-MASS-EXPLICIT-SIGNATURE:SM-Q-U2
$c/\bar c$ & $1.270500\allowbreak{}838641\allowbreak{}911\;\mathrm{GeV}/c^2$\textsuperscript{\(\dagger\)}; carta angular, firma nominal de ruta (61,8) & $1.2729\pm0.0045\;\mathrm{GeV}/c^2$; masa $\overline{\mathrm{MS}}$ evaluada en $\mu=\overline{m}_q$; 90\% C.L. & Diferencia entre valores centrales: $-1884.799558\allowbreak{}558\;\mathrm{ppm}$\textsuperscript{\(\dagger\)}. No es un residual metrológico hasta fijar el mismo esquema y la misma escala. \\
% HMT-MASS-EXPLICIT-SIGNATURE:SM-Q-D2
$s/\bar s$ & $92.940093\allowbreak{}907845\allowbreak{}64\;\mathrm{MeV}/c^2$\textsuperscript{\(\dagger\)}; carta angular, firma nominal de ruta (41,-3) & $92.9\pm0.7\;\mathrm{MeV}/c^2$; masa $\overline{\mathrm{MS}}$ a $\mu=2\,\mathrm{GeV}$; 90\% C.L. & Diferencia entre valores centrales: $431.581354\allowbreak{}635\;\mathrm{ppm}$\textsuperscript{\(\dagger\)}. No es un residual metrológico hasta fijar el mismo esquema y la misma escala. \\
% HMT-MASS-EXPLICIT-SIGNATURE:SM-Q-U3
$t/\bar t$ & $172.597479\allowbreak{}544019\allowbreak{}1\;\mathrm{GeV}/c^2$\textsuperscript{\(\dagger\)}; carta angular, firma nominal de ruta (100,-1) & $172.60\pm0.27\;\mathrm{GeV}/c^2$; masa reconstruida directamente a partir de la cinemática de sucesos & Diferencia entre valores centrales: $-14.602873\allowbreak{}585\;\mathrm{ppm}$\textsuperscript{\(\dagger\)}. No es un residual metrológico hasta fijar el mismo esquema y la misma escala. \\
% HMT-MASS-EXPLICIT-SIGNATURE:SM-Q-D3
$b/\bar b$ & $4.190119\allowbreak{}763299\allowbreak{}790\;\mathrm{GeV}/c^2$\textsuperscript{\(\dagger\)}; carta angular, firma nominal de ruta (71,-5) & $4.186\pm0.006\;\mathrm{GeV}/c^2$; masa $\overline{\mathrm{MS}}$ evaluada en $\mu=\overline{m}_q$; 90\% C.L. & Diferencia entre valores centrales: $984.176612\allowbreak{}467\;\mathrm{ppm}$\textsuperscript{\(\dagger\)}. No es un residual metrológico hasta fijar el mismo esquema y la misma escala. \\
% HMT-MASS-EXPLICIT-SIGNATURE:SM-GA-W
$W^\pm$ & $80.378964\allowbreak{}909704\allowbreak{}547024\allowbreak{}86\;\mathrm{GeV}/c^2$\textsuperscript{\(\dagger\)}; torre E3, firma (94,-1,0,-2,4) & $80.3625\pm0.0077\;\mathrm{GeV}/c^2$; promedio de masa de la instantánea PDG sellada & Diferencia relativa: $204.882995\allowbreak{}234\;\mathrm{ppm}$\textsuperscript{\(\dagger\)}; $2.138299\,\sigma$. La cifra HMT es la evaluación exacta de la firma declarada. \\
% HMT-MASS-EXPLICIT-SIGNATURE:SM-GA-Z
$Z^0$ & $91.187533\allowbreak{}444313\allowbreak{}407844\allowbreak{}85\;\mathrm{GeV}/c^2$\textsuperscript{\(\dagger\)}; torre E3, firma (95,-1,-11,0,-4) & $91.1879\pm0.0020\;\mathrm{GeV}/c^2$; promedio de masa de la instantánea PDG sellada & Diferencia relativa: $-4.019784\allowbreak{}276\;\mathrm{ppm}$\textsuperscript{\(\dagger\)}; $-0.186362\,\sigma$. La cifra HMT es la evaluación exacta de la firma declarada. \\
% HMT-MASS-EXPLICIT-SIGNATURE:SM-H-1
Higgs $H$ & $125.292466\allowbreak{}042536\allowbreak{}133\;\mathrm{GeV}/c^2$\textsuperscript{\(\dagger\)}; carta angular, firma nominal de ruta (97,9) & $125.13\pm0.11\;\mathrm{GeV}/c^2$; promedio de masa de la instantánea PDG sellada & Diferencia relativa: $1298.378027\allowbreak{}140\;\mathrm{ppm}$\textsuperscript{\(\dagger\)}; $1.454206\,\sigma$. La cifra HMT es la evaluación exacta de la firma declarada. \\
% HMT-MASS-EXPLICIT-SIGNATURE:E3B-002
$\pi^0$\par firma $(44,\allowbreak -4,\allowbreak -25,\allowbreak 1,\allowbreak -2)$ & $134.976724\allowbreak{}327872\allowbreak{}125739\allowbreak{}384878\ldots\;\mathrm{MeV}/c^2$\textsuperscript{\(\dagger\)} & $134.976827\allowbreak{}767684\allowbreak{}71\pm 0.000485\allowbreak{}679282\allowbreak{}284233\allowbreak{}51\;\mathrm{MeV}/c^2$\par\footnotesize Catálogo PDG 2026. & Diferencia relativa $-0.766352\allowbreak{}375\,\mathrm{ppm}$; $-0.212979\,\sigma$. Para la firma compuesta declarada, la evaluación se deduce exactamente del carácter beta E3. El selector prospectivo de esa firma es un mapa anterior e independiente de la comparación numérica. \\
% HMT-MASS-EXPLICIT-SIGNATURE:E3B-003
$\pi^\pm$\par firma $(44,\allowbreak 1,\allowbreak -2,\allowbreak 2,\allowbreak -1)$ & $139.570485\allowbreak{}171572\allowbreak{}191374\allowbreak{}734364\ldots\;\mathrm{MeV}/c^2$\textsuperscript{\(\dagger\)} & $139.570390\allowbreak{}983681\allowbreak{}31\pm 0.000182\allowbreak{}007160\allowbreak{}408260\allowbreak{}09\;\mathrm{MeV}/c^2$\par\footnotesize Catálogo PDG 2026. & Diferencia relativa $0.674841\allowbreak{}491\,\mathrm{ppm}$; $0.517495\,\sigma$. Para la firma compuesta declarada, la evaluación se deduce exactamente del carácter beta E3. El selector prospectivo de esa firma es un mapa anterior e independiente de la comparación numérica. \\
% HMT-MASS-EXPLICIT-SIGNATURE:E3B-004
$K^\pm$\par firma $(54,\allowbreak -1,\allowbreak 17,\allowbreak -2,\allowbreak 2)$ & $493.677085\allowbreak{}649530\allowbreak{}612475\allowbreak{}723054\ldots\;\mathrm{MeV}/c^2$\textsuperscript{\(\dagger\)} & $493.676599\allowbreak{}458040\allowbreak{}58\pm 0.015405\allowbreak{}665226\allowbreak{}071509\;\mathrm{MeV}/c^2$\par\footnotesize Catálogo PDG 2026. & Diferencia relativa $0.984838\allowbreak{}030\,\mathrm{ppm}$; $0.031559\,\sigma$. Para la firma compuesta declarada, la evaluación se deduce exactamente del carácter beta E3. El selector prospectivo de esa firma es un mapa anterior e independiente de la comparación numérica. \\
% HMT-MASS-EXPLICIT-SIGNATURE:E3B-005
$K^0$\par firma $(54,\allowbreak 0,\allowbreak 32,\allowbreak 3,\allowbreak -2)$ & $497.611186\allowbreak{}497750\allowbreak{}457358\allowbreak{}558672\ldots\;\mathrm{MeV}/c^2$\textsuperscript{\(\dagger\)} & $497.610767\allowbreak{}531257\allowbreak{}52\pm 0.012990\allowbreak{}199756\allowbreak{}4242\;\mathrm{MeV}/c^2$\par\footnotesize Catálogo PDG 2026. & Diferencia relativa $0.841956\allowbreak{}244\,\mathrm{ppm}$; $0.032252\,\sigma$. Para la firma compuesta declarada, la evaluación se deduce exactamente del carácter beta E3. El selector prospectivo de esa firma es un mapa anterior e independiente de la comparación numérica. \\
% HMT-MASS-EXPLICIT-SIGNATURE:E3B-006
protón $p$\par firma $(59,\allowbreak 0,\allowbreak 9,\allowbreak 1,\allowbreak 0)$ & $938.271751\allowbreak{}546984\allowbreak{}898298\allowbreak{}936787\ldots\;\mathrm{MeV}/c^2$\textsuperscript{\(\dagger\)} & $938.272089\allowbreak{}430000\allowbreak{}05\pm 0.000000\allowbreak{}290000\allowbreak{}000000\allowbreak{}00\;\mathrm{MeV}/c^2$\par\footnotesize Catálogo PDG 2026. & Diferencia relativa $-0.360111\allowbreak{}974\,\mathrm{ppm}$; $-1165.113845\,\sigma$. Para la firma compuesta declarada, la evaluación se deduce exactamente del carácter beta E3. El selector prospectivo de esa firma es un mapa anterior e independiente de la comparación numérica. \\
% HMT-MASS-EXPLICIT-SIGNATURE:E3B-007
neutrón $n$\par firma $(59,\allowbreak 1,\allowbreak -34,\allowbreak 0,\allowbreak 1)$ & $939.565810\allowbreak{}472507\allowbreak{}023312\allowbreak{}664431\ldots\;\mathrm{MeV}/c^2$\textsuperscript{\(\dagger\)} & $939.565421\allowbreak{}939999\allowbreak{}96\pm 0.000000\allowbreak{}480000\allowbreak{}000000\allowbreak{}00\;\mathrm{MeV}/c^2$\par\footnotesize Catálogo PDG 2026. & Diferencia relativa $0.413523\allowbreak{}633\,\mathrm{ppm}$; $809.442723\,\sigma$. Para la firma compuesta declarada, la evaluación se deduce exactamente del carácter beta E3. El selector prospectivo de esa firma es un mapa anterior e independiente de la comparación numérica. \\
\end{longtable}
\endgroup


\Needspace{12\baselineskip}
\paragraph{Invariantes nativos de los sectores topológicos y virtuales.}
Los sectores Fibonacci, Ising, \(\mathbb Z/6\) y las secciones virtuales se comparan
en sus codominios propios. Fibonacci publica el anillo basado y su dimensión de
Frobenius--Perron; los datos \(F/R\) permanecen en la realización posterior. Ising
conserva su anillo de fusión y separa la realización Majorana; \(\mathbb Z/6\)
publica sus caracteres de trenzado; la virtualidad publica horizontes de
persistencia. Esta
tabla mantiene esos invariantes en el mismo mapa de realizaciones sin traducirlos
artificialmente a una unidad de masa.

\begingroup
\setlength{\tabcolsep}{4.5pt}
\renewcommand{\arraystretch}{1.28}
\begin{longtable}{@{}>{\raggedright\arraybackslash}p{.15\textwidth}>{\raggedright\arraybackslash}p{.31\textwidth}>{\raggedright\arraybackslash}p{.25\textwidth}>{\raggedright\arraybackslash}p{.23\textwidth}@{}}
\caption{Sectores topológicos y virtuales: invariantes nativos de fusión, trenzado y persistencia.}\label{tab:masas-invariantes-topologicos-virtuales-rev3}\\
\toprule
\textbf{Sector} & \textbf{Objeto y relaciones} & \textbf{Invariante nativo} & \textbf{Alcance físico} \\
\midrule
\endfirsthead
\toprule
\textbf{Sector} & \textbf{Objeto y relaciones} & \textbf{Invariante nativo} & \textbf{Alcance físico} \\
\midrule
\endhead
\midrule
\multicolumn{4}{r}{\footnotesize Continúa en la página siguiente.}\\
\endfoot
\bottomrule
\endlastfoot
Fibonacci & $K=\mathbb Z1\oplus\mathbb Z\tau$, $\tau^2=1+\tau$. & $\operatorname{FPdim}(\tau)=\varphi$ en el anillo basado. & Las matrices $F/R$, pentágono, hexágono, Hamiltoniano y brecha pertenecen a la realización posterior. \\
Ising & $K=\mathbb Z\{1,\psi,\sigma\}$, $\psi^2=1$, $\psi\sigma=\sigma$, $\sigma^2=1+\psi$. & $\operatorname{FPdim}(\sigma)=\sqrt{2}$; representación de trenzado. & La realización Majorana conserva su propio morfismo físico. \\
$\mathbb Z/6$ & $q_{\epsilon,t}=(-1)^\epsilon\omega_3^t=\zeta_6^k$. & Carácter cuadrático y representación de trenzado sobre los seis sectores. & La observación pertinente es topológica, no una masa universal. \\
Secciones virtuales & Sistema inverso de prolongaciones con horizonte $L(s_N)$. & Registro finito, firma, acarreo y clase de obstrucción. & El observable nativo es la persistencia de la prolongación. \\
\end{longtable}
\endgroup


Los neutrinos permanecen en el codominio de autovalores y mezcla descrito por sus
fichas operatoriales; PMNS realiza el transporte entre las bases de interacción y
masa, mientras la escala pertenece al espectro del operador neutrínico. Las seis
evaluaciones compuestas se obtienen de firmas explícitas mediante el morfismo
general de ligadura hacia la carta escalar. Los sectores topológicos y las secciones virtuales mantienen sus
observables propios de fusión, trenzado y persistencia.

\subsection{Construcción causal y cobertura de las clases de masa}

La arquitectura APP--TRIT--TPK, la ley de masas, el lector bidireccional y
las torres determinan conjuntamente el atlas, los registros cronológicos
enriquecidos, las firmas, el carácter y la clausura basal. Su confluencia en
\eqref{eq:ley-operatorial-general-masas-rectoras} proporciona una única
ecuación de secciones de masa. El dominio completo es el atlas \(81/56/13/324\) y cada realización especifica, según su codominio,
y especifica, según el codominio de cada sector, clasificadores, celdas, rutas,
orientaciones, palabras dodecafásicas, firmas y morfismos de compatibilidad. Los
espectros \(K_D/K_\Omega\) refinan las fibras del atlas, mientras las rutas
cronológicas de \(108\) estados asociadas a \(W\), \(Z\) y \(H\) construyen
directamente sus firmas y caracteres bosónicos.

\endgroup
\begingroup
\let\chapter\section
\let\section\subsection
\let\subsection\subsubsection
\let\subsubsection\paragraph
% Comparación trasladada detrás de la ley y de la ontología sectorial.
% Se excluyen únicamente la portada autónoma y el metalenguaje de edición.
\section{Evaluación cuantitativa posterior a la fijación del grafo generativo}
La tabla siguiente coloca en columnas contiguas toda coordenada escalar que la teoría permite evaluar de forma reproducible. La referencia experimental se abre después de fijar la salida HMT--MD. Una fila rotulada como evaluación condicionada afirma exactamente el valor de la función una vez fijada la firma; no atribuye a esa firma un selector físico que no haya sido construido. Los quarks se imprimen sin residual porque una comparación numérica requiere que esquema y escala de renormalización coincidan.

El censo numérico de esta tabla es exhaustivo respecto de las coordenadas escalares actualmente reproducibles: un electrón beta cerrado, diecisiete evaluaciones condicionadas ---once elementales y seis compuestas--- y dos salidas nulas exactas. Por tanto contiene veinte filas numéricas. Este número no es el dominio de la ley. Las clases cuya salida natural sigue siendo un operador, un espectro, una órbita o una multisección aparecen después con ese objeto completo, sin rellenar su columna con la referencia externa. El inventario experimental de 471 masas y 384 anchuras se conserva íntegro como contraste estado por estado; no se hace pasar su cardinal por un censo de predicciones HMT.

\Needspace{7\baselineskip}
\begingroup
\footnotesize
\setlength{\tabcolsep}{2.1pt}
\begin{longtable}{L{1.17cm}L{2.65cm}L{3.12cm}L{2.18cm}L{5.69cm}}
\toprule
\rowcolor{HMTNavy}\HMTtablehead{Estado} & \HMTtablehead{Valor HMT reproducible (MeV/\allowbreak{}c\({}^{2}\))} & \HMTtablehead{Referencia posterior} & \HMTtablehead{Diferencia relativa (ppm)} & \HMTtablehead{Naturaleza y alcance} \\
\midrule
\endfirsthead
\multicolumn{5}{l}{\sffamily\scriptsize\color{HMTGray}Continuación de la tabla}\\[-1pt]
\toprule
\rowcolor{HMTNavy}\HMTtablehead{Estado} & \HMTtablehead{Valor HMT reproducible (MeV/\allowbreak{}c\({}^{2}\))} & \HMTtablehead{Referencia posterior} & \HMTtablehead{Diferencia relativa (ppm)} & \HMTtablehead{Naturaleza y alcance} \\
\midrule
\endhead
\midrule
\endfoot
\bottomrule
\endlastfoot
e-/\allowbreak{}e+ & 0.510998\allowbreak{}950690\allowbreak{}000809 & 0.510998\allowbreak{}95069±\allowbreak{}0.000000\allowbreak{}00016 MeV/\allowbreak{}c\({}^{2}\) & 1.582406\allowbreak{}884×10\({}^{-9}\) & coordenada beta cerrada en la fibra central \\
\rowcolor{HMTLight}μ-/\allowbreak{}μ+ & 105.658360\allowbreak{}294435\allowbreak{}829303 & 105.658375\allowbreak{}5±\allowbreak{}0.000002\allowbreak{}3 MeV/\allowbreak{}c\({}^{2}\) & -0.143912\allowbreak{}53 & refinamiento E3 condicionado a la firma \\
τ-/\allowbreak{}τ+ & 1776.860917\allowbreak{}631432\allowbreak{}295595 & 1776.932465\allowbreak{}134090\allowbreak{}1 MeV/\allowbreak{}c\({}^{2}\) (valor central usado; presentación redondeada: 1776.93±\allowbreak{}0.09 MeV/\allowbreak{}c\({}^{2}\)) & -40.264615\allowbreak{}601 & refinamiento E3 condicionado a la firma \\
\rowcolor{HMTLight}u/\allowbreak{}\(\bar{\mathrm{u}}\) & 2.165924\allowbreak{}536173\allowbreak{}907663 & 2.16±\allowbreak{}0.07 MeV/\allowbreak{}c\({}^{2}\) & --- & carta angular; no se ha declarado firma Z\({}^{5}\) ni refinamiento E3 \\
d/\allowbreak{}\(\bar{\mathrm{d}}\) & 4.679550\allowbreak{}162948\allowbreak{}874172 & 4.70±\allowbreak{}0.07 MeV/\allowbreak{}c\({}^{2}\) & --- & carta angular; no se ha declarado firma Z\({}^{5}\) ni refinamiento E3 \\
\rowcolor{HMTLight}c/\allowbreak{}\(\bar{\mathrm{c}}\) & 1270.500838\allowbreak{}641911\allowbreak{}135286 & 1.2729±\allowbreak{}0.0045 GeV/\allowbreak{}c\({}^{2}\) & --- & carta angular; no se ha declarado firma Z\({}^{5}\) ni refinamiento E3 \\
s/\allowbreak{}\(\bar{\mathrm{s}}\) & 92.940093\allowbreak{}907845\allowbreak{}640641 & 92.9 ±\allowbreak{}0.7 MeV/\allowbreak{}c\({}^{2}\) & --- & carta angular; no se ha declarado firma Z\({}^{5}\) ni refinamiento E3 \\
\rowcolor{HMTLight}t/\allowbreak{}\(\bar{\mathrm{t}}\) & 172597.479544\allowbreak{}019136\allowbreak{}261367 & 172.60±\allowbreak{}0.27 GeV/\allowbreak{}c\({}^{2}\) & --- & carta angular; no se ha declarado firma Z\({}^{5}\) ni refinamiento E3 \\
b/\allowbreak{}\(\bar{\mathrm{b}}\) & 4190.119763\allowbreak{}299790\allowbreak{}218075 & 4.186±\allowbreak{}0.006 GeV/\allowbreak{}c\({}^{2}\) & --- & carta angular; no se ha declarado firma Z\({}^{5}\) ni refinamiento E3 \\
\rowcolor{HMTLight}γ & 0 & <1×10\({}^{-18}\) eV/\allowbreak{}c\({}^{2}\) & --- & salida nula exacta \\
g & 0 & --- & --- & salida nula exacta \\
\rowcolor{HMTLight}W+/\allowbreak{}W- & 80378.964909\allowbreak{}704547\allowbreak{}024865 & 80.3625 ±\allowbreak{}0.0077 GeV/\allowbreak{}c\({}^{2}\) & 204.882995\allowbreak{}235 & refinamiento E3 condicionado a la firma \\
Z0 & 91187.533444\allowbreak{}313407\allowbreak{}844855 & 91.187873\allowbreak{}297221\allowbreak{}344 GeV/\allowbreak{}c\({}^{2}\) (valor central usado; presentación redondeada: 91.1879±\allowbreak{}0.0020 GeV/\allowbreak{}c\({}^{2}\)) & -3.726952\allowbreak{}89 & refinamiento E3 condicionado a la firma \\
\rowcolor{HMTLight}H & 125292.466042\allowbreak{}536133\allowbreak{}000433 & 125.130943\allowbreak{}828161\allowbreak{}51 GeV/\allowbreak{}c\({}^{2}\) (valor central usado; presentación redondeada: 125.13±\allowbreak{}0.11 GeV/\allowbreak{}c\({}^{2}\)) & 1290.825509\allowbreak{}927 & carta angular; no se ha declarado firma Z\({}^{5}\) ni refinamiento E3 \\
π\({}^{0}\) & 134.976724\allowbreak{}327872\allowbreak{}125739 & 134.976827\allowbreak{}767685 MeV/\allowbreak{}c\({}^{2}\) & -0.766352\allowbreak{}378 & refinamiento E3 condicionado a la firma compuesta \\
\rowcolor{HMTLight}π±\allowbreak{} & 139.570485\allowbreak{}171572\allowbreak{}191375 & 139.570390\allowbreak{}983681 MeV/\allowbreak{}c\({}^{2}\) & 0.674841\allowbreak{}494 & refinamiento E3 condicionado a la firma compuesta \\
K±\allowbreak{} & 493.677085\allowbreak{}649530\allowbreak{}612476 & 493.676599\allowbreak{}458041 MeV/\allowbreak{}c\({}^{2}\) & 0.984838\allowbreak{}03 & refinamiento E3 condicionado a la firma compuesta \\
\rowcolor{HMTLight}K\({}^{0}\) & 497.611186\allowbreak{}497750\allowbreak{}457359 & 497.610767\allowbreak{}531258 MeV/\allowbreak{}c\({}^{2}\) & 0.841956\allowbreak{}243 & refinamiento E3 condicionado a la firma compuesta \\
protón & 938.271751\allowbreak{}546984\allowbreak{}898299 & 938.272089\allowbreak{}43 MeV/\allowbreak{}c\({}^{2}\) & -0.360111\allowbreak{}975 & refinamiento E3 condicionado a la firma compuesta \\
\rowcolor{HMTLight}neutrón & 939.565810\allowbreak{}472507\allowbreak{}023313 & 939.565421\allowbreak{}94 MeV/\allowbreak{}c\({}^{2}\) & 0.413523\allowbreak{}633 & refinamiento E3 condicionado a la firma compuesta \\
\end{longtable}
\endgroup
\Needspace{7\baselineskip}
\section{Evaluación individual de \(20\) secciones materiales y trazabilidad de sus coordenadas metrológicas}
\Needspace{7\baselineskip}
\subsection{Par electrón–positrón: ruta central, conjugación de hojas y carácter de masa}
\begin{itemize}
\item \textbf{Identificador catalográfico:} SM-LC-1.
\item \textbf{Tipo de descriptor:} selector basal y lectores bidireccionales.
\item \textbf{Descriptor:} selector (-22,+15); K\_D=\allowbreak{}K\_Ω=\allowbreak{}(80,54,6).
\item \textbf{Etapa:} beta bidireccional; coordenada inicial 0.510998\allowbreak{}922488\allowbreak{}066073 MeV/c\({}^{2}\).
\item \textbf{Evaluación HMT--MD reproducible:} 0.510998\allowbreak{}950690\allowbreak{}000809 MeV/c\({}^{2}\).
\item \textbf{Naturaleza matemática:} coordenada beta cerrada en la fibra central.
\item \textbf{Referencia posterior:} 0.510998\allowbreak{}95069±\allowbreak{}0.000000\allowbreak{}00016 MeV/c\({}^{2}\).
\item \textbf{Identidad y metrología externas:} Particle Data Group (PDG); PDG2026:\allowbreak{}observable:\allowbreak{}S003:\allowbreak{}S003M:\allowbreak{}332838; esquema no conservado en el corte documental; escala no conservada en el corte documental; nivel de confianza no conservado en el corte documental.
\item \textbf{Diferencia relativa:} 1.582406\allowbreak{}884×10\({}^{-9}\) ppm.
\item \textbf{Alcance:} derivación escalar cerrada.
\end{itemize}
\Needspace{7\baselineskip}
\subsection{Par muón–antimuón: conjugación de rutas y carácter leptónico}
\begin{itemize}
\item \textbf{Identificador catalográfico:} SM-LC-2.
\item \textbf{Tipo de descriptor:} firma completa σ\(\in\)Z\({}^{5}\).
\item \textbf{Descriptor:} (42,-3,8,0,2).
\item \textbf{Etapa:} torre beta E3; coordenada inicial 105.564059\allowbreak{}012859\allowbreak{}26472 MeV/c\({}^{2}\).
\item \textbf{Evaluación HMT--MD reproducible:} 105.658360\allowbreak{}294435\allowbreak{}829303 MeV/c\({}^{2}\).
\item \textbf{Naturaleza matemática:} refinamiento E3 condicionado a la firma.
\item \textbf{Referencia posterior:} 105.658375\allowbreak{}5±\allowbreak{}0.000002\allowbreak{}3 MeV/c\({}^{2}\).
\item \textbf{Identidad y metrología externas:} Particle Data Group (PDG); PDG2026:\allowbreak{}observable:\allowbreak{}S004:\allowbreak{}S004M:\allowbreak{}332849; esquema no conservado en el corte documental; escala no conservada en el corte documental; nivel de confianza no conservado en el corte documental.
\item \textbf{Diferencia relativa:} -0.143912\allowbreak{}53 ppm.
\item \textbf{Alcance:} evaluación exacta al fijar la firma; la selección física de esa firma constituye una operación independiente.
\end{itemize}
\Needspace{7\baselineskip}
\subsection{Par tau–antitau: conjugación de rutas y carácter leptónico}
\begin{itemize}
\item \textbf{Identificador catalográfico:} SM-LC-3.
\item \textbf{Tipo de descriptor:} firma completa σ\(\in\)Z\({}^{5}\).
\item \textbf{Descriptor:} (64,0,25,-1,6).
\item \textbf{Etapa:} torre beta E3; coordenada inicial 1771.945393\allowbreak{}886234\allowbreak{}110148 MeV/c\({}^{2}\).
\item \textbf{Evaluación HMT--MD reproducible:} 1776.860917\allowbreak{}631432\allowbreak{}295595 MeV/c\({}^{2}\).
\item \textbf{Naturaleza matemática:} refinamiento E3 condicionado a la firma.
\item \textbf{Referencia posterior:} 1776.932465\allowbreak{}134090\allowbreak{}1 MeV/c\({}^{2}\) (valor central usado; presentación redondeada: 1776.93±\allowbreak{}0.09 MeV/c\({}^{2}\)).
\item \textbf{Identidad y metrología externas:} Particle Data Group (PDG); PDG2026:\allowbreak{}observable:\allowbreak{}S035:\allowbreak{}S035M:\allowbreak{}332893; esquema no conservado en el corte documental; escala no conservada en el corte documental; nivel de confianza no conservado en el corte documental.
\item \textbf{Diferencia relativa:} -40.264615\allowbreak{}601 ppm.
\item \textbf{Alcance:} evaluación exacta al fijar la firma; la selección física de esa firma constituye una operación independiente.
\end{itemize}
\Needspace{7\baselineskip}
\subsection{Quark arriba y antiquark arriba}
\begin{itemize}
\item \textbf{Identificador catalográfico:} SM-Q-U1.
\item \textbf{Tipo de descriptor:} coindexación angular (n\_A,s\_C)\(\in\)Z\({}^{2}\).
\item \textbf{Descriptor:} (11,7).
\item \textbf{Etapa:} carta angular condicionada; coordenada inicial 2.165924\allowbreak{}536173\allowbreak{}907663 MeV/c\({}^{2}\).
\item \textbf{Evaluación HMT--MD reproducible:} 2.165924\allowbreak{}536173\allowbreak{}907663 MeV/c\({}^{2}\).
\item \textbf{Naturaleza matemática:} carta angular; no se ha declarado firma Z\({}^{5}\) ni refinamiento E3.
\item \textbf{Referencia posterior:} 2.16±\allowbreak{}0.07 MeV/c\({}^{2}\).
\item \textbf{Identidad y metrología externas:} Particle Data Group (PDG); PDG2026:\allowbreak{}observable:\allowbreak{}Q002:\allowbreak{}Q002M:\allowbreak{}333588; esquema no conservado en el corte documental; escala no conservada en el corte documental; nivel de confianza 90\%.
\item \textbf{Diferencia relativa:} --- ppm.
\item \textbf{Alcance:} evaluación condicionada; la firma no se ha elegido desde el valor comparado; esquema y escala externos requeridos para el residual.
\end{itemize}
\Needspace{7\baselineskip}
\subsection{Quark abajo y antiquark abajo}
\begin{itemize}
\item \textbf{Identificador catalográfico:} SM-Q-D1.
\item \textbf{Tipo de descriptor:} coindexación angular (n\_A,s\_C)\(\in\)Z\({}^{2}\).
\item \textbf{Descriptor:} (17,8).
\item \textbf{Etapa:} carta angular condicionada; coordenada inicial 4.679550\allowbreak{}162948\allowbreak{}874172 MeV/c\({}^{2}\).
\item \textbf{Evaluación HMT--MD reproducible:} 4.679550\allowbreak{}162948\allowbreak{}874172 MeV/c\({}^{2}\).
\item \textbf{Naturaleza matemática:} carta angular; no se ha declarado firma Z\({}^{5}\) ni refinamiento E3.
\item \textbf{Referencia posterior:} 4.70±\allowbreak{}0.07 MeV/c\({}^{2}\).
\item \textbf{Identidad y metrología externas:} Particle Data Group (PDG); PDG2026:\allowbreak{}observable:\allowbreak{}Q001:\allowbreak{}Q001M:\allowbreak{}333589; esquema no conservado en el corte documental; escala no conservada en el corte documental; nivel de confianza 90\%.
\item \textbf{Diferencia relativa:} --- ppm.
\item \textbf{Alcance:} evaluación condicionada; la firma no se ha elegido desde el valor comparado; esquema y escala externos requeridos para el residual.
\end{itemize}
\Needspace{7\baselineskip}
\subsection{Quark encanto y antiquark encanto}
\begin{itemize}
\item \textbf{Identificador catalográfico:} SM-Q-U2.
\item \textbf{Tipo de descriptor:} coindexación angular (n\_A,s\_C)\(\in\)Z\({}^{2}\).
\item \textbf{Descriptor:} (61,8).
\item \textbf{Etapa:} carta angular condicionada; coordenada inicial 1270.500838\allowbreak{}641911\allowbreak{}135286 MeV/c\({}^{2}\).
\item \textbf{Evaluación HMT--MD reproducible:} 1270.500838\allowbreak{}641911\allowbreak{}135286 MeV/c\({}^{2}\).
\item \textbf{Naturaleza matemática:} carta angular; no se ha declarado firma Z\({}^{5}\) ni refinamiento E3.
\item \textbf{Referencia posterior:} 1.2729±\allowbreak{}0.0045 GeV/c\({}^{2}\).
\item \textbf{Identidad y metrología externas:} Particle Data Group (PDG); PDG2026:\allowbreak{}observable:\allowbreak{}Q004:\allowbreak{}Q004M:\allowbreak{}333604; esquema no conservado en el corte documental; escala no conservada en el corte documental; nivel de confianza 90\%.
\item \textbf{Diferencia relativa:} --- ppm.
\item \textbf{Alcance:} evaluación condicionada; la firma no se ha elegido desde el valor comparado; esquema y escala externos requeridos para el residual.
\end{itemize}
\Needspace{7\baselineskip}
\subsection{Quark extraño y antiquark extraño}
\begin{itemize}
\item \textbf{Identificador catalográfico:} SM-Q-D2.
\item \textbf{Tipo de descriptor:} coindexación angular (n\_A,s\_C)\(\in\)Z\({}^{2}\).
\item \textbf{Descriptor:} (41,-3).
\item \textbf{Etapa:} carta angular condicionada; coordenada inicial 92.940093\allowbreak{}907845\allowbreak{}640641 MeV/c\({}^{2}\).
\item \textbf{Evaluación HMT--MD reproducible:} 92.940093\allowbreak{}907845\allowbreak{}640641 MeV/c\({}^{2}\).
\item \textbf{Naturaleza matemática:} carta angular; no se ha declarado firma Z\({}^{5}\) ni refinamiento E3.
\item \textbf{Referencia posterior:} 92.9 ±\allowbreak{}0.7 MeV/c\({}^{2}\).
\item \textbf{Identidad y metrología externas:} Particle Data Group (PDG); PDG2026:\allowbreak{}observable:\allowbreak{}Q003:\allowbreak{}Q003M:\allowbreak{}333590; esquema no conservado en el corte documental; escala no conservada en el corte documental; nivel de confianza 90\%.
\item \textbf{Diferencia relativa:} --- ppm.
\item \textbf{Alcance:} evaluación condicionada; la firma no se ha elegido desde el valor comparado; esquema y escala externos requeridos para el residual.
\end{itemize}
\Needspace{7\baselineskip}
\subsection{Quark cima y antiquark cima}
\begin{itemize}
\item \textbf{Identificador catalográfico:} SM-Q-U3.
\item \textbf{Tipo de descriptor:} coindexación angular (n\_A,s\_C)\(\in\)Z\({}^{2}\).
\item \textbf{Descriptor:} (100,-1).
\item \textbf{Etapa:} carta angular condicionada; coordenada inicial 172597.479544\allowbreak{}019136\allowbreak{}261367 MeV/c\({}^{2}\).
\item \textbf{Evaluación HMT--MD reproducible:} 172597.479544\allowbreak{}019136\allowbreak{}261367 MeV/c\({}^{2}\).
\item \textbf{Naturaleza matemática:} carta angular; no se ha declarado firma Z\({}^{5}\) ni refinamiento E3.
\item \textbf{Referencia posterior:} 172.60±\allowbreak{}0.27 GeV/c\({}^{2}\).
\item \textbf{Identidad y metrología externas:} Particle Data Group (PDG); PDG2026:\allowbreak{}observable:\allowbreak{}Q007:\allowbreak{}Q007TP:\allowbreak{}333614; esquema no conservado en el corte documental; escala no conservada en el corte documental; nivel de confianza no conservado en el corte documental.
\item \textbf{Diferencia relativa:} --- ppm.
\item \textbf{Alcance:} evaluación condicionada; la firma no se ha elegido desde el valor comparado; esquema y escala externos requeridos para el residual.
\end{itemize}
\Needspace{7\baselineskip}
\subsection{Quark fondo y antiquark fondo}
\begin{itemize}
\item \textbf{Identificador catalográfico:} SM-Q-D3.
\item \textbf{Tipo de descriptor:} coindexación angular (n\_A,s\_C)\(\in\)Z\({}^{2}\).
\item \textbf{Descriptor:} (71,-5).
\item \textbf{Etapa:} carta angular condicionada; coordenada inicial 4190.119763\allowbreak{}299790\allowbreak{}218075 MeV/c\({}^{2}\).
\item \textbf{Evaluación HMT--MD reproducible:} 4190.119763\allowbreak{}299790\allowbreak{}218075 MeV/c\({}^{2}\).
\item \textbf{Naturaleza matemática:} carta angular; no se ha declarado firma Z\({}^{5}\) ni refinamiento E3.
\item \textbf{Referencia posterior:} 4.186±\allowbreak{}0.006 GeV/c\({}^{2}\).
\item \textbf{Identidad y metrología externas:} Particle Data Group (PDG); PDG2026:\allowbreak{}observable:\allowbreak{}Q005:\allowbreak{}Q005M:\allowbreak{}333611; esquema no conservado en el corte documental; escala no conservada en el corte documental; nivel de confianza 90\%.
\item \textbf{Diferencia relativa:} --- ppm.
\item \textbf{Alcance:} evaluación condicionada; la firma no se ha elegido desde el valor comparado; esquema y escala externos requeridos para el residual.
\end{itemize}
\Needspace{7\baselineskip}
\subsection{Descriptor espectral reproducible del fotón: sector nulo, helicidad y carácter electromagnético}
\begin{itemize}
\item \textbf{Identificador catalográfico:} SM-GA-EM.
\item \textbf{Tipo de descriptor:} invariante sectorial nulo.
\item \textbf{Descriptor:} carta nula.
\item \textbf{Etapa:} carta nula; coordenada inicial 0 MeV/c\({}^{2}\).
\item \textbf{Evaluación HMT--MD reproducible:} 0 MeV/c\({}^{2}\).
\item \textbf{Naturaleza matemática:} salida nula exacta.
\item \textbf{Referencia posterior:} <1×10\({}^{-18}\) eV/c\({}^{2}\).
\item \textbf{Identidad y metrología externas:} Particle Data Group (PDG); PDG2026:\allowbreak{}observable:\allowbreak{}S000:\allowbreak{}S000M:\allowbreak{}332461; esquema no conservado en el corte documental; escala no conservada en el corte documental; nivel de confianza no conservado en el corte documental.
\item \textbf{Diferencia relativa:} --- ppm.
\item \textbf{Alcance:} exacto en el sector nulo; la cota externa no se interpreta como una masa positiva.
\item \textbf{Cota fotónica:} referencia externa ilustrativa no comparable; el corte identifica el observable pero no conserva el nivel de confianza.
\end{itemize}
\Needspace{7\baselineskip}
\subsection{Realización del gluón como sección de la representación adjunta de \(SU(3)\)}
\begin{itemize}
\item \textbf{Identificador catalográfico:} SM-GA-GL.
\item \textbf{Tipo de descriptor:} invariante sectorial nulo.
\item \textbf{Descriptor:} carta nula.
\item \textbf{Etapa:} carta nula; coordenada inicial 0 MeV/c\({}^{2}\).
\item \textbf{Evaluación HMT--MD reproducible:} 0 MeV/c\({}^{2}\).
\item \textbf{Naturaleza matemática:} salida nula exacta.
\item \textbf{Referencia posterior:} ---.
\item \textbf{Identidad y metrología externas:} fuente no consignada; observable no individualizado; esquema no consignado; escala no consignada; nivel de confianza no consignado.
\item \textbf{Diferencia relativa:} --- ppm.
\item \textbf{Alcance:} exacto en el sector nulo; la nulidad pertenece a la carta de calibre de color.
\end{itemize}
\Needspace{7\baselineskip}
\subsection{Bosones \texorpdfstring{\(W^{+}\)}{W+} y \texorpdfstring{\(W^{-}\)}{W−}}
\begin{itemize}
\item \textbf{Identificador catalográfico:} SM-GA-W.
\item \textbf{Tipo de descriptor:} firma completa σ\(\in\)Z\({}^{5}\).
\item \textbf{Descriptor:} (94,-1,0,-2,4).
\item \textbf{Etapa:} torre beta E3; coordenada inicial 80381.592550\allowbreak{}220666\allowbreak{}304144 MeV/c\({}^{2}\).
\item \textbf{Evaluación HMT--MD reproducible:} 80378.964909\allowbreak{}704547\allowbreak{}024865 MeV/c\({}^{2}\).
\item \textbf{Naturaleza matemática:} refinamiento E3 condicionado a la firma.
\item \textbf{Referencia posterior:} 80.3625 ±\allowbreak{}0.0077 GeV/c\({}^{2}\).
\item \textbf{Identidad y metrología externas:} Particle Data Group (PDG); PDG2026:\allowbreak{}observable:\allowbreak{}S043:\allowbreak{}S043M:\allowbreak{}332464; esquema no conservado en el corte documental; escala no conservada en el corte documental; nivel de confianza no conservado en el corte documental.
\item \textbf{Diferencia relativa:} 204.882995\allowbreak{}235 ppm.
\item \textbf{Alcance:} evaluación exacta al fijar la firma; la selección física de esa firma constituye una operación independiente.
\item \textbf{Convención resonante:} el corte externo no fija conjuntamente convención de polo, Breit--Wigner o sobre capa, anchura y covarianza; el residual se interpreta sólo como información escalar.
\end{itemize}
\Needspace{7\baselineskip}
\subsection{Descriptor espectral reproducible del bosón \(Z^0\): mezcla neutra, polarización y carácter electrodébil}
\begin{itemize}
\item \textbf{Identificador catalográfico:} SM-GA-Z.
\item \textbf{Tipo de descriptor:} firma completa σ\(\in\)Z\({}^{5}\).
\item \textbf{Descriptor:} (95,-1,-11,0,-4).
\item \textbf{Etapa:} torre beta E3; coordenada inicial 91299.748286\allowbreak{}598170\allowbreak{}117317 MeV/c\({}^{2}\).
\item \textbf{Evaluación HMT--MD reproducible:} 91187.533444\allowbreak{}313407\allowbreak{}844855 MeV/c\({}^{2}\).
\item \textbf{Naturaleza matemática:} refinamiento E3 condicionado a la firma.
\item \textbf{Referencia posterior:} 91.187873\allowbreak{}297221\allowbreak{}344 GeV/c\({}^{2}\) (valor central usado; presentación redondeada: 91.1879±\allowbreak{}0.0020 GeV/c\({}^{2}\)).
\item \textbf{Identidad y metrología externas:} Particle Data Group (PDG); PDG2026:\allowbreak{}observable:\allowbreak{}S044:\allowbreak{}S044M:\allowbreak{}332513; esquema no conservado en el corte documental; escala no conservada en el corte documental; nivel de confianza no conservado en el corte documental.
\item \textbf{Diferencia relativa:} -3.726952\allowbreak{}89 ppm.
\item \textbf{Alcance:} evaluación exacta al fijar la firma; la selección física de esa firma constituye una operación independiente.
\item \textbf{Convención resonante:} el corte externo no fija conjuntamente convención de polo, Breit--Wigner o sobre capa, anchura y covarianza; el residual se interpreta sólo como información escalar.
\end{itemize}
\Needspace{7\baselineskip}
\subsection{Realización del bosón de Higgs: modo escalar, acoplamiento y carácter de masa}
\begin{itemize}
\item \textbf{Identificador catalográfico:} SM-H-1.
\item \textbf{Tipo de descriptor:} coindexación angular (n\_A,s\_C)\(\in\)Z\({}^{2}\).
\item \textbf{Descriptor:} (97,9).
\item \textbf{Etapa:} carta angular condicionada; coordenada inicial 125292.466042\allowbreak{}536133\allowbreak{}000433 MeV/c\({}^{2}\).
\item \textbf{Evaluación HMT--MD reproducible:} 125292.466042\allowbreak{}536133\allowbreak{}000433 MeV/c\({}^{2}\).
\item \textbf{Naturaleza matemática:} carta angular; no se ha declarado firma Z\({}^{5}\) ni refinamiento E3.
\item \textbf{Referencia posterior:} 125.130943\allowbreak{}828161\allowbreak{}51 GeV/c\({}^{2}\) (valor central usado; presentación redondeada: 125.13±\allowbreak{}0.11 GeV/c\({}^{2}\)).
\item \textbf{Identidad y metrología externas:} Particle Data Group (PDG); PDG2026:\allowbreak{}observable:\allowbreak{}S126:\allowbreak{}S126M:\allowbreak{}332733; esquema no conservado en el corte documental; escala no conservada en el corte documental; nivel de confianza no conservado en el corte documental.
\item \textbf{Diferencia relativa:} 1290.825509\allowbreak{}927 ppm.
\item \textbf{Alcance:} evaluación condicionada; la firma no se ha elegido desde el valor comparado.
\item \textbf{Convención resonante:} el corte externo no fija conjuntamente convención de polo, Breit--Wigner o sobre capa, anchura y covarianza; el residual se interpreta sólo como información escalar.
\end{itemize}
\Needspace{7\baselineskip}
\subsection{Pión neutro \texorpdfstring{\(\pi^{0}\)}{pi0}}
\begin{itemize}
\item \textbf{Identificador catalográfico:} CMP-PI0.
\item \textbf{Tipo de descriptor:} firma completa σ\(\in\)Z\({}^{5}\).
\item \textbf{Descriptor:} (44,-4,-25,1,-2).
\item \textbf{Etapa:} torre beta E3 compuesta; coordenada inicial 135.350257\allowbreak{}251504\allowbreak{}283349 MeV/c\({}^{2}\).
\item \textbf{Evaluación HMT--MD reproducible:} 134.976724\allowbreak{}327872\allowbreak{}125739 MeV/c\({}^{2}\).
\item \textbf{Naturaleza matemática:} refinamiento E3 condicionado a la firma compuesta.
\item \textbf{Referencia posterior:} 134.976827\allowbreak{}767685 MeV/c\({}^{2}\).
\item \textbf{Identidad y metrología externas:} referencia compuesta conservada por el suplemento E3; la convención, la covarianza y la edición se consignan en la fuente metrológica de referencia.
\item \textbf{Diferencia relativa:} -0.766352\allowbreak{}378 ppm.
\item \textbf{Alcance:} evaluación exacta al fijar la firma compuesta; el generador de ligadura permanece como morfismo propio.
\end{itemize}
\Needspace{7\baselineskip}
\subsection{Piones cargados \(\pi^\pm\): composición quark–antiquark, carga y masa}
\begin{itemize}
\item \textbf{Identificador catalográfico:} CMP-PIPM.
\item \textbf{Tipo de descriptor:} firma completa σ\(\in\)Z\({}^{5}\).
\item \textbf{Descriptor:} (44,1,-2,2,-1).
\item \textbf{Etapa:} torre beta E3 compuesta; coordenada inicial 139.596265\allowbreak{}529971\allowbreak{}781015 MeV/c\({}^{2}\).
\item \textbf{Evaluación HMT--MD reproducible:} 139.570485\allowbreak{}171572\allowbreak{}191375 MeV/c\({}^{2}\).
\item \textbf{Naturaleza matemática:} refinamiento E3 condicionado a la firma compuesta.
\item \textbf{Referencia posterior:} 139.570390\allowbreak{}983681 MeV/c\({}^{2}\).
\item \textbf{Identidad y metrología externas:} referencia compuesta conservada por el suplemento E3; la convención, la covarianza y la edición se consignan en la fuente metrológica de referencia.
\item \textbf{Diferencia relativa:} 0.674841\allowbreak{}494 ppm.
\item \textbf{Alcance:} evaluación exacta al fijar la firma compuesta; el generador de ligadura permanece como morfismo propio.
\end{itemize}
\Needspace{7\baselineskip}
\subsection{Kaones cargados \(K^\pm\): composición con extrañeza, carga y masa}
\begin{itemize}
\item \textbf{Identificador catalográfico:} CMP-KPM.
\item \textbf{Tipo de descriptor:} firma completa σ\(\in\)Z\({}^{5}\).
\item \textbf{Descriptor:} (54,-1,17,-2,2).
\item \textbf{Etapa:} torre beta E3 compuesta; coordenada inicial 492.762503\allowbreak{}210140\allowbreak{}547854 MeV/c\({}^{2}\).
\item \textbf{Evaluación HMT--MD reproducible:} 493.677085\allowbreak{}649530\allowbreak{}612476 MeV/c\({}^{2}\).
\item \textbf{Naturaleza matemática:} refinamiento E3 condicionado a la firma compuesta.
\item \textbf{Referencia posterior:} 493.676599\allowbreak{}458041 MeV/c\({}^{2}\).
\item \textbf{Identidad y metrología externas:} referencia compuesta conservada por el suplemento E3; la convención, la covarianza y la edición se consignan en la fuente metrológica de referencia.
\item \textbf{Diferencia relativa:} 0.984838\allowbreak{}03 ppm.
\item \textbf{Alcance:} evaluación exacta al fijar la firma compuesta; el generador de ligadura permanece como morfismo propio.
\end{itemize}
\Needspace{7\baselineskip}
\subsection{Kaón neutro \texorpdfstring{\(K^{0}\)}{K0}}
\begin{itemize}
\item \textbf{Identificador catalográfico:} CMP-K0.
\item \textbf{Tipo de descriptor:} firma completa σ\(\in\)Z\({}^{5}\).
\item \textbf{Descriptor:} (54,0,32,3,-2).
\item \textbf{Etapa:} torre beta E3 compuesta; coordenada inicial 495.816066\allowbreak{}594426\allowbreak{}768595 MeV/c\({}^{2}\).
\item \textbf{Evaluación HMT--MD reproducible:} 497.611186\allowbreak{}497750\allowbreak{}457359 MeV/c\({}^{2}\).
\item \textbf{Naturaleza matemática:} refinamiento E3 condicionado a la firma compuesta.
\item \textbf{Referencia posterior:} 497.610767\allowbreak{}531258 MeV/c\({}^{2}\).
\item \textbf{Identidad y metrología externas:} referencia compuesta conservada por el suplemento E3; la convención, la covarianza y la edición se consignan en la fuente metrológica de referencia.
\item \textbf{Diferencia relativa:} 0.841956\allowbreak{}243 ppm.
\item \textbf{Alcance:} evaluación exacta al fijar la firma compuesta; el generador de ligadura permanece como morfismo propio.
\end{itemize}
\Needspace{7\baselineskip}
\subsection{Protón: composición bariónica, neutralización de color y carácter de masa}
\begin{itemize}
\item \textbf{Identificador catalográfico:} CMP-P.
\item \textbf{Tipo de descriptor:} firma completa σ\(\in\)Z\({}^{5}\).
\item \textbf{Descriptor:} (59,0,9,1,0).
\item \textbf{Etapa:} torre beta E3 compuesta; coordenada inicial 937.314779\allowbreak{}258699\allowbreak{}634563 MeV/c\({}^{2}\).
\item \textbf{Evaluación HMT--MD reproducible:} 938.271751\allowbreak{}546984\allowbreak{}898299 MeV/c\({}^{2}\).
\item \textbf{Naturaleza matemática:} refinamiento E3 condicionado a la firma compuesta.
\item \textbf{Referencia posterior:} 938.272089\allowbreak{}43 MeV/c\({}^{2}\).
\item \textbf{Identidad y metrología externas:} referencia compuesta conservada por el suplemento E3; la convención, la covarianza y la edición se consignan en la fuente metrológica de referencia.
\item \textbf{Diferencia relativa:} -0.360111\allowbreak{}975 ppm.
\item \textbf{Alcance:} evaluación exacta al fijar la firma compuesta; el generador de ligadura permanece como morfismo propio.
\end{itemize}
\Needspace{7\baselineskip}
\subsection{Neutrón: composición bariónica, neutralización de color y carácter de masa}
\begin{itemize}
\item \textbf{Identificador catalográfico:} CMP-N.
\item \textbf{Tipo de descriptor:} firma completa σ\(\in\)Z\({}^{5}\).
\item \textbf{Descriptor:} (59,1,-34,0,1).
\item \textbf{Etapa:} torre beta E3 compuesta; coordenada inicial 943.123155\allowbreak{}648642\allowbreak{}014982 MeV/c\({}^{2}\).
\item \textbf{Evaluación HMT--MD reproducible:} 939.565810\allowbreak{}472507\allowbreak{}023313 MeV/c\({}^{2}\).
\item \textbf{Naturaleza matemática:} refinamiento E3 condicionado a la firma compuesta.
\item \textbf{Referencia posterior:} 939.565421\allowbreak{}94 MeV/c\({}^{2}\).
\item \textbf{Identidad y metrología externas:} referencia compuesta conservada por el suplemento E3; la convención, la covarianza y la edición se consignan en la fuente metrológica de referencia.
\item \textbf{Diferencia relativa:} 0.413523\allowbreak{}633 ppm.
\item \textbf{Alcance:} evaluación exacta al fijar la firma compuesta; el generador de ligadura permanece como morfismo propio.
\end{itemize}
% Las tablas históricas abreviadas permanecen conservadas a continuación en
% este propietario, pero quedan fuera del tronco científico vigente: no
% definen el dominio 81/56/13/324, la firma, el carácter ni las torres.
\endinput
\clearpage
\begin{landscape}
\section{Malla histórica extendida conservada}
\begingroup
\fontsize{6.8}{7.7}\selectfont
\setlength{\tabcolsep}{1.15pt}
\renewcommand{\arraystretch}{1.06}
\begin{longtable}{L{2.93cm}L{2.93cm}L{2.93cm}L{2.93cm}L{2.93cm}L{2.93cm}L{2.93cm}L{2.93cm}}
\toprule
\rowcolor{HMTNavy}\HMTtablehead{Estado} & \HMTtablehead{n} & \HMTtablehead{k} & \HMTtablehead{t} & \HMTtablehead{Evaluación HMT histórica (MeV/\allowbreak{}c\({}^{2}\))} & \HMTtablehead{Referencia posterior (MeV/\allowbreak{}c\({}^{2}\))} & \HMTtablehead{Diferencia consignada (ppm)} & \HMTtablehead{Fuerza matemática} \\
\midrule
\endfirsthead
\multicolumn{8}{l}{\sffamily\scriptsize\color{HMTGray}Continuación de la tabla}\\[-1pt]
\toprule
\rowcolor{HMTNavy}\HMTtablehead{Estado} & \HMTtablehead{n} & \HMTtablehead{k} & \HMTtablehead{t} & \HMTtablehead{Evaluación HMT histórica (MeV/\allowbreak{}c\({}^{2}\))} & \HMTtablehead{Referencia posterior (MeV/\allowbreak{}c\({}^{2}\))} & \HMTtablehead{Diferencia consignada (ppm)} & \HMTtablehead{Fuerza matemática} \\
\midrule
\endhead
\midrule
\endfoot
\bottomrule
\endlastfoot
mu & 64 & 146 & 1 & 105.656529 & 105.658374 & -17.461938\allowbreak{}2274 & evidencia numérica histórica condicionada; no constituye una regla física de selección anterior al contraste \\
\rowcolor{HMTLight}tau & 64 & 349 & -1 & 1776.835 & 1776.86 & -14.069763\allowbreak{}5154 & evidencia numérica histórica condicionada; no constituye una regla física de selección anterior al contraste \\
π±\allowbreak{} & 37 & 125 & 0 & 139.570582 & 139.57039 & 1.375649\allowbreak{}94982 & evidencia numérica histórica condicionada; no constituye una regla física de selección anterior al contraste \\
\rowcolor{HMTLight}K±\allowbreak{} & 52 & 177 & 1 & 493.679 & 493.677 & 4.051231\allowbreak{}87833 & evidencia numérica histórica condicionada; no constituye una regla física de selección anterior al contraste \\
K\({}^{0}\) & 53 & 178 & 1 & 497.616 & 497.611 & 10.048009\allowbreak{}3889 & evidencia numérica histórica condicionada; no constituye una regla física de selección anterior al contraste \\
\rowcolor{HMTLight}p & 55 & 195 & 1 & 938.27244 & 938.272081 & 0.382618\allowbreak{}226919 & evidencia numérica histórica condicionada; no constituye una regla física de selección anterior al contraste \\
n & 55 & 197 & 0 & 939.538 & 939.565413 & -29.176254\allowbreak{}9161 & evidencia numérica histórica condicionada; no constituye una regla física de selección anterior al contraste \\
\rowcolor{HMTLight}eta & 46 & 158 & 0 & 547.836 & 547.862 & -47.457206\allowbreak{}3768 & evidencia numérica histórica condicionada; no constituye una regla física de selección anterior al contraste \\
ρ\({}^{0}\) & 50 & 176 & 0 & 775.267 & 775.26 & 9.029228\allowbreak{}90385 & evidencia numérica histórica condicionada; no constituye una regla física de selección anterior al contraste \\
\rowcolor{HMTLight}φ(1020) & 54 & 190 & 0 & 1019.475 & 1019.461 & 13.732747\allowbreak{}0104 & evidencia numérica histórica condicionada; no constituye una regla física de selección anterior al contraste \\
J/\allowbreak{}ψ & 64 & 237 & 0 & 3096.979 & 3096.916 & 20.342818\allowbreak{}4684 & evidencia numérica histórica condicionada; no constituye una regla física de selección anterior al contraste \\
\rowcolor{HMTLight}Υ(1S) & 89 & 334 & 0 & 9460.34 & 9460.3 & 4.228195\allowbreak{}72318 & evidencia numérica histórica condicionada; no constituye una regla física de selección anterior al contraste \\
\end{longtable}
\endgroup
\end{landscape}
Estas doce cifras no se fusionan con las veinte salidas de la tabla rectora. La malla conserva cinco etiquetas adicionales ---η, ρ\({}^{0}\), φ(1020), J/ψ y Υ(1S)---, pero aún no construye el mapa estado compuesto→ruta y registro genealógico→(n,k,t). El cuanto interno CT108, 0.058177\allowbreak{}641733\allowbreak{}144319\allowbreak{}230789\allowbreak{}692282\allowbreak{}953757 en la normalización hbar\_int=\allowbreak{}c\_int=\allowbreak{}t0=\allowbreak{}1, es una sección de la recta de masa y no una especie.

La tabla no delimita el dominio. El dominio interno continúa siendo el atlas de 81 celdas, 56 familias, 13 multisecciones y 324 rutas orientadas; las realizaciones físicas tipadas y el inventario de 471 masas y 384 anchuras se desarrollan después.

\Needspace{7\baselineskip}

\endgroup

\section{Correspondencia entre el atlas \(81/56/13/324\) y los inventarios exhaustivos de \(471\) masas y \(384\) anchuras}
\label{chap:portal-atlas-inventarios-masas}
El \hyperref[chap:inventario-universal-471-masas-384-anchuras]{inventario de
471 masas y 384 anchuras} es un catálogo metrológico externo con destino
tipado. No es un censo de predicciones ni el cardinal del dominio material.
